\documentclass[letter,11pt]{article}

%LAYOUT & SPACING
\usepackage[margin=1.2 in, top=1 in, bottom= 1.2 in]{geometry}

\usepackage[normalem]{ulem}

\usepackage{tikz}
\usetikzlibrary{hobby}
\usepackage{tikz-cd}
\usepackage{adjustbox}
\usepackage{fbox}

\usetikzlibrary{matrix, calc, arrows}


%\usepackage{natbib}
\usepackage{bibentry}
% \usepackage{fancyhdr}
\usepackage[linktocpage=true]{hyperref}

%FOOTNOTE

\usepackage{lipsum}
\newcommand\blfootnote[1]{%
  \begingroup
  \renewcommand\thefootnote{}\footnote{#1}%
  \addtocounter{footnote}{-1}%
  \endgroup
}

%EQUATION LAYOUT
\makeatletter
\g@addto@macro\normalsize{
  \setlength\abovedisplayskip{8pt}
  \setlength\belowdisplayskip{8pt}
  \setlength\abovedisplayshortskip{8pt}
  \setlength\belowdisplayshortskip{8pt}
  }
\makeatother

%TABLE OF CONTENT LAYOUT
\usepackage{tocloft}
\renewcommand\cftsecfont{\bfseries\footnotesize}
\renewcommand\cftsubsecfont{\footnotesize}
\renewcommand\cftsubsubsecfont{\footnotesize}
\renewcommand\cftsecpagefont{\bfseries\footnotesize}
\renewcommand\cftsubsecpagefont{\footnotesize}
\renewcommand\cftsubsubsecpagefont{\footnotesize}
\renewcommand\cftsecafterpnum{\vskip0pt}
\renewcommand\cftsubsecafterpnum{\vskip0pt}
\renewcommand\cftsubsubsecafterpnum{\vskip0pt}
\setlength{\cftbeforesecskip}{0pt}
\setlength\cftparskip{0pt}
\setlength\cftaftertoctitleskip{7pt}

%FOOTNOTE LAYOUT
\interfootnotelinepenalty=10000

%ALPHABETS
\usepackage[english]{babel}
\usepackage{amsfonts}
\usepackage{mathrsfs}
\usepackage{bbm}
\usepackage{latexsym}
\usepackage{math dots}
\usepackage{amssymb}
\usepackage{mathtools}

%ENVIRONMENTAL PACKAGES
\usepackage{enumitem}
\setlist{nolistsep}
\usepackage{amsthm}
%\usepackage{relsize}
%\usepackage{nicefrac}
\usepackage[capitalize]{cleveref}
\crefformat{footnote}{#2\footnotemark[#1]#3} %Adds footnote capability to cref

%Roman numerial enumeration in lists
\def\theenumi{\roman{enumi}}
\def\labelenumi{(\theenumi)}

% % LABEL CHECKING: checks for unused labels, modified to handle cleverref
% \usepackage[norefs]{refcheck}
% %%% Infrastructure    
% \makeatletter
% \newcommand{\refcheckize}[1]{%
%   \expandafter\let\csname @@\string#1\endcsname#1%
%   \expandafter\DeclareRobustCommand\csname relax\string#1\endcsname[1]{%
%     \csname @@\string#1\endcsname{##1}\wrtusdrf{##1}}%
%   \expandafter\let\expandafter#1\csname relax\string#1\endcsname
% }
% \makeatother
% %%%
% \refcheckize{\cref}
% \refcheckize{\Cref}
% % END LABEL CHECKING

% eqnitem is for itemize or enumerate and removes vertical space before an equation
\newcommand\eqnitem[1][]{%
  \ifx\relax#1\relax  \item \else \item[#1] \fi
  \abovedisplayskip=0pt\abovedisplayshortskip=0pt~\vspace*{-\baselineskip}}

%AMS THEOREMS
\newtheoremstyle{plain}{3mm}{3mm}{\slshape}{}{\bfseries}{.}{.5em}{}
\newtheoremstyle{definition}{2mm}{2mm}{}{}{\bfseries}{.}{.5em}{}
\theoremstyle{plain}
\newtheorem*{thm}{theorem}	
\newtheorem{theorem}{Theorem}
\newtheorem{fact}[theorem]{Fact}
\newtheorem*{fct}{Fact}
\newtheorem{lemma}[theorem]{Lemma}
\newtheorem{proposition}[theorem]{Proposition}
\newtheorem{corollary}[theorem]{Corollary}
\newtheorem{conjecture}[theorem]{Conjecture}
\newtheorem{problem}[theorem]{Problem}
\newtheorem{claim}[theorem]{Claim}
\newtheorem{question}[theorem]{Question}
\theoremstyle{definition}
\newtheorem{definition}[theorem]{Definition}
\newtheorem{remark}[theorem]{Remark}
\newtheorem{example}[theorem]{Example}
\newtheorem{examples}[theorem]{Examples}

%MAIN THEOREMS
\theoremstyle{plain}
\newcounter{MainTheoremCounter}
\renewcommand*{\theMainTheoremCounter}{\Alph{MainTheoremCounter}}
\newtheorem{Maintheorem}[MainTheoremCounter]{Theorem}
\newtheorem{Mainquestion}[MainTheoremCounter]{Question}

%NAMED THEOREMS
\theoremstyle{plain}
\newtheorem*{namedthm}{\namedthmname}
\newcounter{namedthm}
\makeatletter
	\newenvironment{named}[2]
	{\def\namedthmname{#1}
	\refstepcounter{namedthm}
	\namedthm[#2]\def\@currentlabel{#1}}
	{\endnamedthm}
\makeatother
%Example:
%\begin{named}{Incompleteness Theorem}{\cite{Godel1931}}
%\label{thm_incompleteness}
% statement of theorem
%\end{named}
%In 1931 Kurt G{\"o}del proved the \ref{theorem_incopleteness}.

%APPENDIX LAYOUT
\usepackage{chngcntr}
\counterwithin{theorem}{section}

%NUMBERWITHIN
\numberwithin{equation}{section}

%ALLOWDISPLAYBREAKS
\allowdisplaybreaks

%COLORS
\usepackage{xcolor}
\definecolor{Color2}{rgb}{0.78, 0.11, 0.0}%Scarlet
\hypersetup{citecolor = black,colorlinks,
			linkcolor = black,
			urlcolor = Color2}

%SECTION TITLE LAYOUT
\usepackage{titlesec}
\titleformat{\section}
  {\Large\center\bfseries}
  {\thesection.}{.7em}{}
\titlespacing*{\section}{0pt}{3.5ex plus 0ex minus 0ex}{1.5ex plus 0ex}
\titleformat{\subsection}
  {\large\center\bfseries}
  {\thesubsection.}{.7em}{}
\titlespacing*{\subsection}{0pt}{3.5ex plus 0ex minus 0ex}{1.5ex plus 0ex}
\titleformat{\subsubsection}
  {\center\bfseries}
  {\thesubsubsection.}{.7em}{}
\titlespacing*{\subsubsection}{0pt}{3.5ex plus 0ex minus 0ex}{1.5ex plus 0ex}

%CONTENTS TITLE LAYOUT
\addto\captionsenglish{\renewcommand{\contentsname}{\normalsize\bfseries Contents}}

%DOCUMENT TITLE LAYOUT
\usepackage{titling}
\setlength{\droptitle}{-4em}

%BIBLIOGRAPHY LAYOUT
%\usepackage{doi}
%\usepackage[backend=biber,style=alphabetic]{biblatex}
%\addbibresource{mynewlibrary.bib}

%Affiliations package
\usepackage{authblk}

%PEEPS
\newcommand{\Caratheodory}{Carath\'{e}odory}
\newcommand{\Cech}{\v{C}ech}
\newcommand{\Cesaro}{Ces\`{a}ro}
\newcommand{\Erdos}{Erd\H{o}s}
\newcommand{\Folner}{F\o{}lner}
\newcommand{\Poincare}{Poincar\'{e}}
\newcommand{\Szemeredi}{Szemer\'{e}di}
\newcommand{\Beiglbock}{Beiglb\"{o}ck}
\newcommand{\Fejer}{Fej\'{e}r}
\newcommand{\Sarkozy}{S\'{a}rk\"{o}zy}
\newcommand{\Katai}{K\'{a}tai}
\newcommand{\Mobius}{M\"{o}bius}
\newcommand{\KulagaPrzymus}{Ku{\l}aga-Przymus}
\newcommand{\Lemanczyk}{Lema\'{n}czyk}
\newcommand{\Halasz}{Hal\'{a}sz}
\newcommand{\Gauss}{Gau{\ss}}
\newcommand{\Turan}{Tur{\'a}n}
\hyphenation{Min-kow-ski}

%STANDARD COMMANDS
\newcommand{\concat}{\hspace{-.02em}{\mathbin{\raisebox{.75ex}{\scalebox{.9}{${\smallfrown}$}}}}\hspace{-.02em}}
\newcommand{\supp}{{\normalfont\text{supp}}\,}
\newcommand{\Oh}{{\rm O}}
\newcommand{\oh}{{\rm o}}
\newcommand{\eps}{\epsilon}
\newcommand{\N}{\mathbb{N}}
\newcommand{\Z}{\mathbb{Z}}
\newcommand{\R}{\mathbb{R}}
\newcommand{\C}{\mathbb{C}}
\newcommand{\Q}{\mathbb{Q}}
\newcommand{\T}{\mathbb{T}}
\newcommand{\Nz}{\N_{\geq 0}}
\newcommand{\defeq}{\vcentcolon=}
\newcommand{\eqdef}{=\vcentcolon}
\newcommand\restr[2]{{ \left.\kern-\nulldelimiterspace #1 \right|_{#2}}}

%LAYOUT COMMANDS
\newcommand{\define}[1]{{\itshape #1}}

%RENEWED COMMANDS
\renewcommand{\epsilon}{\varepsilon}
\renewcommand{\leq}{\leqslant}
\renewcommand{\geq}{\geqslant}
\renewcommand{\setminus}{\backslash}
\renewcommand{\Re}{{\rm Re}}
\renewcommand{\Im}{{\rm Im}}
\renewcommand{\P}{\mathbb{P}}

%MEASURE THEORY COMMANDS
\newcommand{\E}{\mathbb{E}}

%EDITING COMMANDS
\newcommand{\edit}[3]{\color{#1}{#3}\color{black}\marginpar{\textcolor{#1}{[[#2]]}}}
\newcommand{\dgg}[1]{\edit{blue}{DGG}{#1}}
\newcommand{\anh}[1]{\edit{red}{ANH}{#1}}
\newcommand{\dggtodo}[1]{\edit{blue}{DGG TODO}{#1}}
\newcommand{\anhtodo}[1]{\edit{red}{ANH TODO}{#1}}

\newcommand{\old}[1]{\edit{orange}{OLD}{#1}}
\usepackage[normalem]{ulem}
\newcommand{\soutcomment}[1]{\textcolor{red}{\sout{#1}}}


%DOCUMENT SPECIFIC COMMANDS

\newcommand{\orb}{\overline{o}}

\DeclareMathOperator{\upclose}{\uparrow}
\DeclareMathOperator{\bigupclose}{\big\uparrow}

\newcommand{\syndetic}{\mathcal{S}}
\newcommand{\thick}{\mathcal{T}}

\newcommand{\dsyndetic}{d\mathcal{S}}
\newcommand{\dthick}{d\mathcal{T}}

\newcommand{\dcsyndetic}{dc\mathcal{S}}
\newcommand{\dcthick}{dc\mathcal{T}}


\newcommand{\dsyndeticdistal}{d\mathcal{S}_{\text{distal}}}
\newcommand{\dcsyndeticdistal}{dc\mathcal{S}_{\text{distal}}}
\newcommand{\dthickdistal}{d\mathcal{T}_{\text{distal}}}

\newcommand{\dcthickdistal}{dc\mathcal{T}_{\text{distal}}}

\newcommand{\PS}{\mathcal{PS}}
\newcommand{\dPS}{d\mathcal{PS}}
\newcommand{\dcPS}{dc\mathcal{PS}}

\newcommand{\starters}{\mathcal{I}}
\newcommand{\cubes}{\mathcal{Q}}
\newcommand{\sd}{\text{\sc{sd}}}

\newcommand{\central}{\mathcal{C}}
\newcommand{\IP}{\mathcal{IP}}
\newcommand{\IPzero}{\mathcal{IP}_0}

\newcommand{\IR}{\mathcal{IR}}



\newcommand{\paper}[1]{\cite{#1} \bibentry{#1}}

\newcommand{\filter}{\mathscr{F}}
\newcommand{\filtertwo}{\mathscr{G}}
\newcommand{\filterthree}{\mathscr{H}}

\newcommand{\class}{\mathscr{F}}
\newcommand{\classone}{\mathscr{F}}
\newcommand{\classtwo}{\mathscr{G}}
\newcommand{\classcap}{\sqcap}
\newcommand{\classcup}{\sqcup}

\newcommand{\family}{\mathscr{F}}
\newcommand{\familyone}{\mathscr{F}}
\newcommand{\familytwo}{\mathscr{G}}
\newcommand{\familythree}{\mathscr{H}}
\newcommand{\familycap}{\sqcap}

\newcommand{\collection}{\mathscr{C}}


\newcommand{\odometer}{Z}
\newcommand{\orotation}{\rho}


\newcommand{\polyret}{R}

% \newcommand\restr[2]{{ \left.\kern-\nulldelimiterspace #1 \right|_{#2}}}

\newcommand{\mpc}{\text{mpc}}
\newcommand{\FS}{\text{FS}}


\newcommand{\bpunch}{\beta}
\newcommand{\apunch}{\alpha}

\begin{document}

\title{Dynamically syndetic sets and the combinatorics of syndetic, idempotent filters}

\author[1]{Daniel Glasscock}
\author[2]{Anh N. Le}

\affil[1]{\small Dept. of Mathematics and Statistics, U. of Massachusetts Lowell, Lowell, MA, USA}
\affil[2]{\small Dept. of Mathematics,
	U. of Denver, Denver, CO, USA}

\date{}

\maketitle

%End comments

\blfootnote{2020 \emph{Mathematics Subject Classification.}  Primary: 37B20. Secondary: 37B05.}


\blfootnote{\emph{Key words and phrases. }
return time sets in minimal topological dynamical systems,
syndetic sets,
thick sets,
piecewise syndetic sets,
central sets,
IP sets,
%thickly syndetic sets,
% central syndetic sets,
% strongly central sets,
dynamically syndetic sets,
% dynamically central syndetic sets,
dynamically thick sets,
% dynamically piecewise syndetic sets,
sets of pointwise recurrence,
%translation-invariant filters,
%minimal topological dynamical systems,
topological recurrence,
measurable recurrence,
%sets of polynomial recurrence,
idempotent filters,
%ultrafilters,
idempotent ultrafilters}


\begin{abstract}
A subset of the positive integers is \emph{dynamically central syndetic} if it contains the set of times that a point returns to a neighborhood of itself under a minimal transformation of a compact metric space.  These sets are part of the highly-influential link between dynamics and combinatorics forged by Furstenberg and Weiss in the 1970's.  Our main result is a characterization of dynamically central syndetic sets as precisely those sets that belong to syndetic, idempotent filters.
Idempotent filters are combinatorial objects that abound in ergodic Ramsey theory but have been largely unnoticed and unexplored.  We develop the algebra of these objects for the proof of the main theorem and with an eye toward future applications.

The main result is best contextualized as a ``global'' analogue to the ``local'' characterization of Furstenberg's central sets as members of minimal, idempotent ultrafilters.  It leads to a dual characterization of sets of topological pointwise recurrence, allowing us to answer a question of Glasner, Tsankov, Weiss, and Zucker.  We draw numerous striking contrasts between pointwise recurrence and set recurrence, a topic with a long history in the subject and its applications, and answer four questions posed by Host, Kra, and Maass. We also show that the intersection of a dynamically central syndetic set with a set of pointwise recurrence must be piecewise syndetic, generalizing results of Dong, Glasner, Huang, Shao, Weiss, and Ye.
\end{abstract}

\newpage

\setcounter{tocdepth}{2}
\tableofcontents


\section{Introduction}

\subsection{Part I: Dynamically syndetic sets}
\label{sec_part_one_intro}

A subset of the positive integers, $\N$, is called \emph{dynamically syndetic} if it contains a set of the form
\[R_T(x,U) \defeq \big\{ n \in \N \ \big| \ T^n x \in U \big\}, \quad x \in X, \quad U \subseteq X \text{ nonempty, open},\]
where $(X,d_X)$ is a compact metric space and the map $T: X \to X$ generates a minimal action of the semigroup $(\N,+)$ on $X$ by continuous maps (see \cref{sec_top_dynamics}). In this case, the pair $(X,T)$ is called a \emph{minimal topological dynamical system}.  If $x \in U$, such a set is called \emph{dynamically central syndetic}.  Dynamically syndetic and dynamically central syndetic sets are the main objects of study in this paper.


Examples of dynamically syndetic sets abound in combinatorics and number theory.  Periodic sets and their almost-periodic and higher-order generalizations, namely Bohr and nil-Bohr sets, are dynamically syndetic.  The images of Beatty sequences and the set of those positive integers with even $2$-adic valuation are examples of dynamically central syndetic sets. We give further examples in \cref{sec_dsyndetic} and an instructive non-example at the beginning of \cref{sec_intro_char_of_ds_sets} below.



\subsubsection{Motivation: structure in dynamics and combinatorics}

Dynamically syndetic sets are structural objects that form part of the bridge between dynamics and combinatorics.  We will mention two highly influential sets of results that exemplify this.

Furstenberg and Weiss \cite{Furstenberg-ErgodicBehavior, Furstenberg_Weiss-topological-dynamics} pioneered the use of ergodic theory and topological dynamics as tools in Ramsey theory and combinatorial number theory.  A \emph{correspondence principle} -- a device used to create a dynamical object out of a combinatorial one -- makes available a wide range of tools and ideas from dynamics.
For example, a topological correspondence principle can be used to reduce van der Waerden's theorem -- ostensibly concerning arbitrary finite partitions of $\N$ in its original formulation \cite{van_der_waerden_1927} or arbitrary syndetic subsets of $\N$ in an equivalent formulation \cite{kakeya_morimoto_1930} -- to a statement concerning dynamically syndetic sets.
The first step in Furstenberg's highly influential proof of \Szemeredi{}'s theorem \cite{Furstenberg-ErgodicBehavior} is also of this nature, transferring the problem from combinatorics to one about measure preserving systems.  This set of ideas has had a huge impact on the field and continues to yield dividends \cite{ward_ergodic_theory_interactions_survey}.


Rotations on compact abelian groups are among the simplest examples of dynamical systems, and the dynamically syndetic sets they generate, \emph{Bohr sets}, are  among the simplest examples of dynamically syndetic sets.  Nilsystems -- algebraic generalizations of compact group rotations -- give rise to dynamically syndetic sets called \emph{nil-Bohr sets}.
Since groundbreaking work of Host and Kra \cite{host_kra_2005} and Ziegler \cite{ziegler_2007}, nilsystems and nil-Bohr sets have been used to describe the ``structured'' components in powerful decomposition theorems in ergodic theory \cite{Bergelson_Host_Kra05, frant15, leibman_2010}, additive combinatorics and number theory \cite{Green_Tao10, Green_Tao12, Green_Tao12a, Green_Tao_Ziegler12} and, recently, in topological dynamics \cite{glasner_huang_shao_weiss_ye_2020}.
Thus, dynamically syndetic sets arising from highly structured systems have become essential in defining structure, even in non-dynamical settings.


In their full generality, the families of dynamically (central) syndetic sets have appeared explicitly in the literature in a few places.  In various combinations, Dong, Huang, Shao, and Ye \cite{dong_shao_ye_2012,Huang_Shao_Ye_2019,huang_ye_2005} have studied these classes (they call them $m$-sets  and $sm$-sets) in the context of disjointness between topological dynamical systems.  The family of dynamically central syndetic sets was studied in a combinatorial context by Bergelson, Hindman, and Strauss \cite{bergelson_hindman_strauss_2012} in connection to the quality of sets of return times of points to open sets under polynomial iterates of an irrational rotation.  A subfamily of dynamically syndetic sets was considered by Kennedy, Raum, and Salomon \cite{kennedy_raum_salomon_2022}; we will discuss some of their results more carefully in the context of our own below.



\subsubsection{The main result: a characterization of dynamically syndetic sets}
\label{sec_intro_char_of_ds_sets}

A subset of $\N$ is \emph{syndetic} if it has bounded gaps, that is, if there exists $N \in \N$ such that it has nonempty intersection with every interval of length $N$.  Syndeticity is a generalization of periodicity that is intimately connected to minimality in topological dynamics.  That dynamically syndetic sets are syndetic, for example, is a consequence of a well-known characterization of minimal actions: the system $(X,T)$ is minimal if and only if every set of the form $R_T(x,U)$ is syndetic.

It is not the case, however, that all syndetic sets are dynamically syndetic.  Indeed, here is an instructive example to keep in mind:
\begin{align}
    \label{eqn_intro_even_odd}
    \bigg( 2 \N \cap \bigcup_{\substack{n = 0 \\ \text{$n$ even}}}^\infty \big[ 2^{n}, 2^{n + 1}\big) \bigg) \cup \bigg( \big( 2\N + 1 \big) \cap \bigcup_{\substack{n = 0 \\ \text{$n$ odd}}}^\infty \big[ 2^{n}, 2^{n + 1}\big) \bigg).
\end{align}
That this set is syndetic is obvious.  That this set is not dynamically syndetic is less obvious.  We leave the discovery of a short argument to the intrepid reader, opting instead to show this as an application of \cref{mainthm_ds_equivalents} below.

This example raises a question that has appeared explicitly and implicitly in a number of different contexts in the literature: how do we know if a given set is dynamically syndetic? The following theorem -- the main result of our paper -- answers this question by giving necessary and sufficient conditions for a set to be dynamically syndetic.  \cref{mainthm_ds_equivalents} is proved in \cref{sec_char_of_ds_and_dcs}.


\begin{Maintheorem}
\label{mainthm_ds_equivalents}
    Let $A \subseteq \N$.  The following are equivalent.
    \begin{enumerate}
        \item
        \label{item_intro_ds_def_condition}
        The set $A$ is dynamically syndetic.

        \item
        \label{item_intro_ds_combo_condition}
        There exists a nonempty subset $B \subseteq A$ that satisfies: for all finite $F \subseteq B$, the set $\bigcap_{f \in F} (B-f)$ is syndetic.

        \item
        \label{item_intro_translate_belongs_to_sif}
        There exists $n \in \N$ for which the set $A - n$ belongs to a syndetic, idempotent filter on $\N$.
    \end{enumerate}
\end{Maintheorem}

Point \eqref{item_intro_ds_combo_condition} characterizes dynamical syndeticity in terms of a simple combinatorial property.  Let us use it to see that the set in \eqref{eqn_intro_even_odd}, call it $A$, is not dynamically syndetic.  If $B \subseteq A$ satisfies the condition in \eqref{item_intro_ds_combo_condition}, then it must be syndetic.  Every syndetic subset of $A$, in particular, $B$, must contain both an even integer, $b_0$, and an odd integer, $b_1$.  It is quick to see, however, that the set $(A - b_0) \cap (A-b_1)$ is not syndetic, contradicting what we had supposed about $B$.  In addition to giving a way to check whether a set is dynamically syndetic, the combinatorial condition in \eqref{item_intro_ds_combo_condition} allows us to characterize members of the dual class, the class of dynamically thick sets, as we define and discuss in \cref{sec_intro_part_two}.  

It is worth comparing point \eqref{item_intro_ds_combo_condition} in \cref{mainthm_ds_equivalents} with a recent result of Kennedy, Raum, and Salomon \cite[Thm. 8.2]{kennedy_raum_salomon_2022}. In the context of the positive integers, it is a consequence of their result that a set $A \subseteq \N$ is dynamically syndetic if and only if there exists $B \subseteq A$ such that for all finite $F_1 \subseteq B$ and $F_2 \subseteq \N \setminus B$, the set 
\begin{align}\
\label{eq:new_B'}
    \bigcap_{f_1 \in F_1} \big(B - f_1\big)  \cap  \bigcap_{f_2 \in F_2} \big((\N \setminus B) - f_2\big)
\end{align}
is syndetic.  It is not hard to see that a set $B \subseteq \N$ satisfies the condition in \eqref{eq:new_B'} if and only if its indicator function $1_B$ is a uniformly recurrent point in the full shift $(\{0,1\}^{\N}, \sigma)$; see \cref{sec_top_dynamics} for the notation and terminology.  That a set is dynamically syndetic if and only if it contains such a subset follows then from, eg., \cite[Prop. 2.3]{huang_ye_2005} (see also \cref{thm:simple_equivalent_dS}).  The condition in \cref{mainthm_ds_equivalents} \eqref{item_intro_ds_combo_condition} is substantially weaker than that in \eqref{eq:new_B'}.  In particular, it does not imply that the point $1_B$ is uniformly recurrent.  This weaker condition translates to an easier-to-verify sufficient condition for dynamical syndeticity.

The following is an analogous characterization for dynamically central syndetic sets. \cref{intro_mainthm_full_chars_of_dcsyndetic} is proved in \cref{sec_char_of_ds_and_dcs}.

\begin{Maintheorem}
    \label{intro_mainthm_full_chars_of_dcsyndetic}
    Let $A \subseteq \N$.  The following are equivalent.
    \begin{enumerate}
        \item \label{intro_item_def_of_dcs}
        The set $A$ is dynamically central syndetic.

        \item \label{intro_item_dcs_combo_condition}
        There exists a subset $B \subseteq A$ that satisfies: for all finite $F \subseteq B$, the set $B \cap \bigcap_{f \in F} (B-f)$ is syndetic.

        \item \label{intro_item_member_of_sif}
        The set $A$ belongs to a syndetic, idempotent filter on $\N$.
    \end{enumerate}
\end{Maintheorem}

The novelty and difficulty of Theorems \ref{mainthm_ds_equivalents} and \ref{intro_mainthm_full_chars_of_dcsyndetic} lie in the third conditions, ones that characterize dynamical syndeticity in terms of membership in a combinatorial object called a syndetic, idempotent filter.  This connection is, in fact, behind most of the other results in this paper.  We turn our attention toward these objects now.

\subsubsection{A new tool: the algebra of idempotent families}
\label{sec_intro_algebra_of_filters}

A \emph{family} is an upward-closed collection of sets.  A \emph{filter} is a family $\family$ satisfying: for all members $A_1$ and $A_2$ of $\family$, the set $A_1 \cap A_2$ belongs to $\family$.  Filters are combinatorial objects membership in which formalizes what it means for a set to be ``large'' in some way. The collection of all neighborhoods of a given point in a topological space is the quintessential example of a filter in topology.  We will see that an analogue of this filter for minimal topological dynamical systems is inextricably linked to the notion of dynamical syndeticity.



Let us briefly narrow the discussion from filters to ultrafilters.  \emph{Ultrafilters on $\N$} -- maximal filters on $\N$, whose existence is guaranteed by Zorn's lemma -- are combinatorial objects with a deep and important history in the study of topological dynamics, especially in applications to combinatorial number theory \cite{bergelson_2010}. The set $\beta \N$ of ultrafilters -- identifiable with the Stone-\Cech{} compactification of $\N$ -- is a multifaceted topological and algebraic object, the structure of which is behind many of the aforementioned applications.  For the purposes of this introduction, it suffices to know that $\beta \N$ is a compact Hausdorff space, and that there is a natural way to lift addition from $\N$ to $\beta \N$, turning $\beta \N$ into a semigroup: if $\family$ and $\familytwo$ are ultrafilters, then
\begin{align}
\label{eqn_sum_of_families_intro}
    \family + \familytwo \defeq \Big\{ A \subseteq \N \ \Big| \ \big\{ n \in \N \ \big| \ A - n \in \familytwo \big\} \in \family \Big\}
\end{align}
is an ultrafilter.  The repeated addition of the principal ultrafilter $\{A \subseteq \N \ | \ 1 \in A\}$ on the left turns $\beta \N$ into a (non-metrizable) topological dynamical system.

An ultrafilter $\family$ is called \emph{idempotent} if $\family = \family + \family$.  Idempotency of ultrafilters is closely related to proximality in topological dynamics \cite[Sec. 19.3]{hindman_strauss_book_2012} and to finite sums sets in additive combinatorics \cite[Sec. 5.2]{hindman_strauss_book_2012}.  An ultrafilter is called \emph{minimal} if it is \emph{uniformly recurrent} (see \cref{sec_top_dynamics}) under the dynamics just described.  Minimality of ultrafilters is closely related to uniform recurrence \cite[Sec. 19.3]{hindman_strauss_book_2012} and piecewise syndeticity \cite[Sec. 4.4]{hindman_strauss_book_2012}.  At the intersection of all of these ideas are Furstenberg's \emph{central sets} \cite[Def. 8.3]{furstenberg_book_1981}, combinatorially-rich, dynamically-defined subsets of $\N$ introduced to aid in the topological dynamical approach to problems in Ramsey theory. It is a result of Bergelson, Hindman, and Weiss \cite{bergelson_hindman_1990} that a subset of $\N$ is central if and only if it belongs to a minimal, idempotent ultrafilter on $\N$.

There is no obstruction to extending the familiar definition of sums of ultrafilters in \eqref{eqn_sum_of_families_intro} to arbitrary families: if $\family$ and $\familytwo$ are families, the sum $\family + \familytwo$ defined in \eqref{eqn_sum_of_families_intro} is again a family.  But while there are books written about the algebra of ultrafilters (we cite, eg., \cite{hindman_strauss_book_2012} frequently in this work), the algebra of families appears to be much less explored. For filters, the sum in \eqref{eqn_sum_of_families_intro} appears, to our knowledge, first in Berglund and Hindman \cite{berglund_hindman_1984}. We define a family $\family$ to be \emph{idempotent} if $\family \subseteq \family + \family$.  (See \cref{sec_trans_inv_and_idempotent_families} for a discussion on containment versus equality.) This definition generalizes the notion of ``idempotent filter'' that appears implicitly in Papazyan \cite{papazyan_1989} and explicitly in Krautzberger \cite{Krautzberger_2009}.

In the setting of topological dynamics, the canonical example of an idempotent filter is one of the form
\begin{align}
\label{eqn_canonical_sif}
    \family \defeq \upclose \big\{ R_T(x,U) \ \big| \ U \subseteq X \text{ is open and contains $x$} \big\},
\end{align}
where $x$ is a point in a dynamical system $(X,T)$ and the upward arrow indicates the upward closure of the collection of sets.  That $\family$ is a filter is clear.  To see that this is an idempotent filter, one needs only to show and interpret the following: for all open $U \subseteq X$ containing $x$,
\[R_T(x,U) \subseteq \big\{n \in \N \ \big| \ R_T(x,U) - n \in \family\big\}.\]
If $n \in R_T(x,U)$, then $x \in T^{-n}U$, whereby $R_T(x,U) - n = R_T(x,T^{-n}U)$ is a member of $\family$.  Since $R_T(x,U) \in \family$, we see from the inclusion that $\{n \in \N \ | \ R_T(x,U) - n \in \family\} \in \family$, demonstrating that $R_T(x,U) \in \family + \family$.  Since every member of $\family$ contains a set of the form $R_T(x,U)$, this shows that $\family \subseteq \family + \family$.

A \emph{syndetic filter} is a filter whose members are all syndetic.  If $(X,T)$ is a minimal system, the filter in \eqref{eqn_canonical_sif} is an example of a syndetic, idempotent filter.  Thus, we see that every dynamically central syndetic set belongs to a syndetic, idempotent filter. \cref{intro_mainthm_full_chars_of_dcsyndetic} \eqref{intro_item_member_of_sif} gives the converse: every member of a syndetic, idempotent filter is dynamically central syndetic.  If ultrafilters are ``local'' combinatorial objects, then syndetic filters are ``global'' combinatorial objects.  One way to contextualize \cref{intro_mainthm_full_chars_of_dcsyndetic}, then, is as a global analogue to the local characterization of central sets mentioned above.\\


There are several other ``large,'' idempotent filters that arise naturally in the subject.  All three of the following examples play an important role in the results in this paper. See \cref{sec_prelims} for definitions of the families $\central^*, \IP^*$, and $\PS^*$ (thickly syndetic sets).
\begin{enumerate}
    \item By a theorem of Furstenberg \cite[Thm. 9.11]{furstenberg_book_1981}, if $x$ is a distal point, the filter defined in \eqref{eqn_canonical_sif} is a $\central^*$, idempotent filter.  Such filters feature prominently in \cref{sec_filter_algebra_for_central_duals}.

    \item Very general families of sets of return times in topological dynamics and ergodic theory are $\IP^*$, idempotent filters. Our demonstration of this in a concrete example in \cref{lemma_localization_is_syndetic_idempotent_filter} is behind the proof of \cref{thm:dcps_multiple_recurrence} below. 

    \item The family of thickly syndetic sets is a $\PS^*$, translation-invariant filter (see \cref{lemma_condition_on_subfamily_of_ps_star}).  This filter is behind some of the main applications of our results, such as in \cref{cor:largest_family_satisfies_cap}.
\end{enumerate}
We believe that these examples and the results in this paper demonstrate the usefulness of the algebra of families of large sets and motivate its further development.



\subsubsection{A peek into the proofs of Theorems \ref{mainthm_ds_equivalents} and \ref{intro_mainthm_full_chars_of_dcsyndetic}}
\label{sec_intro_into_the_proof}

Let us briefly discuss the proofs of Theorems \ref{mainthm_ds_equivalents} and \ref{intro_mainthm_full_chars_of_dcsyndetic}.  The full proofs, finished in \cref{sec_char_of_ds_and_dcs}, occupy a substantial portion of Part I of this paper. Since any dynamically syndetic set can be translated to become dynamically central syndetic (\cref{lemma_translates_of_dsyndetic_sets}), \cref{mainthm_ds_equivalents} follows relatively easily from \cref{intro_mainthm_full_chars_of_dcsyndetic}.

That the combinatorial condition in \cref{intro_mainthm_full_chars_of_dcsyndetic} \eqref{intro_item_dcs_combo_condition} characterizes membership in a syndetic, idempotent filter, \eqref{intro_item_member_of_sif}, is proved in \cref{sec_combo_condition_for_central_syndetic}.  The proof is purely combinatorial and highlights some of the family algebra discussed in the previous section.  Readers familiar with the characterization of $\IP$ sets (see \cref{sec_ip_and_central_sets}) as members of idempotent ultrafilters will immediately recognize the ideas at play.

That dynamically central syndetic sets belong to a syndetic, idempotent filter (that \eqref{intro_item_def_of_dcs} implies \eqref{intro_item_member_of_sif}) is simple and was discussed in the previous section. The challenge lies in the converse, proved in \cref{sec_char_of_dy_central_synd_sets}. Since this is the most technical part of the paper, let us briefly describe the main idea.  A more detailed outline is given in \cref{sec_main_idea_of_shift_punch}. 


We begin with a set $A \subseteq \N$ that satisfies the combinatorial condition in \eqref{intro_item_dcs_combo_condition}.  Were the point $1_{A \cup \{0\}}$ uniformly recurrent in the full symbolic shift $(\{0,1\}^{\N \cup \{0\}},\sigma)$, it would follow from \cref{thm:simple_equivalent_dcS} that $A$ is dynamically central syndetic.  It is not true in general, however, that $1_{A \cup \{0\}}$ is uniformly recurrent.  The main idea is to modify the symbolic shift by composing it with a ``punch'' map to form a related system that we call ``shift-punch.''  Under the shift-punch dynamics, the point $1_{A \cup \{0\}}$ is uniformly recurrent.  It is this uniform recurrence, then, that allows us to find a subset of $A$ that is dynamically central syndetic.


\subsection{Part II: Dynamically thick sets}
\label{sec_intro_part_two}

Dual to the families of dynamically (central) syndetic sets are the families of dynamically thick sets and sets of pointwise recurrence.  A set $A \subseteq \N$ is \emph{dynamically thick} if for all minimal systems $(X,T)$ and all points $x \in X$, the set $\{T^n x \ | \ n \in A\}$ is dense in $X$.  The set $A$ is a \emph{set of (minimal topological) pointwise recurrence} if for all minimal systems $(X,T)$, all points $x \in X$, and all open $U \subseteq X$ containing $x$, there exists $a \in A$ such that $T^a x \in U$.

\emph{Thick} sets -- subsets of $\N$ that contain arbitrary long intervals -- are the most immediate examples of dynamically thick sets.  It is quick to check that a set is dynamically thick if and only if its complement is not dynamically syndetic.  Thus, the complement of the set in \eqref{eqn_intro_even_odd} (and, in fact, the set itself) is an example of a dynamically thick set that is, coincidentally, not thick.  Further examples of dynamically thick sets and sets of pointwise recurrence are given in \cref{sec_dy_thick_def_and_char}.


\subsubsection{Motivation: set and point recurrence}
\label{sec_intro_set_and_pt_recurrence}

The phenomenon of recurrence in dynamical systems has been studied extensively in numerous forms.  We will discuss two of those forms most relevant to this work: set recurrence and point recurrence in ergodic theory and topological dynamics.


In a probability measure space $(X, \mu)$ with a measure-preserving transformation $T: X \to X$, a set $E \subseteq X$ of positive measure is said to recur at time $n \in \N$ if $\mu(E \cap T^{-n} E) > 0$.  A set $A \subseteq \N$ is called a \emph{set of measurable recurrence} if in every probability measure preserving system, every set of positive measure recurs at some time in $A$.  The question of which times a set of positive measure can be guaranteed to recur -- that is, the question of which sets are sets of measurable recurrence -- is important both conceptually and for applications.  
Szemer\'edi's theorem \cite{Szemeredi-on_the_sets}, a cornerstone of additive combinatorics, states that every set of positive integers of positive upper density contains arbitrarily long arithmetic progressions. Furstenberg \cite{Furstenberg-ErgodicBehavior} famously reproved this theorem by connecting the existence of arithmetic progressions in sets of positive upper density to the existence of sets of multiple measurable recurrence.
Since then, this connection has been explored and exploited to great effect \cite{ward_ergodic_theory_interactions_survey}.

Set recurrence in topological dynamics is also well-studied.  A set $A \subseteq \N$ is a \emph{set of topological recurrence} if for all minimal systems $(X,T)$ and all nonempty, open sets $U \subseteq X$, there exists $n \in A$ such that the set $U \cap T^{-n}U$ is nonempty.  The connection between topological dynamics and partition combinatorics forged by Furstenberg and Weiss \cite{Furstenberg_Weiss-topological-dynamics} parallels the connection between measurable dynamics and density combinatorics described above and has similarly been influential in additive combinatorics and Ramsey theory.\\


While the dynamical behavior of points in measurable dynamics is necessarily limited to generic points, meaningful statements can be made about the dynamical behavior of all points in topological dynamics.  This is manifest in the ``for all $x \in X$'' quantifier in the definitions of dynamically thick sets and sets of pointwise recurrence.  Despite the fact that these definitions are naturally analogous to those in the topic of set recurrence, far less has been written on them.  Many of our results in this paper are best understood as showing that the nature of pointwise recurrence differs substantially from that of set recurrence.


It is a consequence of a result of Pavlov \cite{Pavlov-2008} that a set of pointwise recurrence must have positive upper Banach density.  Indeed, he showed that if a set $A \subseteq \N$ has zero upper Banach density ($\limsup_{N \to \infty} \max_{n \in \N} |A \cap \{n, \ldots, n + N\}| / N = 0$) then the set $A$ is not a set of pointwise recurrence: there is a totally minimal, totally uniquely ergodic, and topologically mixing system $(X, T)$ and uncountably many points $x \in X$ for which $x \not\in \overline{\{T^n x \ \big| \ n \in A\}}$.  That dynamically thick sets and sets of pointwise recurrence are \emph{piecewise syndetic} (the intersection of a syndetic set and a thick set) are the results of Huang and Ye \cite[Thm. 2.4]{huang_ye_2005} and Dong, Song, Ye \cite[Prop. 4.4]{dong_shao_ye_2012}, respectively.  We improve on both of these results in \cref{thm:relations_C_dPS_PS-intro}, discussed below.


Pointwise recurrence is implicitly studied in \cite{Glasner-Weiss-Interpolation,Koutsogiannis_Le_Moreira_Pavlov_Richter_interpolation_sets} in the context of interpolation sets for minimal systems. A set $I \subseteq \N_0$ is \emph{interpolation set for minimal systems} if for every bounded function $f: I \to \C$, there exists a minimal system $(X, T)$, a point $x \in X$, and a continuous function $F: X \to \C$ such that $f(n) = F(T^n x)$ for all $n \in I$. It was proved in \cite{Glasner-Weiss-Interpolation, Koutsogiannis_Le_Moreira_Pavlov_Richter_interpolation_sets} that $I \subseteq \N_0$ is an interpolation set for minimal systems if and only if $I$ is not piecewise syndetic. If $A \subseteq \N$ is not piecewise syndetic, then $I \defeq A \cup \{0\}$ is not piecewise syndetic, so there is minimal system $(X, T)$, a point $x \in X$, and a continuous function $F: X \to \C$ such that 
\[
    F(T^n x) = \begin{cases}
            1 \text{ if } n = 0 \\
            0 \text{ if } n \in A 
    \end{cases}.
\]
It follows that $A$ is not a set of pointwise recurrence. Contrapositively, a set of pointwise recurrence must be piecewise syndetic, recovering the results in \cite[Prop. 4.4]{dong_shao_ye_2012} and, hence, in \cite[Thm. 2.4]{huang_ye_2005} as well.
    
As far as we know, the term ``set of pointwise recurrence'' was first introduced by Host, Kra, and Maass \cite{HostKraMaass2016}, who compared and contrasted several different types of recurrence: single and multiple topological recurrence, topological recurrence for nilsystems, pointwise recurrence, etc. They observed that the family of sets of pointwise recurrence is properly contained in the family of sets of topological recurrence.  We answer in this paper several of the questions they raise.

Finally, the families of dynamically thick sets and sets of pointwise recurrence have appeared explicitly in the literature in the context of disjointness.  We discuss these appearances in \cref{sec_intro_disjointness} below.

\subsubsection{Combinatorial characterizations of dynamically thick sets}

There is a useful and well-known combinatorial characterization for sets of topological recurrence (cf. \cite[Thm. 2.4]{bergelson_mccutcheon_1998}): a set $A \subseteq \N$ is a set of topological recurrence if and only if for all syndetic sets $S' \subseteq S \subseteq \N$, the set $A \cap (S - S')$ is nonempty.  Host, Kra, and Maass \cite[Question 2.13]{HostKraMaass2016} asked whether there is a purely combinatorial characterization of sets of pointwise recurrence.  Our next result, proved in \cref{sec_dthick_comb_characterizations}, gives an answer to this question.


\begin{Maintheorem}
\label{mainthm_dct_combo_characterizations}
    Let $A \subseteq \N$.  The following are equivalent.
    \begin{enumerate}
        \item
        \label{item_intro_dct_def_char}
        $A$ is a set of pointwise recurrence.

        \item
        \label{item_intro_dct_sif_char}
        For all $B \supseteq A$, there exists a finite set $F \subseteq \N \setminus B$ such that the set $B \cup \big( B-F \big)$ is thick.

        \item
        \label{item_intro_condition_one_dct}
        For all syndetic $S \subseteq \N$, there exists a finite set $F \subseteq A$ such that for all syndetic $S' \subseteq S$,
        \[F \cap (S - S') \neq \emptyset.\]
    \end{enumerate}
\end{Maintheorem}

The statement in \eqref{item_dct_sif_char} is derived from the combinatorial characterization of dynamically central syndetic sets in \cref{intro_mainthm_full_chars_of_dcsyndetic} and can be used to produce examples of sets of pointwise recurrence.  We demonstrate this in the analogous context of dynamically thick sets below.  The statement in \eqref{item_condition_one_dct} is best understood as a kind of quantitative strengthening of the combinatorial characterization of sets of topological recurrence given at the beginning of this section.

The following is an analogue of \cref{mainthm_dct_combo_characterizations} for dynamically thick sets and is proved in \cref{sec_dthick_comb_characterizations}.

\begin{Maintheorem}
\label{mainthm_dt_characterizations}
    Let $A \subseteq \N$.  The following are equivalent.
    \begin{enumerate}
        \item 
        \label{item_intro_dthick_def}
        $A$ is dynamically thick.

        \item 
        \label{item:intro_dyn_thick_idempotent}
        For all $\N \supsetneq B \supseteq A$, there exists a finite set $F \subseteq \N \setminus B$ such that the set $B-F$ is thick.

        \item \label{item:intro_complement_thick_dt_characterization_v2}
        For all piecewise syndetic $S \subseteq \N$, there exists a finite set $F \subseteq A$ such that the set $S-F$ is thick.
    \end{enumerate}
\end{Maintheorem}

In \cite[Question 9.6]{glasner_tsankov_weiss_zucker_2021}, Glasner, Tsankov, Weiss and Zucker call for an explicit characterization of dynamically thick sets (they call them \emph{dense orbit sets}) in countable, discrete groups.  \cref{mainthm_dt_characterizations} gives an answer to their call for subsets of $\N$.  Kennedy, Raum, and Salomon \cite{kennedy_raum_salomon_2022} also provide an answer by translating their characterization of dynamically syndetic sets into one for dynamically thick sets.  Because our characterization of dynamically syndetic sets provides a simpler-to-verify sufficient condition (as discussed in \cref{sec_intro_char_of_ds_sets}), \cref{mainthm_dt_characterizations} reveals more about dynamically thick subsets of $\N$.



The statement in \cref{mainthm_dt_characterizations} \eqref{item:intro_dyn_thick_idempotent} can be used to produce novel examples of dynamically thick sets, such as
\begin{equation}
\label{eq:example_dct_first}
    A \defeq \bigcup_{i=1}^{\infty} \big((p_i \N + 1) \cap H_i \big),
\end{equation}
where $(p_i)_{i \in \N}$ is a sequence of distinct primes and $(H_i)_{i \in \N}$ is a sequence of thick sets. Heuristically, the set $A$ is dynamically thick because it contains sufficiently long pieces of return-time sets from a sufficiently rich collection of systems, in this case, rotations on prime many points.  \cref{thm_bessel} makes this heuristic precise and allows us to generalize the example in \eqref{eq:example_dct_first} by showing that no minimal topological dynamical system can be disjoint from an infinite collection of disjoint systems.  A collection of further examples of dynamically thick sets is given in \cref{sec_dy_thick_examples}.  


The dynamically thick set $A$ in \eqref{eq:example_dct_first} allows us to answer another question posed by Host, Kra, and Maass \cite[Question 3.11]{HostKraMaass2016}. They asked whether every set of pointwise recurrence for distal systems is an $\IP$ set (see \cref{sec_top_dynamics,sec_ip_and_central_sets} for the requisite definitions).  If the thick sets $H_i$ are chosen to be sufficiently spaced (this is clarified in \cref{rem:combinatorial_primes}), the set $A$ is not an $\IP$ set. It is still, however, dynamically thick -- and, hence, a set of pointwise recurrence -- giving a negative answer to the question.  (Their question was first answered in \cite{Koutsogiannis_Le_Moreira_Pavlov_Richter_interpolation_sets}, where an example of a dynamically thick set for distal systems that does not contain the configuration $\{x,y,x+y\}$ is shown.)



\subsubsection{Dynamical thickness in the context of disjointness}
\label{sec_intro_disjointness}

The notions of pointwise recurrence and dynamical thickness are linked to the topics of product recurrence and disjointness, respectively.  We will point out where these topics have appeared together in the literature and what our results have to say about them.

A point $x$ in a dynamical system $(X,T)$ is \emph{recurrent} if for all open $U \subseteq X$ containing $x$, the set $R_T(x,U)$ is nonempty.  The following, then, is immediate from the definitions: for all open $U \subseteq X$ containing $x$, the set $R_T(x,U)$ is a set of pointwise recurrence if and only if for all minimal system $(Y,S)$ and all $y \in Y$, the point $(x,y)$ is recurrent in the product system $(X \times Y, T \times S)$. Dong, Shao, and Ye \cite{dong_shao_ye_2012} call points $x$ with this property \emph{$\mathcal{F}_s$-PR points}.  Distal points (see \cref{example_dcps_star_point_returns}) are examples of $\mathcal{F}_s$-PR points, as are transitive points in the full shift (see below).

Given two dynamically central syndetic sets $A$ and $B$, there exist minimal systems $(X,T)$ and $(Y,S)$, points $x \in X$ and $y \in Y$, and open sets $U \subseteq X$ and $Y \subseteq Y$ containing $x$ and $y$, respectively, such that $R_T(x,U) \subseteq A$ and $R_S(y,V) \subseteq B$.  If $A$ and $B$ are disjoint, then $R_T(x,U) \cap R_S(y,V) = \emptyset$. In this case, the point $(x,y)$ is not recurrent in the product system $(X \times Y, T \times S)$.  We show in the following theorem, proved in \cref{sec_partitioning_dcs_sets}, that such pairs of sets exist in abundance.


\begin{Maintheorem}
\label{mainthm_partition_of_dcs_sets}
    Every dynamically central syndetic set can be partitioned into infinitely many, pairwise disjoint dynamically central syndetic sets.
\end{Maintheorem}

While \cref{mainthm_partition_of_dcs_sets} concerns dynamically central syndetic sets, its consequences for the family of sets of pointwise recurrence are most relevant to this discussion.  Host, Kra, and Maass \cite[Question 6.5]{HostKraMaass2016} asked whether or not the family of sets of pointwise recurrence is partition regular.
A strong negative answer, given in \cref{thm_partition_of_N_into_disjoint_non_sets_of_rec}, follows immediately from \cref{mainthm_partition_of_dcs_sets}: if $\N = A \cup B$, where $A$ and $B$ are disjoint, dynamically central syndetic sets, then $\N$ has been partitioned into two sets, neither of which is a set of pointwise recurrence. It is well-known that sets of topological and measurable recurrence are partition regular, so \cref{mainthm_partition_of_dcs_sets} offers yet another contrast between the nature of set and pointwise recurrence.\\


A \emph{joining} of two topological dynamical systems $(X,T)$ and $(Y,S)$ is a nonempty, closed set $Z \subseteq X \times Y$ that satisfies: $(T \times S) Z \subseteq Z$, $\pi_X Z = X$, and $\pi_Y Z = Y$, where $\pi_X, \pi_Y$ are the projections from $X \times Y$ to $X$ and $Y$, respectively. Furstenberg \cite[Thm. II.2]{Furstenberg-disjointness} defined two topological systems $(X,T)$ and $(Y,S)$ to be \emph{disjoint} if their only joining is $X \times Y$.

To relate the notions of disjointness and dynamical thickness, let $(X,T)$ be a system and $x \in X$ be a point with a dense orbit.  The following is simple to verify from the definitions: the system $(X,T)$ is disjoint from every minimal system if and only if for all nonempty, open $U \subseteq X$, the set $R_T(x,U)$ is dynamically thick.  It is interesting to compare this to the characterization of recurrent points in the product $X \times Y$ given above.

Furstenberg proved \cite[Thm. II.2]{Furstenberg-disjointness} that the full shift $(\{0,1\}^{\N}, \sigma)$ is disjoint from every minimal system.  This result spurred a number of works in which the notion of dynamical thickness arises either directly or indirectly \cite{dong_shao_ye_2012, glasner_tsankov_weiss_zucker_2021, Huang_Shao_Ye_2019, huang_ye_2005,kennedy_raum_salomon_2022}.  We provide a quick explanation of Furstenberg's result in \cref{rmk_transitive_point_visit_times} by combining the characterization in the previous paragraph with the example in \eqref{eq:example_dct_first} above.



This connection between dynamical thickness and disjointness; examples of dynamically thick sets such as the one in \eqref{eq:example_dct_first} and its generalizations in \cref{sec_dy_thick_examples}; and the call from \cite[Question 9.6]{glasner_tsankov_weiss_zucker_2021} led us to wonder whether or not there is a ``structure theorem'' for dynamically thick sets. Drawing on the results from Part III of this paper, we are able to offer the following description of an arbitrary dynamically thick set, proved in \cref{sec_general_form_of_dthick}.  A collection $\collection$ of subsets of $\N$ is \emph{robustly syndetic} if for all dynamically syndetic sets $A \subseteq \N$, there exists $B \in \collection$ such that the set $A \cap B$ is syndetic.

\begin{Maintheorem}
\label{mainthm_structure_for_dy_thick_sets}
\label{cor_weak_structure_result}
    Let $A \subseteq \N$.  The set $A$ is dynamically thick if and only if there exist a robustly syndetic collection $\collection$ of subsets of $\N$ and, for each $B \in \collection$, a thick set $H_B \subseteq \N$ such that
\begin{align}
\label{eqn_A_contains_family_cap_thick}
    A = \bigcup_{B \in \collection} \big( B \cap H_B\big).
\end{align}
\end{Maintheorem}

In all of the examples of dynamically thick sets we know, there is much more that can be said about the robustly syndetic collection $\collection$.  More needs to be said: we could not even determine using \cref{mainthm_structure_for_dy_thick_sets} whether or not a dynamically thick set can be partitioned into two dynamically thick sets.  We speculate more precisely in \cref{sec_dy_thick_sets} on how \cref{mainthm_structure_for_dy_thick_sets} might be improved.




\subsection{Part III: Dynamically piecewise syndetic sets}

In the last part of this paper, we introduce and study the family of sets formed by intersections of dynamically (central) syndetic and dynamically thick sets (sets of pointwise recurrence).  
For sets $B$ and $C$ of positive integers, the set $B \cap C$ is called 
\begin{itemize}
    \item \emph{dynamically piecewise syndetic} if $B$ is dynamically syndetic and $C$ is dynamically thick, and

    \item \emph{dynamically central piecewise syndetic} if $B$ is dynamically central syndetic and $C$ is a set of pointwise recurrence.
\end{itemize}
To our knowledge, the families of dynamically (central) piecewise syndetic sets have not appeared before in the literature.  As we discuss next, their study is motivated by the well-established notions of piecewise syndeticity and centrality.


\subsubsection{Motivation: combinatorial and dynamical piecewise syndeticity}

Syndeticity and thickness are \emph{dual}: for every syndetic set $A$ and every thick set $B$, the intersection $A \cap B$ is nonempty.  As mentioned in \cref{sec_intro_set_and_pt_recurrence}, sets of this form are called piecewise syndetic.  Piecewise syndeticity is a ``local'' analogue of syndeticity that has the advantage of being partition regular: if the union of two sets is piecewise syndetic, one of the sets must be piecewise syndetic.  The family of piecewise syndetic sets and the subfamily of central sets (mentioned in \cref{sec_intro_algebra_of_filters}) are important to the connection between combinatorics and dynamics (via correspondence principles) and to describing algebraic structure in the Stone-\Cech{} compactification of $\N$.

The families of dynamically syndetic sets and dynamically thick sets are dual in the same sense, motivating the study of their intersections.  By definition, then, dynamically piecewise syndetic sets are nonempty, but what more can be said?  The same question applies equally well to the family of dynamically central piecewise syndetic sets.  What are these families' relationships to each other and to the families of piecewise syndetic and central sets?  In addressing these questions, we generalize some of the aforementioned results on dynamically syndetic and thick sets in the literature and and give some novel combinatorial applications.



\subsubsection{Dynamically piecewise syndetic sets are piecewise syndetic}

Because the families of dynamically (central) piecewise syndetic sets have not been considered before, our first results concern the relationships between these families and with the families of piecewise syndetic and central sets. Looking ahead, \cref{fig:containmentdiagram} at the end of \cref{sec_prelims} summarizes the relationships between the families studied in this paper with the established families of subsets of $\N$ in the topic.

There are dynamically syndetic sets that are not dynamically central syndetic, and there are sets of pointwise recurrence that are not dynamically thick.  Therefore, from the definitions, there is no apparent relationship between the families of dynamically piecewise syndetic and dynamically central piecewise syndetic sets.  As our terminology suggests, however, we prove that dynamically central piecewise syndetic sets are dynamically piecewise syndetic.

We could not find a proof of this fact that does not also elucidate the relationship between the families of dynamically piecewise syndetic and piecewise syndetic sets.  We show that, perhaps surprisingly, these families are the same.  The following theorem summarizes these relationships and is proved in \cref{sec_dcps_relations}.


\begin{Maintheorem}
\label{thm:relations_C_dPS_PS-intro}
    Denote by $\central$, $\PS$, $\dPS$, and $\dcPS$ the families of central, piecewise syndetic, dynamically piecewise syndetic, and dynamically central piecewise syndetic sets, respectively.  The following inclusions hold:
    \[
        \central \subsetneq \dcPS \subsetneq \dPS = \PS.
    \]
    In particular, a set is dynamically piecewise syndetic if and only if it is piecewise syndetic.
\end{Maintheorem}

That dynamically piecewise syndetic sets are piecewise syndetic offers a striking contrast to the analogy with set recurrence given in \cref{sec_intro_set_and_pt_recurrence}.  Indeed, sets of measurable recurrence -- and thus, in particular, the intersection of a set of measurable recurrence with the set of times that a particular set of positive measure recurs -- are not necessarily even of positive Banach density (and, hence, are not piecewise syndetic).  By a result of Furstenberg \cite{furstenberg_book_1981} and S\'{a}rk\"{o}zy \cite{sarkozy_1978}, the set of perfect squares, for example, is such a set of measurable recurrence.

We discussed in \cref{sec_intro_set_and_pt_recurrence} how the results in \cite[Prop. 4.4]{dong_shao_ye_2012}, \cite{Glasner-Weiss-Interpolation}, and \cite[Thm. 2.4]{huang_ye_2005} show that the families of dynamically thick sets and sets of pointwise recurrence are piecewise syndetic.  \cref{thm:relations_C_dPS_PS-intro} generalizes these results by showing that the families of dynamically (central) piecewise syndetic sets are piecewise syndetic.  It follows by the abstract family results in \cref{lemma_class_cap_is_pr} that a set of pointwise recurrence modified on a non-piecewise syndetic set is still a set of pointwise recurrence.  The same is true for all of the new families described in this paper; see \cref{cor:largest_family_satisfies_cap}.

There are two keys to the proof of \cref{thm:relations_C_dPS_PS-intro}. The first is our characterization of dynamically central syndetic sets as members of syndetic, idempotent filters in \cref{intro_mainthm_full_chars_of_dcsyndetic}, and the second is a collection of abstract results on large, idempotent filters (see, eg., \cref{prop:general_abstract_intersection_filter_nonsense}).

\subsubsection{Sets of polynomial recurrence and combinatorial configurations in dynamically central piecewise syndetic sets}

A set $A \subseteq \N$ is called a \emph{set of multiple topological recurrence} if for all $k \in \N$, all minimal systems $(X, T)$, and all nonempty, open sets $U \subseteq X$, there exists $n \in A$ such that
\[
    U \cap T^{-n} U \cap \cdots \cap T^{-kn} U \neq \emptyset.
\]
As with the sets of topological recurrence defined in \cref{sec_intro_set_and_pt_recurrence}, understanding which sets are sets of multiple topological recurrence leads to various applications, for example, generalizations of the classical van der Waerden's theorem \cite{van_der_waerden_1927}. While it is known that a set of topological recurrence is not necessarily a set of multiple topological recurrence \cite{frantzikinakis_lesigne_wierdl_2006}, Host, Kra, and Maass \cite[Question 2.14]{HostKraMaass2016} asked whether or not a set of pointwise recurrence is a set of multiple topological recurrence.



We show in the next theorem a strong positive answer: dynamically central piecewise syndetic sets (and, hence, sets of pointwise recurrence) are sets of polynomial multiple measurable recurrence for commuting transformations.  A set $A \subseteq \N$ is a \emph{set of polynomial multiple measurable recurrence for commuting transformations} if for all $k, \ell \in \N$, all commuting measure preserving systems $(X, \mu, T_1, \ldots, T_k)$ (see \cref{sec_dcps_and_set_recurrence}), all $E \subseteq X$ with $\mu(E) > 0$, and, all polynomials $p_{i,j} \in \Q[x]$ with $p_i(0) = 0$ and $p_i(\N_0) \subseteq \N_0$, $1 \leq i \leq k$, $1 \leq j \leq \ell$, there exists $n \in A$ such that
\[
    \mu\Big( \big(T_1^{p_{1,1}(n)} T_2^{p_{2,1}(n)} \cdots T_k^{p_{k,1}(n)} \big)^{-1} E \cap \cdots \cap \big(T_1^{p_{1,\ell}(n)} T_2^{p_{2,\ell}(n)} \cdots T_k^{p_{k,\ell}(n)} \big)^{-1} E \Big) > 0.
\]
Because any minimal topological dynamical system admits an invariant measure of full support, sets of polynomial multiple measurable recurrence for commuting transformations are sets of multiple topological recurrence. Therefore, the following theorem answers \cite[Question 2.14]{HostKraMaass2016} positively.


\begin{Maintheorem}
\label{thm:dcps_multiple_recurrence}
    Dynamically central piecewise syndetic sets are sets of polynomial multiple measurable recurrence for commuting transformations.
\end{Maintheorem}

The proof of \cref{thm:dcps_multiple_recurrence} in \cref{sec_dcps_and_set_recurrence} appeals to the $\IP$ polynomial \Szemeredi{} theorem of Bergelson and McCutcheon \cite{Bergelson-McCutcheon-IPpolynomialSzemeredi} to show that the family of times of polynomial returns in ergodic theory is an $\IP^*$, idempotent filter (see \cref{lemma_localization_is_syndetic_idempotent_filter}). \cref{intro_mainthm_full_chars_of_dcsyndetic} is then invoked to prove that every member of this $\IP^*$, idempotent filter intersects all dynamically central piecewise syndetic sets.


Combining \cref{thm:dcps_multiple_recurrence} with the fact that dynamically central piecewise syndetic sets are piecewise syndetic (\cref{thm:relations_C_dPS_PS-intro}), we show in  \cref{thm:Brauer_and_dcPS} that every such set contains Brauer-type polynomial configurations.  The extent to which dynamically central piecewise syndetic sets are combinatorially rich remains an interesting open question that we speculate on in \cref{sec_open_quests_dps_sets}.


\subsection{Organization of the paper}

Notation, terminology, and preliminaries are given in \cref{sec_prelims}, including material on the algebra of families of subsets of $\N$.  The main body of work in this paper is then divided into three parts, which comprise Sections \ref{sec_dsyndetic}, \ref{sec_dthick}, and \ref{sec_dps_sets}, respectively:
\begin{itemize}
    \item[--] Part I: Dynamically syndetic sets, in which we prove Theorems \ref{mainthm_ds_equivalents},  \ref{intro_mainthm_full_chars_of_dcsyndetic}, and \ref{mainthm_partition_of_dcs_sets};
    \item[--] Part II: Dynamically thick sets, in which we prove Theorems \ref{mainthm_dct_combo_characterizations}, \ref{mainthm_dt_characterizations}, and \ref{mainthm_structure_for_dy_thick_sets};
    \item[--] Part III: Dynamically piecewise syndetic sets, in which we prove Theorems \ref{thm:relations_C_dPS_PS-intro} and \ref{thm:dcps_multiple_recurrence}.
\end{itemize}
We conclude the paper with a collection of open problems and directions in \cref{sec:open_questions}.  A map of the main results and their logical dependencies is drawn in \cref{fig_main_results}.



\begin{figure}[omitted: a diagram of this paper's results, not reproduced here]\label{fig_main_results}\end{figure}




\subsection{Acknowledgments}

Thanks goes to Andy Zucker and Josh Frisch for helpful discussions at the 2024 Southeastern Logic Symposium Conference and for the references they provided.  Thanks also goes to Neil Hindman for correspondence regarding translates of central sets.  The authors are indebted to Eli Glasner, John Johnson, and Mauro Di Nasso for their assistance in correcting and improving several of the historical references and claims made in a first draft of the paper. Last but not least, thanks go to Andreas Koutsogiannis, Joel Moreira, Florian Richter, and Donald Robertson, with whom the authors were working when a number of the ideas in this paper sharpened. 

\section{Preliminaries}
\label{sec_prelims}


The set of integers, non-negative integers, and positive integers are denoted by $\Z$, $\N_0$, and $\N$, respectively.  The power set of $\N$ is denoted $\mathcal{P}(\N)$.  For $N \in \N_0$, we define $[N]$ to be $\{0, 1, \ldots, N - 1\}$.

For $A \subseteq \N$ and $n \in \N$, we define
\begin{align*}
    A+n &\defeq \{a + n \ | \ a \in A \}, & nA & \defeq \{ na \ | \ a \in A\}, \\
    A-n &\defeq \{m \in \N \ | \ m + n \in A \}, & A/n & \defeq \{ m \in \N \ | \ mn \in A\}.
\end{align*}
For $A, B \subseteq \N$, we define
\begin{align*}
    A+B &\defeq \bigcup_{b \in B} (A + b), & A-B &\defeq \bigcup_{b \in B} (A-b).
\end{align*}
We follow the convention that the empty union of subsets of $\N$ is empty while the empty intersection of subsets of $\N$ is equal to $\N$. Note that all set operations on subsets of $\N$ result in subsets of $\N$.




\subsection{Families of positive integers}
\label{sec_combinatorics}

Let $\class$ be a collection of subsets of $\N$.  The \emph{upward closure} of $\class$ is
\[\upclose \class \defeq \big\{ B \subseteq \N \ \big| \ \exists A \in \class, \ B \supseteq A \big\}.\]
The set $\class$ is \emph{upward-closed} if $\class = \upclose \class$.  We will call an upward-closed collection of subsets of $\N$ a \emph{family}.  Note that both $\emptyset$ and $\mathcal{P}(\N)$ are families.  A family is \emph{proper} if it is neither $\emptyset$ nor $\mathcal{P}(\N)$, that is, it is both nonempty and does not contain the empty set.

\subsubsection{Dual families and intersections}
\label{sec_dual_families}

Much of the algebra described in this subsection and the next goes back to work of Choquet \cite{choquet_1947} and Schmidt \cite{schmidt_1952,schmidt_1953}.  A good modern reference for this material can be found in \cite[Section 2]{christopherson_johnson_2022}.

The \emph{dual} of a family $\family$ of subsets of $\N$ is
\[\family^* \defeq \big\{ B \subseteq \N \ \big | \ \forall A \in \family, \ B \cap A \neq \emptyset \big\},\]
the collection of all those sets which have nonempty intersection with all elements of $\family$.  The following facts are quick to check and will be used without mention:
\begin{enumerate}
    \item $(\family^*)^* = \family$, whereby it makes sense to say that $\family$ and $\family^*$ are \emph{dual};
    \item the empty family $\emptyset$ and $\mathcal{P}(\N)$ are dual, while the dual of a proper family is a proper family;
    \item $A \in \family$ if and only if $\N \setminus A \not\in \family^*$;
    \item $\family \subseteq \familytwo$ if and only if $\familytwo^* \subseteq \family^*$.\\
\end{enumerate}

For nonempty families $\familyone$ and $\familytwo$, we define
\[\familyone \familycap \familytwo \defeq \big\{ A \cap B \ \big| \ A \in \familyone \text{ and } B \in \familytwo \big\}.\]
We will adopt the convention that $\emptyset \classcap \family = \family \classcap \emptyset = \family$ for all families $\family$.  The following facts are quick to check for all families $\family$ and $\familytwo$:
\begin{enumerate}
\setcounter{enumi}{4}
    \item the collection $\family \classcap \familytwo$ is a family that contains both $\family$ and $\familytwo$ (despite what is suggested by the notation);
    \item if $\familyone$ is a proper family, then $\familyone \classcap \familyone^*$ is a proper family;
    \item if $\family_1 \subseteq \familytwo_1$ and $\family_2 \subseteq \familytwo_2$, then $\family_1 \classcap \family_2 \subseteq \familytwo_1 \classcap \familytwo_2$.
\end{enumerate}


\subsubsection{Filters, partition regularity, and ultrafilters}
\label{sec_filter_pr_and_ufs}

There are two complementary notions of largeness for families that feature prominently in this topic: that of being partition regular and that of being a filter.  A family $\family$ is \emph{partition regular} if for all $A \in \family$ and all finite partitions $A = A_1 \cup \cdots \cup A_k$, some piece $A_i$ belongs to $\family$.  The family $\family$ is a \emph{filter} if for all $A_1, A_2 \in \family$, the set $A_1 \cap A_2 \in \family$.  These notions are dual in the sense of fact \eqref{item_pr_filter_duality} below.



If $P$ is a property of families, we call a family that has property $P$ and that is also a filter a \emph{$P$ filter}. For example, we will consider proper filters, partition regular filters, translation-invariant filters, syndetic filters, and so on. Note, for example, that $\emptyset$ and $\mathcal{P}(\N)$ are both partition regular filters.



The following facts are quick to check:
\begin{enumerate}
    \item \label{item_pr_filter_duality} a family $\family$ is partition regular if and only if the dual family $\family^*$ is a filter;
    \item if $\family$ is a proper filter, then $\family \subseteq \family^*$;
    \item \label{item_filter_containment} if $\familyone$ and $\familytwo$ are filters, then $\familyone \classcap \familytwo$ is the smallest filter containing both $\familyone$ and $\familytwo$; in particular, if $\family \subseteq \familytwo$, then $\family \classcap \familytwo = \familytwo$.\\ 
\end{enumerate}

Given a family $\family$, we will frequently consider the family $\family \classcap \family^*$.  The follow lemma records some useful facts about it.

\begin{lemma}
\label{lemma_class_cap_is_pr}
    Let $\family$ be a family and define $\familytwo \defeq \family \classcap \family^*$.  The family $\familytwo$ is partition regular and its dual, $\familytwo^*$, is a filter.  Moreover, for $A \subseteq \N$, the following are all equivalent to the set $A$ belonging to $\familytwo^*$:
    \begin{enumerate}
        \item \label{item_a_cap_b_is_in_family} for all $B \in \family$, the set $A \cap B \in \family$;
        \item for all $B \in \family^*$, the set $A \cap B \in \family^*$;
        \item for all $B \in \familytwo$, the set $A \cap B \in \familytwo$;
        \item for all $B \in \familytwo^*$, the set $A \cap B \in \familytwo^*$.
    \end{enumerate}
\end{lemma}
\begin{proof}
    If $\family$ is not proper, then $\familytwo = \mathcal{P}(\N)$, which is partition regular.  Otherwise, the partition regularity of $\familytwo$ is shown in \cite[Prop. 2.5 (h)]{christopherson_johnson_2022}.  That its dual, $\familytwo^*$, is a filter is mentioned as item \eqref{item_pr_filter_duality} before the statement of this lemma.

    We will show that \eqref{item_a_cap_b_is_in_family} is equivalent to $A$ belonging $\familytwo^*$; the other statements follow analogously and are left to the reader.  Suppose $A \in \familytwo^*$ and $B \in \family$.  To see that $A \cap B \in \family$, we will show that for all $C \in \family^*$, the set $A \cap B \cap C$ is nonempty.  Let $C \in \family^*$.  The set $B \cap C \in \familytwo$, and since $A \in \familytwo^*$, we have that $A \cap B \cap C \neq \emptyset$, as desired.
    
    Conversely, suppose that \eqref{item_a_cap_b_is_in_family} holds.  To see that $A \in \familytwo^*$, we will show that for all $D \in \familytwo$, the set $A \cap D \neq \emptyset$.  Let $D \in \familytwo$.  There exists $B \in \family$ and $C \in \family^*$ such that $D = B \cap C$.  Then $A \cap D = A \cap B \cap C$.  By assumption, $A \cap B \in \family$, whereby $ A \cap B \cap C \neq \emptyset$, as desired.
\end{proof}


A proper family is called an \emph{ultrafilter} if it satisfies one (equivalently, all) of the properties in the following lemma.  This lemma is well-known (cf. \cite[Thm. 3.6]{hindman_strauss_book_2012}), but its proof is short, so we provide it here.

\begin{lemma}
\label{lemma_equiv_ultrafilter}
    Let $\family$ be a proper family of subsets of $\N$.  The following are equivalent:
    \begin{enumerate}
        \item
        \label{item_family_is_maximal}
        the family $\family$ is a proper filter and is maximal (by containment) amongst proper filters;
        
        \item
        \label{item_family_is_self_dual}
        the family $\family$ is contained in a proper filter and is self-dual, ie., $\family = \family^*$;
        
        \item
        \label{item_family_is_pr_and_filter}
        the family $\family$ is partition regular and a filter.
    \end{enumerate}
\end{lemma}

\begin{proof}
    \eqref{item_family_is_maximal} $\implies$ \eqref{item_family_is_self_dual} \ Suppose $\family$ is a proper filter and is maximal amongst proper filters.  Since $\family$ is a proper filter, we have that $\family \subseteq \family^*$.  Let $B \in \family^*$.  It is quick to check that $\familytwo \defeq \upclose \{A \cap B \ | \ A \in \family \}$ is a proper filter that contains $\family$.  Since $\family$ is maximal, we have that $\family = \familytwo$, whereby $B \in \family$.  Since $B \in \family^*$ was arbitrary, we have that $\family = \family^*$.

    \eqref{item_family_is_self_dual} $\implies$ \eqref{item_family_is_maximal} \ Suppose that $\family$ is contained in a proper filter and is self-dual. To conclude \eqref{item_family_is_maximal}, it suffices to show that for all proper filters $\familytwo$ containing $\family$, we have $\familytwo = \family$.
    
    Let $\familytwo$ be a proper filter containing $\family$. Since $\familytwo$ is a proper filter, we have that $\familytwo \subseteq \familytwo^*$.  Since $\family \subseteq \familytwo$, we have that $\familytwo^* \subseteq \family^* = \family \subseteq \familytwo \subseteq \familytwo^*$.  Therefore, $\family = \familytwo$, as desired.

    \eqref{item_family_is_self_dual} $\implies$ \eqref{item_family_is_pr_and_filter} \ Suppose that $\family$ is contained in a proper filter and is self-dual.  It is shown in the proof of \eqref{item_family_is_self_dual} implies \eqref{item_family_is_maximal} that $\family$ is a proper filter. It follows from fact \eqref{item_pr_filter_duality} from the beginning of this section, then, that $\family^*$ is partition regular.  By assumption, $\family^* = \family$, hence $\family$ is partition regular, as desired.

    \eqref{item_family_is_pr_and_filter} $\implies$ \eqref{item_family_is_self_dual} \ Suppose $\family$ is a proper filter that is partition regular.  Since $\family$ is proper, we have that $\family \subseteq \family^*$.  Since $\family$ is proper and partition regular, the family $\family^*$ is a proper filter, whereby $\family^* \subseteq (\family^*)^* = \family$.  Thus, $\family = \family^*$, as desired.
\end{proof}

Ultrafilters exist by a standard application of Zorn's Lemma.  The space of ultrafilters on $\N$ is an important object that we will make use of in this paper.  We delay further discussion until \cref{sec_topology_and_ultrafilters}.



\subsubsection{Translation algebra}
\label{sec_family_algebra}

For $A \subseteq \N$ and families $\family$ and $\familytwo$ of subsets of $\N$, we define
\begin{align*}
    A - \family &\defeq \big\{ n \in \N \ \big| \ A-n \in \family \big\},\\
    \familyone + \familytwo &\defeq \big\{ B \subseteq \N \ \big| \ B - \familytwo \in \family \big\}.
\end{align*}
Thus, we see that $n \in A - \family$ if and only if $A - n \in \family$ and that $B \in \family + \familytwo$ if and only $B - \familytwo \in \family$.  It is simple to check that $\family + \familytwo$ is a family of subsets of $\N$.

The definition of $\family + \familytwo$ appears for filters in \cite{berglund_hindman_1984,papazyan_1989} and, when $\family$ and $\familytwo$ are ultrafilters, agrees with the usual definition of sums of ultrafilters (cf. \cite[Thm. 4.12]{hindman_strauss_book_2012}).  Family sums will appear throughout the paper.  It is useful to note that, loosely speaking, the containment $\family \subseteq \familytwo + \familythree$ means: for all $A \in \family$, there are $\familytwo$ many positive integers $n$ for which the set $A-n$ belongs to $\familythree$.

The following lemma records a number of algebraic facts that will be useful later on.

\begin{lemma}
    \label{lemma_filter_algebra}
    \label{lemma_fam_dual_sums}
    \label{lemma_filter_sums_containment}
    Let $n \in \N$, $A, B \subseteq \N$, and $\filter, \filtertwo, \filter_1, \filter_2, \filtertwo_1, \filtertwo_2$ be families.
    \begin{enumerate}
        \item
        \label{item_family_prop_zero}
        If $\filter$ and $\filtertwo$ are proper, then $\filter + \filtertwo$ is proper.
        
        \item
        \label{item_family_prop_one_new}
        $(A - \filter) \cap (B - \filtertwo) \subseteq (A \cap B) - (\filter \classcap \filtertwo)$, with equality if $\filter = \filtertwo$ is a filter.

        \item \label{item_family_prop_two}
        If $\filter$ and $\filtertwo$ are filters, then $\filter + \filtertwo$ is a filter.

        \item \label{item_family_prop_three}
        $(A -n) - \filter = (A - \filter) - n$.

        \item \label{item_family_prop_four}
        $A - (\filter + \filtertwo) = (A - \filtertwo) - \filter$.

        \item \label{item_family_diff_containment}
        If $\family \subseteq \familytwo$, then $A - \family \subseteq A - \familytwo$.
        
        \item \label{item_filter_sums_containment}
        If $\filter_1 \subseteq \filtertwo_1$ and $\filter_2 \subseteq \filtertwo_2$, then $\filter_1 + \filter_2 \subseteq \filtertwo_1 + \filtertwo_2$.

        \item \label{item_sum_dual}
        $(\familyone + \familytwo)^* = \familyone^* + \familytwo^*$.

        \item \label{item_cap_sums}
        $(\familyone_1 + \familyone_2) \classcap (\familytwo_1 + \familytwo_2) \subseteq (\familyone_1 \classcap \familytwo_1) + (\familyone_2 \classcap \familytwo_2)$.
    \end{enumerate}
\end{lemma}

\begin{proof}
    \eqref{item_family_prop_zero} \ Suppose $\family$ and $\familytwo$ are proper.  Since $\emptyset \not\in \familytwo$, we have $\emptyset - \familytwo = \emptyset$. Since $\emptyset \not\in \family$, we have $\emptyset \not\in \family + \familytwo$. Since $\family$ and $\familytwo$ both contain $\N$, we see that $\N \in \family + \familytwo$.  Since $\family + \familytwo$ does not contain $\emptyset$ but contains $\N$, it is proper.
    
    \eqref{item_family_prop_one_new} \  Suppose $n \in (A - \filter) \cap (B - \filtertwo)$.  Since $A - n \in \filter$ and $B - n \in \filtertwo$, we have that $(A \cap B) - n = (A-n) \cap (B-n) \in \filter \classcap \filtertwo$.  Thus, $n \in (A \cap B) - (\filter \classcap \filtertwo)$, as desired.

    If $\family = \familytwo$ is a filter, then $\family \classcap \familytwo = \family$.  We wish to show in this case that $(A \cap B) - \family \subseteq (A- \family) \cap (B-\family)$.  Suppose $n \in (A \cap B) - \family$ so that $(A - n) \cap (B-n) \in \family$.  Since $\family$ is upward closed, we have that $A -n \in \family$ and $B-n \in \family$.  Therefore, $n \in (A - \family) \cap (B-\family)$, as desired.

    \eqref{item_family_prop_two} \ Suppose $\filter$ and $\filtertwo$ are filters.  Since $\filter + \filtertwo$ is a family, we need only to show that for all $A, B \in \filter + \filtertwo$, the set $A \cap B \in \filter + \filtertwo$.  Let $A, B \in \filter + \filtertwo$.  Thus, $A - \filtertwo, B - \filtertwo \in \filter$, whereby $(A-\filtertwo) \cap (B-\filtertwo) \in \filter$.  By \eqref{item_family_prop_one_new}, it follows that $(A \cap B) - \filtertwo \in \filter$, whereby $A \cap B \in \filter + \filtertwo$, as desired.

    \eqref{item_family_prop_three} \ Note that $m \in (A -n) - \filter$ if and only if $A-n-m \in \filter$, which happens if and only if $A-(m+n) \in \filter$.  We see that this happens if and only if $m+n \in A-\filter$, that is, if and only if $m \in (A-\filter) - n$, as desired.

    \eqref{item_family_prop_four} \ We see that $n \in A - (\filter + \filtertwo)$ if and only if $A-n \in \filter + \filtertwo$, which happens if and only if $(A-n) - \filtertwo \in \filter$.  By \eqref{item_family_prop_three}, this happens if and only if $(A-\filtertwo) - n \in \filter$, which happens if and only if $n \in (A - \filtertwo) - \filter$, as desired.

    \eqref{item_family_diff_containment} \ If $n \in A - \family$, then $A - n \in \family \subseteq \familytwo$, whereby $n \in A - \familytwo$.
    
    \eqref{item_filter_sums_containment} \ Suppose $A \in \filter_1 + \filter_2$.  Then $A-\filter_2 \in \filter_1 \subseteq \filtertwo_1$.  Since $\filter_2 \subseteq \filtertwo_2$, we have that $A - \filter_2 \subseteq A - \filtertwo_2$.  Since $\filtertwo_1$ is upward closed and $A- \filter_2 \in \filtertwo_1$, we see that $A - \filtertwo_2 \in \filtertwo_1$.  Therefore, $A \in \filtertwo_1 + \filtertwo_2$, as desired.
    
    \eqref{item_sum_dual} We will show first that $\familyone^* + \familytwo^* \subseteq (\familyone + \familytwo)^*$.  Suppose $A \in \familyone^* + \familytwo^*$, and let $B \in \familyone + \familytwo$.  Since $A - \familytwo^* \in \familyone^*$ and $B - \familytwo \in \familyone$, we have that there exists $n \in (A - \familytwo^*) \cap (B - \familytwo)$.  We see that $A - n \in \familytwo^*$ and $B - n \in \familytwo$.  It follows that $(A-n) \cap (B-n) \neq \emptyset$, whereby $A \cap B \neq \emptyset$, as desired.

    Now we can write
    \[\familyone^* + \familytwo^* \subseteq (\familyone + \familytwo)^* = \big((\familyone^*)^* + (\familytwo^*)^* \big)^* \subseteq \big((\familyone^* + \familytwo^*)^* \big)^* = \familyone^* + \familytwo^*,\]
    where the first containment follows from the previous paragraph and the second containment follows from the previous paragraph applied to the dual classes.  This shows the desired equality.

    \eqref{item_cap_sums} \ Suppose $\family_1 + \family_2 = \emptyset$ or $\familytwo_1 + \familytwo_2 = \emptyset$.  If $\family_1 + \family_2 = \emptyset$, then either $\family_1 = \emptyset$ or $\family_2 = \emptyset$.  In the first case, the desired inclusion simplifies to $\familytwo_1 + \familytwo_2 \subseteq \familytwo_1 + (\family_2 \classcap \familytwo_2)$, which follows from \eqref{item_filter_sums_containment}.  The second case is shown similarly.  The case that $\familytwo_1 + \familytwo_2 = \emptyset$ is handled analogously.

    Otherwise, we have that $\family_1 + \family_2 \neq \emptyset$ and $\familytwo_1 + \familytwo_2 \neq \emptyset$.  To see the desired inclusion, let $C \in (\familyone_1 + \familyone_2) \classcap (\familytwo_1 + \familytwo_2)$. There exists $A \in \familyone_1 + \familyone_2$ and $B \in \familytwo_1 + \familytwo_2$ such that $C = A \cap B$. It follows from point \eqref{item_family_prop_one_new} that $(A-\family_2) \cap (B-\familytwo_2) \subseteq (A \cap B) - (\family_2 \classcap \familytwo_2)$.  Since $A-\family_2 \in \family_1$ and $B-\familytwo_2 \in \familytwo_1$, we see that that $(A \cap B) - (\family_2 \classcap \familytwo_2) \in \family_1 \classcap \familytwo_1$, whereby $A \cap B \in (\family_1 \classcap \familytwo_1) + (\family_2 \classcap \familytwo_2)$, as desired.
\end{proof}




\subsubsection{Translation-invariant and idempotent families}
\label{sec_trans_inv_and_idempotent_families}

For $m \in \N$ and a family $\family$ of subsets of $\N$, we define
\[\family - m \defeq \big\{ B - m \ \big| \ B \in \family \big\}.\]
The family $\family$ is \emph{translation invariant} if $\family - n \subseteq \family$ for all $n \in \N$. In the algebra of the previous section, this is equivalent to saying that $\family \subseteq \{\N\} + \family$.  The family $\filter$ is \emph{idempotent} if $\filter \subseteq \filter + \filter$, equivalently, if for all $A \in \filter$, $A - \filter \in \filter$. (We comment on why the set inclusion in the definition is not an equality in \cref{rmk_why_set_inclusion_not_equality} below.) Note that all translation-invariant families are idempotent.  We will give some less trivial examples of idempotent families in \cref{sec_families_of_sets}.



\begin{lemma}
\label{lemma_idempotent_classcap_is_idempotent}
    Let $\family$ and $\familytwo$ be families.
    \begin{enumerate}
        \item
        \label{item_class_cap_of_idempotent_is_idempotent}
        If $\family$ and $\familytwo$ are idempotent, then $\family \classcap \familytwo$ is idempotent.

        \item
        \label{item_class_cap_of_trans_inv}
        If $\family$ and $\familytwo$ are translation invariant, then $\family \classcap \familytwo$ is translation invariant.
    \end{enumerate}
\end{lemma}

\begin{proof}
    \eqref{item_class_cap_of_idempotent_is_idempotent} \ We appeal to \cref{lemma_fam_dual_sums} \eqref{item_cap_sums} with $\familyone_1 = \familyone_2 = \familyone$ and $\familytwo_1 = \familytwo_2 = \familytwo$.  Thus,
    \[\familyone \classcap \familytwo \subseteq (\familyone + \familyone) \classcap (\familytwo + \familytwo) \subseteq (\familyone \classcap \familytwo) + (\familyone \classcap \familytwo),\]
    as desired.  Statement \eqref{item_class_cap_of_trans_inv} follows similarly and is left to the reader.
\end{proof}

\begin{remark}
    \label{rmk_showing_idempotency}
    The following situation arises several times in this paper.  A family $\family$ is defined as the upward closure of a collection of special subsets of $\N$.  To show that $\family$ is idempotent, it suffices to show that for all special sets $A \subseteq \N$, the set $A - \filter$ belongs to $\family$.  Indeed, if $B \in \filter$, then $B$ contains some special set $A$.  If we show that $A-\filter \in \filter$, then since $A- \filter \subseteq B - \filter$, we have that $B - \filter \in \filter$, as desired.
\end{remark}




Idempotent ultrafilters are important objects with a deep history in the subject (cf. \cite{bergelson_2010}).  It appears that there are two different definitions of ``idempotent ultrafilter.''  In the literature, a family $\family$ is an idempotent ultrafilter if it is an ultrafilter and satisfies $\family = \family + \family$.  According to our definitions, an idempotent ultrafilter is an idempotent family (ie. $\family \subseteq \family + \family$) that is also an ultrafilter. We show in the following lemma these two definition are, in fact, the same.  We will use facts from \cref{sec_topology_and_ultrafilters}.


\begin{theorem}
\label{thm_equiv_forms_of_max_idempot_filter}
    Let $\family$ be a family.  The following are equivalent:
    \begin{enumerate}
        \item
        \label{item_max_prop_idemp_filter}
        the family $\family$ is a proper, idempotent filter and is maximal (by containment) in the set of proper, idempotent filters;

        \item
        \label{item_ultra_idemp}
        the family $\family$ is an ultrafilter that satisfies $\family \subseteq \family + \family$;

        \item
        \label{item_ultra_strong_idemp}
        the family $\family$ is an ultrafilter that satisfies $\family = \family + \family$.
    \end{enumerate}
    In particular, every proper, idempotent filter is contained in an idempotent ultrafilter.
\end{theorem}

\begin{proof}
    \eqref{item_max_prop_idemp_filter} $\implies$ \eqref{item_ultra_idemp} \ We will prove first that every proper, idempotent filter is contained in an idempotent family that is an ultrafilter.  Supposing this is so, if $\family$ is a maximal proper, idempotent filter, then it is contained in an idempotent family $\familytwo$ that is an ultrafilter.  By the maximality of $\family$ as a proper, idempotent filter, we have that $\family = \familytwo$, giving \eqref{item_ultra_idemp}, as desired.

    Suppose $\family$ is a proper, idempotent filter and is maximal as such. We claim that the set $\overline{\filter} \defeq \bigcap_{A \in \family} \overline{A}$ is a nonempty, compact subsemigroup of $(\beta \N, +)$.  Because ultrafilters are maximal filters, there is an ultrafilter containing $\family$, and so $\overline{\filter}$ is nonempty.   The set $\overline{\filter}$ is compact because it is closed and $\beta \N$ is compact.  Suppose $q_1, q_2 \in \overline{\filter}$.  Since $\filter \subseteq q_1$ and $\filter \subseteq q_2$, it follows from \cref{lemma_filter_sums_containment} \eqref{item_filter_sums_containment} and the idempotency of $\filter$ that
    \[\filter \subseteq \filter + \filter \subseteq q_1 + q_2.\]
    Therefore, we have that $q_1 + q_2 \in \overline{\filter}$.

    Finally, since $\overline{\filter}$ is a nonempty, compact subsemigroup of $(\beta \N, +)$, by \cite[Thm. 2.5]{hindman_strauss_book_2012}, there exists an idempotent ultrafilter containing $\family$.

    \eqref{item_ultra_idemp} $\implies$ \eqref{item_ultra_strong_idemp} \ Suppose $\family$ is an ultrafilter that satisfies $\family \subseteq \family + \family$.  By \cref{lemma_filter_algebra} \eqref{item_family_prop_zero} and \eqref{item_family_prop_two}, the family $\family + \family$ is a proper filter.  Since $\family$ is a maximal proper filter, we have that $\family = \family + \family$, as desired.

    \eqref{item_ultra_strong_idemp} $\implies$ \eqref{item_max_prop_idemp_filter} \ If $\family$ is an ultrafilter that satisfies $\family = \family + \family$, then it is a proper, idempotent filter.  Since it is maximal amongst proper filters, it is maximal amongst proper, idempotent filters.
\end{proof}



\begin{remark}
\label{rmk_idempotents_without_zorn}
    A routine application of Zorn's lemma gives the existence of maximal proper, idempotent filters.  It follows from \cref{thm_equiv_forms_of_max_idempot_filter} that these are ultrafilters.  The proof given here relies on the fact that a compact, Hausdorff, right-topological semigroup contains an idempotent element (see \cite[Thm. 2.5]{hindman_strauss_book_2012} and the end-of-chapter notes). A short, topology-free proof of this was given by Papazyan \cite{papazyan_1989}, and a proof without the full strength of Zorn's Lemma was given by Di Nasso and Tachtsis \cite{dinasso_tachtsis_2018}.

    Another routine application of Zorn's lemma gives the existence of maximal proper, translation-invariant filters.  It is an easy exercise to show that there are no translation-invariant ultrafilters on $\N$.  Therefore, in contrast to maximal proper, idempotent filters, maximal proper, translation-invariant filters are not ultrafilters.
\end{remark}


\begin{remark}
\label{rmk_why_set_inclusion_not_equality}
    The definition of idempotent family we give in this section matches the one for filters given implicitly in \cite{papazyan_1989} and explicitly in \cite{Krautzberger_2009}.  It may seem more natural to require equality in the definition: $\filter = \filter + \filter$.  While certainly some families do satisfy this -- in particular, idempotent ultrafilters (\cref{thm_equiv_forms_of_max_idempot_filter}) -- equality appears to be rather rare.  Here is an example of a dynamically simple idempotent filter $\filter$ which satisfies $\filter \subsetneq \filter + \filter$.  See \cref{sec_top_dynamics} for the notation.
    
    Fix $\alpha \in \R \setminus \Q$, and consider the irrational rotation $(\R / \Z, T_{\alpha}: x \mapsto x + \alpha)$.  Put
    \[\filter \defeq \upclose \big\{ R(0, B_\eps(0)) \ \big| \ \eps > 0\big\}.\]
    The family $\filter$ is an idempotent filter.  This can be easily checked (see, eg., the proof of the ``only if'' statement of \cref{thm_dcs_iff_cs}).

    Let $(k_i)_{i=1}^\infty \subseteq \N$ be such that the only limit point of $i \mapsto k_i \alpha$ is $0$. For $n \in \N$, the quantity
    \[\eps_n \defeq \min \big\{ \|k_i \alpha - n \alpha \|/2 \ \big| \ i \in \N, \ k_i > n \big\},\]
    where $\| \cdot \|: \R / \Z \to [0,1/2)$ denotes the Euclidean distance to $\Z$, is non-zero.
    We claim that the set
    \[A \defeq \bigcup_{n = 1}^\infty \big( R(0, B_{\eps_n}(0) ) + n \big)\]
    belongs to $\filter + \filter$ but not to $\filter$.  Indeed, for all $n \in \N$, the set $A - n$ contains $R(0, B_{\eps_n}(0) ) \in \filter$, and so $A \in \{\N\} + \filter \subseteq \filter + \filter$.

    To see that $A \not\in \filter$, we will show that for all $\delta > 0$, we have $R(0,B_{\delta}(0)) \not\subseteq A$.  Let $\delta > 0$, and let $i \in \N$ such that $k_i \in R(0,B_{\delta}(0))$.  Now, for all $n < k_i$, we have that $\|k_i \alpha - n \alpha\| \geq 2\eps_n$, whereby $k_i \notin R(0,B_{\eps_n}(0)) + n$.  For all $n \geq k_i$, we have that $k_i \notin R(0,B_{\eps_n}(0)) + n$ since $R(0,B_{\eps_n}(0)) + n \subseteq \{n+1, n+2, \ldots\}$.  Therefore, we have that $k_i \not\in A$, as was to be shown.
\end{remark}


\subsection{Topology and dynamics}

For a set $U$ in a topological space $X$, we denote by $\overline U$, $U^\circ$, and $\partial U \defeq \overline{U} \setminus U^\circ$ the closure, interior, and boundary of $U$, respectively.  In a metric space $(X,d)$, we denote by $B_\eps(x)$ the open ball centered at a point $x \in X$ with radius $\eps > 0$.

\subsubsection{Topological dynamics}
\label{sec_top_dynamics}

A \emph{(topological dynamical) system} $(X,T)$ is a nonempty, compact metric space $(X,d_X)$ together with a continuous self-map $T: X \to X$.  A \emph{subsystem of $(X,T)$} is a system of the form $(Y,T)$, where $Y \subseteq X$ is nonempty, closed (equivalently, compact), and \emph{$T$-invariant}, meaning $TY \subseteq Y$.  A system $(X,T)$ whose only subsystem is itself is called \emph{minimal}.  Given two systems $(X,T)$ and $(Y,S)$, a continuous surjection $\pi: X \to Y$ that satisfies $S \circ \pi = \pi \circ T$ is called a \emph{factor map} of systems.  We write $\pi: (X,T) \to (Y,S)$ to indicate that $\pi$ is a factor map.

All systems in this paper will be considered as actions of the semigroup $(\N, +)$ by continuous maps on a compact metric spaces.  Thus, even in the event that $T$ is a homeomorphism (in which case, the system $(X,T)$ is called \emph{invertible}), the emphasis will be on non-negative iterates of the map $T$.

We will occasionally have need to consider non-metrizable systems.  A \emph{not-necessarily-metrizable system} is a nonempty, compact Hausdorff space $X$ together with a continuous self-map $T: X \to X$.  We will always specify when a system need not be metrizable.  All terminology and notation will be shared without confusion between the metric and non-metric cases.

Given a system $(X,T)$, a point $x \in X$, and a set $A \subseteq \N$, we define
\[T^A x \defeq \big\{ T^a x \ \big| \ a \in A \big\}.\]
The \emph{orbit of $x$} is the set $T^{\N} x$ and the \emph{orbit closure of $x$} is the set $\overline{T^{\N}x}$. Note that $\overline{T^{\N}x}$ is a nonempty, closed, $T$-invariant subset of $X$, whereby $(\overline{T^{\N}x}, T)$ is a subsystem of $(X,T)$.  Thus, we see that $(X,T)$ is minimal if and only if every point $x \in X$ has a dense orbit.


Given a system $(X,T)$, a set $U \subseteq X$, and a point $x \in X$, we write
\begin{align*}
    R(x, U) &\defeq \big\{n \in \N \ \big | \ T^n x \in U \big\}
\end{align*}
for the times at which the point $x$ visits the set $U$. It is a simple algebraic fact that we will use without mention that for $n \in \N_0$,
\[R(x,U) - n = R(T^n x, U) = R(x, T^{-n}U).\]


A set $A \subseteq \N$ is \emph{syndetic} if there exists $N \in \N$ such that $A \cup (A-1) \cup \cdots \cup (A-N) = \N$.  Syndetic sets -- discussed in more context in \cref{sec_syndetic_thick_ps_sets} -- are intimately linked to minimal dynamics.  Indeed, let $(X,T)$ be a not-necessarily-metrizable system.  Both of the facts below are not hard to show following the results in \cite[Ch. 1, Sec. 4]{furstenberg_book_1981}.
\begin{enumerate}
    \item Let $x \in X$.  The subsystem $(\overline{T^{\N_0}x}, T)$ is minimal if and only if for all open $U \subseteq X$ containing $x$, the set $R(x,U)$ is syndetic. In this case, the point $x$ is said to be \emph{uniformly recurrent.}

    \item The system $(X,T)$ is minimal if and only if for all $x \in X$ and all nonempty, open $U \subseteq X$, the set $R(x,U)$ is syndetic.
\end{enumerate}
We will use these facts several times in this paper.

\begin{lemma}
    \label{lemma_isolated_points_and_minimal_systems}
    Let $(X,T)$ be a minimal system.  There exists an isolated point in $X$ if and only if $X$ is finite and the system $(X,T)$ is periodic.
\end{lemma}

\begin{proof}
    If $x \in X$ is isolated, the set $U \defeq \{x\}$ is open.  Since $(X,T)$ is minimal, the set $R(x,U)$ is syndetic. In particular, there exists $n \in \N$ such that $T^n x = x$, whereby $x$ is a periodic point.  Since $(X,T)$ is minimal, the set $X$ is equal to the (finite) orbit of $x$ and $(X,T)$ is periodic.  The converse is trivial.
\end{proof}

The dual notions of distality and proximality arise a few times in this paper.  Let $(X,T)$ be a system and $x, y \in X$.  The points $x$ and $y$ are \emph{proximal} if $\inf_{n \in \N} d_X(T^n x, T^n y) = 0$.  The point $x$ is called \emph{distal} if it is proximal only to itself.  The system $(X,T)$ is called \emph{distal} if every point is distal.

\subsubsection{Symbolic space and full shift}
\label{sec:symbolic_prelim}

We denote the space of 0--1 valued functions on $\N$ by $\{0,1\}^{\N}$. Given the product topology, this is a compact Hausdorff space.  There are many equivalent metrics that generate the product topology on $\{0,1\}^{\N}$.  Instead of specifying one of these metrics explicitly, it suffices for our purposes to fix one and note that two elements of $\{0,1\}^{\N}$ are close if and only if they agree as functions on a long initial interval $\{1, \ldots, N\}$ of $\N$.

Let $\omega \in \{0,1\}^{\N}$ and $i \in \N$.  We write $\omega(i)$ for the value of the function $\omega$ at $i$.  Given $k \in \{0,1\}$, the \emph{cylinder set} of all those functions in $\{0,1\}^{\N}$ that evaluate to $k$ at $i$ is denoted by $[k]_i$.  This is easily checked to be a clopen subset of $\{0,1\}^{\N}$.  The \emph{support} of $\omega$ is
\[\supp (\omega) \defeq \big\{ n \in \N \ \big| \ \omega(n) = 1 \big\}.\]

We define the (left) shift $\sigma: \{0,1\}^{\N} \to \{0,1\}^{\N}$ at the function $\omega$ by the rule $(\sigma \omega)(i) = \omega(i+1)$, $i \in \N$.  The shift is easily seen to be a continuous self-map of $\{0,1\}^{\N}$.  The \emph{full shift} is the system $(\{0,1\}^{\N}, \sigma)$.

It will frequently be useful to consider functions on $\N_0$ instead of $\N$.  Everything written above applies equally well to the space $\{0,1\}^{\N_0}$ and the \emph{full shift} $(\{0,1\}^{\N_0}, \sigma)$.


\subsubsection{Dynamics on the space of ultrafilters}
\label{sec_topology_and_ultrafilters}

The Stone-\Cech{} compactification of $\N$, denoted $\beta \N$, is a compact Hausdorff, right-topological semigroup that provides a convenient universal object in the category of minimal systems.  We summarize here just what we will need; the interested reader is directed to \cite[Ch. 19]{hindman_strauss_book_2012} and the references therein for more information.

As a set, we realize $\beta \N$ as the set of all ultrafilters on $\N$ (cf. \cref{sec_filter_pr_and_ufs}). Sets of the form
\[\overline{A} \defeq \big\{p \in \beta \N \ \big| \ A \in p \big\}, \quad A \subseteq \N,\]
form a base for a non-metrizable, compact Hausdorff topology on $\beta \N$.  Since ultrafilters are families, addition of two ultrafilters is as defined in \cref{sec_family_algebra}.  It follows by combining \cref{lemma_equiv_ultrafilter} \eqref{item_family_is_self_dual} and \cref{lemma_filter_algebra} \eqref{item_family_prop_zero}, \eqref{item_family_prop_two}, and \eqref{item_sum_dual} that the sum of two ultrafilters is an ultrafilter.  Thus, $(\beta \N, +)$ is a compact Hausdorff semigroup.

It is not hard to check that for all $n \in \N$ and $q \in \beta \N$, the maps $\beta \N \to \beta \N$ defined by $p \mapsto n + p$ and $p \mapsto p + q$ are continuous.  Thus, $(\beta \N, +)$ is called a \emph{right-topological} semigroup.  A nonempty subset $L \subseteq \beta \N$ with the property that $p + L \subseteq L$ for all $p \in \beta \N$ is called a \emph{left ideal}.  A left ideal that is minimal amongst all left ideals is called a \emph{minimal left ideal}.  By a routine application of Zorn's lemma, minimal left ideals exist in $(\beta \N, +)$.  If $L \subseteq \beta \N$ is a minimal left ideal, for all $p \in L$, the set $\beta \N + p$ is a left ideal contained in $L$, and hence is equal to $L$.  Since addition on the right by $p$ is continuous and $\beta \N$ is compact, we see that minimal left ideals are compact.

Denoting addition by 1 on the left by $\lambda_1: \beta \N \to \beta \N$, we see that $(\beta \N, \lambda_1)$ is a non-metrizable topological dynamical system.  Fix $q \in \beta \N$.  Since addition on the right by $q$ is continuous and $\N$ is dense in $\beta \N$, we have that the orbit closure of $q$ under $\lambda_1$ is
\[\overline{\lambda_1^{\N} q} = \overline{ \N + q } = \overline{\N} + q = \beta \N + q,\]
a left ideal.  It is not hard to see that $\beta \N + q$ is a minimal left ideal if and only if the system $(\beta \N + q, \lambda_1)$ is minimal. (Indeed, the system $(\beta \N + q, \lambda_1)$ is minimal if and only if for all $p \in \beta \N + q$, the orbit closure $\beta \N + p$ -- a left ideal -- is equal to $\beta \N + q$.)  In this case, the ultrafilter $q$ is called a \emph{minimal ultrafilter}.  We see, then, that the minimal left ideals of the semigroup $(\beta \N, +)$ are precisely the minimal subsystems of $(\beta \N, \lambda_1)$.


Let $(X,T)$ be a not-necessarily-metrizable system, and let $x \in X$.  Consider the map $T^{\displaystyle \cdot} x: \N \to X$ given by $n \mapsto T^n x$.  By the universal property of the Stone-\Cech{} compactification, the map $T^{\displaystyle \cdot} x$ lifts to a continuous map $T^{\displaystyle \cdot} x: \beta \N \to X$.  Concretely, for $p \in \beta \N$, the point $T^p x \in X$ is defined uniquely by the property that for all open $U \subseteq X$ containing $T^p x$, the set $R(x,U)$ is a member of $p$.  The following fact will be useful for us later on: if $U \subseteq X$ is clopen and $p \in \beta \N$, then
\begin{align}
    \label{eqn_ultrafilter_shift_of_return_set}
    R(x,U) - p = R(T^p x, U).
\end{align}
Indeed, for $n \in \N$, we see that $n \in R(x,U) - p$ if and only if $R(T^n x, U) = R(x,U) - n \in p$.  If $R(T^n x, U) \in p$, then
\[T^n T^p x = T^{n + p} x = T^{p + n} x = T^p T^n x \in \overline{U} = U,\]
whereby $n \in R(T^p x, U)$.  On the other hand, if $n \in R(T^p x, U)$, then $T^p T^n x = T^n T^p x \in U$, whereby $R(T^n x, U) \in p$.

Minimal left ideals of $\beta \N$ under addition by 1 are universal minimal dynamical systems.  More precisely, let $L \subseteq \beta \N$ be a minimal left ideal, and let $(X,T)$ be a not-necessarily-metrizable, minimal dynamical system.  Fix $x \in X$.  The map $T^{\displaystyle \cdot} x: \beta \N \to X$ described in the previous paragraph is continuous and satisfies $T^{\lambda_1 p} x = T^{1 + p} x = T (T^p x)$ and hence is a factor map $(\beta \N,\lambda_1) \to (X,T)$ of systems.  Restricted to $L$, by the minimality of $(X,T)$, we see that $T^{\displaystyle \cdot} x: (L, \lambda_1) \to (X,T)$ is surjective and hence is a factor map of minimal systems.


\subsection{Families of sets from dynamics}
\label{sec_families_of_sets}

We outline below those families of subsets of $\N$ that are important in this work.  The following nomenclature will be particularly convenient.  Let $\familytwo$ be a family. A set $A \subseteq \N$ is called a \emph{$\familytwo$ set} if it is a member of $\familytwo$, and a family $\familyone$ is called a \emph{$\familytwo$ family} if all members of $\familyone$ are $\familytwo$ sets, ie., $\familyone \subseteq \familytwo$. We will frequently consider in this work, for example, syndetic filters, ie., filters whose members are all syndetic.

\subsubsection{Syndetic, thick, and piecewise syndetic sets}
\label{sec_syndetic_thick_ps_sets}

A set $A \subseteq \N$ is \dots
\begin{enumerate}
    \item \dots \emph{syndetic} if there exists $N \in \N$ such that
    \[
        A \cup (A-1) \cup \cdots \cup (A-N) = \N;
    \]

    \item \dots \emph{thick} if for all finite $F \subseteq \N$, there exists $n \in \N$ such that $F + n \subseteq A$;

    \item \dots \emph{piecewise syndetic} if there exists $N \in \N$ such that the set
    \[A \cup (A-1) \cup \cdots \cup (A-N)\]
    is thick.
\end{enumerate}
We denote by $\syndetic$, $\thick$, and $\PS$ the families of syndetic, thick, and piecewise syndetic subsets of $\N$, respectively. All three families are translation invariant. It is well-known and not difficult to show that the families of syndetic and thick sets are dual and that
\[\PS = \syndetic \classcap \thick.\]
By \cref{lemma_class_cap_is_pr}, the family $\PS$ is piecewise syndetic and its dual, $\PS^*$, is a filter. It is not hard to show that a set $A$ belongs to $\PS^*$ if and only if for all finite $F \subseteq \N$, there exists a syndetic set $S \subseteq \N$ such that $F + S \subseteq A$.  Thus, members of the filter $\PS^*$ are called \emph{thickly syndetic} sets.

In the following lemmas and in those in the next section, we demonstrate how well-known families of subsets of $\N$ can be characterized by membership in certain types of filters.  These results will be useful for us later on.

\begin{lemma} \leavevmode
\label{lemma_thick_set_filter_char}
    \begin{enumerate}
        \item
        \label{item_proper_filter_iff_thick}
        Every proper, translation-invariant filter is a thick, translation-invariant filter.

        \item
        \label{item_thick_iff_proper_filter}
        A set $A \subseteq \N$ is thick if and only if it is a member of a proper, translation-invariant filter.
    \end{enumerate}
\end{lemma}

\begin{proof}
    \eqref{item_proper_filter_iff_thick} \ Let $\family$ be a proper, translation-invariant filter, and let $A \in \family$.  For all $N \in \N$, the set $\bigcap_{n=0}^N (A-n) \in \family$, and hence is nonempty.  It follows that $A$ contains an interval of length $N+1$, and hence that $A$ is thick.  It follows that $\family$ is thick.

    
    \eqref{item_thick_iff_proper_filter} \ If $A \subseteq \N$ is a member of a proper, translation-invariant filter, then it is thick, as shown in the proof of \eqref{item_proper_filter_iff_thick}.  Conversely, if $A \subseteq \N$ is thick, then the family
    \[\bigupclose \left \{ \bigcap_{n=0}^N (A-n) \ \middle| \ N \in \N \right\}\]
    is quickly checked to be a thick (hence, proper), translation-invariant filter.
\end{proof}

The following lemma gives a useful characterization of the filter of thickly syndetic sets as the largest syndetic, translation-invariant filter.

\begin{lemma}
\label{lemma_condition_on_subfamily_of_ps_star}
    The family $\PS^*$ is a syndetic, translation-invariant filter.  If $\family$ is a syndetic, translation-invariant filter, then $\family \subseteq \PS^*$.
\end{lemma}
\begin{proof}
    That $\PS^*$ is a syndetic, translation-invariant filter was discussed above.
    
    Let $\family$ be a syndetic, translation-invariant filter, and let $A \in \family$.  Since $\family$ is translation invariant, for all $N \in \N$, the sets $A-1$, \dots, $A-N$ all belong to $\family$. Since $\family$ is a syndetic filter, the set $(A - 1) \cap \cdots \cap (A - N)$ belongs to $\family$ and is syndetic.  Since
    \[\{1, \ldots, N\} + \big( (A - 1) \cap \cdots \cap (A - N) \big) \subseteq A,\]
    we see that $A$ is thickly syndetic and hence is a member of $\PS^*$.
\end{proof}

\subsubsection{Finite sums sets and central sets}
\label{sec_ip_and_central_sets}

The families of infinite, finite sums sets and central sets will also appear in this work.  A set $A \subseteq \N$ is called an \emph{$\IP$ set} (also, an \emph{infinite, finite sums set}) if there exists a sequence $(x_i)_{i=1}^\infty \subseteq \N$ such that
\[\FS(x_i)_{i=1}^\infty \defeq \left\{ \sum_{f \in F} x_f \ \middle | \ F \subseteq \N \text{ is finite and nonempty} \right\} \subseteq A.\]
Here ``FS'' stands for ``finite sums.''  It is a consequence of Hindman's theorem \cite{hindman_1974} that the family $\IP$ is partition regular (see \cite[Lemma 2.1]{bergelson_hindman_1993}). Thus, the dual family $\IP^*$ is a filter.

It is well known \cite[Thm. 5.12]{hindman_strauss_book_2012} that a set is an $\IP$ set if and only if it is member of an idempotent ultrafilter.  The following lemma offers a characterization in terms of idempotent filters, showing that the assumption of maximality (recall, ultrafilters are maximal filters) is superfluous.

\begin{lemma}
\label{lemma_ip_set_iff_contained_in_prop_idemp_filter}
    A set $A \subseteq \N$ is an $\IP$ set if and only if it is a member of a proper, idempotent filter.
\end{lemma}

\begin{proof}
    Suppose $A \subseteq \N$ is an $\IP$ set.  There exists $(x_i)_{i=1}^\infty \subseteq \N$ such that $\FS(x_i)_{i=1}^\infty \subseteq A$.  We will check that the family
    \[\family \defeq \upclose \big \{ \ \FS(x_i)_{i=N}^\infty \ \big | \ N \in \N \big\}\]
    is a proper, idempotent filter containing $A$. It is easy to see from the definition that $\family$ is a proper filter.
    
    To see that $\family$ is idempotent, let $A \in \family$.  We must show that $A - \family \in \family$.  By definition, there exists $N \in \N$ such that $\FS(x_i)_{i=N}^\infty \subseteq A$.  We claim that $\FS(x_i)_{i=N}^\infty \subseteq A - \family$, from which it will follow that $A - \family \in \family$.  Let $x \in \FS(x_i)_{i=N}^\infty$.  We must show that $A - x \in \family$.  There exists $F \subseteq \N$ finite with $\min F \geq N$ such that $x = \sum_{f \in F} x_f$.  We see that $A - x \supseteq \FS(x_i)_{i=\max F  + 1}^\infty$, which belongs to $\family$.  Therefore, $A - x \in \family$, as desired.

    Conversely, suppose that $A$ is a member of a proper, idempotent filter $\family$.  Since $A \in \family$ and $A - \family \in \family$, there exists $x_1 \in A \cap (A - \family)$.  Since $A \in \family$ and $A - x_1 \in \family$, we have that $A \cap (A-x_1) \in \family$, and hence that $\big(A \cap (A-x_1) \big) - \family \in \family$.  There exists
    \[x_2 \in A \cap (A-x_1) \cap \big( \big(A \cap (A-x_1) \big) - \family \big).\]
    So far, we have found $x_1, x_2 \in \N$ such that $\{x_1, x_2, x_1 + x_2\} \subseteq A$.  Iterating, we will take $x_3$ from the set
    \[A \cap (A-x_1) \cap \big( \big(A \cap (A-x_1) \big) - x_2 \big) \cap \bigg( \Big( A \cap (A-x_1) \cap \big( \big(A \cap (A-x_1) \big) - x_2 \big) \Big) - \family \bigg)\]
    so that all finite sums from the set $\{x_1, x_2, x_3\}$ belong to $A$. It is clear how to proceed to find $(x_i)_{i=1}^\infty \subseteq \N$ such that $\FS(x_i)_{i=1}^\infty \subseteq A$, as desired.
\end{proof}

We will consider the family of central sets for its importance to the subject and for its connection to the families of dynamically defined sets at the heart of this work.  Furstenberg \cite[Def. 8.3]{furstenberg_book_1981} defined a set $A \subseteq \N$ to be \emph{central} if there exists a system $(X,T)$, a uniformly recurrent point $x \in X$, a point $y \in X$ proximal to $x$, and an open set $U \subseteq X$ containing $x$ such that $A = R(y,U)$.  We denote the family of central sets by $\central$.

Bergelson, Hindman, and Weiss \cite{bergelson_hindman_1990} showed that a subset of $\N$ is central if and only if it is a member of a minimal idempotent ultrafilter on $\N$.  (This characterization of central sets fits in a series of related characterizations in more general settings: by Glasner \cite{glasner_1980} in countable, abelian groups; Bergelson-Hindman-Weiss \cite{bergelson_hindman_1990} in countable semigroups; and Shi-Yang \cite{shi_yang_1996} in all semigroups.) In analogy to $\IP$ sets and \cref{lemma_ip_set_iff_contained_in_prop_idemp_filter}, central sets can be characterized in terms of a special class of idempotent filters.  Defining this class -- ``collectionwise piecewise syndetic,'' idempotent filters -- would require too large a diversion, so we leave it to the interested reader. (See \cite[Def. 14.19]{hindman_strauss_book_2012} for a definition of collectionwise piecewise syndetic.)

It was stated by Furstenberg and follows immediately from the ultrafilter characterization of central sets that the family $\central$ is partition regular. It is well known that thick sets are central and that central sets are piecewise syndetic and $\IP$.  Dual to the family of central sets is the family $\central^*$, a syndetic filter that will appear prominently in \cref{sec_dps_sets}. \cref{fig:containmentdiagram} depicts the relationships between the families of subsets of $\N$ introduced so far.



\begin{figure}[ht]
\centering
\begin{tikzpicture}[>=triangle 60]
  \matrix[matrix of math nodes,column sep={30pt,between origins},row sep={20pt,between origins}](m)
  {
    & & & & & & & & |[name=psstar]|\PS^* = \dPS^* & & & & & & \\
    \\
    & & & & & & |[name=thick]|\thick & &  \\
    & & & & & & & & |[name=dcpsstar]|\dcPS^* \\
    & & & & |[name=dthick]|\dthick & |[name=ipstar]|\IP^* & & & & & \\
    & & & & & & |[name=centralstar]|\central^* & & & & |[name=dcsyndetic]| \dcsyndetic \\
    \\
    & & & & |[name=dcthick]|\dcthick & & & & |[name=central]|\central \\
    & & & & & & & & & |[name=ip]|\IP & |[name=dsyndetic]| \dsyndetic \\
    & & & & & & |[name=dcps]| \dcPS & & & & \\
    & & & & & & & & |[name=syndetic]|\syndetic \\
    \\
    & & & & & & |[name=ps]|\PS = \dPS \\
 };
 
   \draw[-angle 90]
            (psstar) edge (thick)
            (psstar) edge (dcpsstar)
            (thick) edge (dthick)
            (dthick) edge (dcthick)
            (ipstar) edge (centralstar)
            (dcpsstar) edge (centralstar)
            (dcpsstar) edge (dcsyndetic)
            (centralstar) edge (dcthick)
            (thick) edge (central)
            (centralstar) edge (central)
            (dcsyndetic) edge (central)
            (dcsyndetic) edge (dsyndetic)
            (dcthick) edge (dcps)
            (central) edge (dcps)
            (central) edge (ip)
            (centralstar) edge (syndetic)
            (dsyndetic) edge (syndetic)
            (syndetic) edge (ps)
            (dcps) edge (ps)
  ;
\end{tikzpicture}
\caption{Containment amongst the families appearing in this paper.  An arrow $\family \to \familytwo$ indicates that $\family \subsetneq \familytwo$. Those families not prefixed with the letter ``$d$'' are defined in \cref{sec_families_of_sets}.  The families $\dcsyndetic$ / $\dsyndetic$, $\dcthick$ / $\dthick$, and $\dcPS$ / $\dPS$ are the main objects of study in this work and are defined in Sections \ref{sec_dsyndetic}, \ref{sec_dthick}, and \ref{sec_dps_sets}, respectively.  There are eight families and their duals, yielding sixteen families in all, two pairs of which coincide. The diagram's symmetry is explained by the family duality described in \cref{sec_dual_families}.}
\label{fig:containmentdiagram}
\end{figure}


The families $\IP$ and $\central$ are examples of idempotent families that are not translation invariant.  It is a short exercise to verify this directly from the definitions.  Alternatively, one can note two things: first, by the characterizations given above,
\[\IP = \bigcup_{\substack{p \in \beta \N \\ p \text{ idempotent}}} p \qquad \text{ and } \qquad \central = \bigcup_{\substack{p \in \beta \N \text{ minimal} \\ p \text{ idempotent}}} p,\]
and second, a union of idempotent families is idempotent.



\section{Part I: Dynamically syndetic sets}
\label{sec_dsyndetic}

In this section, we define the family of dynamically syndetic sets and give a number of equivalent characterizations of dynamical syndeticity.  The main result is \cref{thm_dcs_iff_cs}, a combinatorial characterization of dynamical syndeticity in terms of syndetic, idempotent filters.

\subsection{Definitions and first characterizations}
\label{sec_dy_synd_def_and_char}

A set $A \subseteq \N$ is \emph{dynamically syndetic} if there exists a minimal topological dynamical system $(X,T)$, a point $x \in X$, and a nonempty, open set $U \subseteq X$ such that $R(x,U) \subseteq A$.  If $x \in U$, the set $A$ is called \emph{dynamically central syndetic}.  The families of dynamically syndetic and dynamically central syndetic subsets of $\N$ are denoted $\dsyndetic$ and $\dcsyndetic$, respectively.

It is clear from the definitions and the facts in \cref{sec_top_dynamics} that
\[\dcsyndetic \subseteq \dsyndetic \subseteq \syndetic.\]
The set of odd, positive integers and the set in \eqref{eqn_intro_even_odd} from the introduction demonstrate that each of these inclusions is proper. (That $2\N-1 \not\in \dcsyndetic$ follows from, eg., \cref{lemma_dilates_of_dsyndetic_sets}.)  Here is a short list of examples of dynamically (central) syndetic sets that may be helpful to have in mind.  Let $\alpha \in \R \setminus \Q$.
\begin{enumerate}
    \item \label{item:irrational_rotation_dcS_definition}
    The irrational rotation $(\R / \Z, x \mapsto x + \alpha)$, where $\alpha$ is interpreted as an element of $\R / \Z$, is a minimal system.  For all $x \in \R / \Z$ and all nonempty, open $U \subseteq \R / \Z$, the set
    \[R(x,U) = \big\{ n \in \N \ \big| \  x + n \alpha  \in U \big\}\]
    is dynamically syndetic and, if $x \in U$, dynamically central syndetic.  More generally, supersets of Bohr (Bohr$_0$) sets -- those arising in this way from compact group rotations -- are dynamically (central) syndetic.

    \item When $\alpha > 1$, the sequence $(\lfloor n \alpha \rfloor)_{n \in \N}$ is called a \emph{Beatty sequence}.  It is not hard to show that the image of this sequence is equal to the set
    \[\big\{ n \in \N \ \big| \ \big\{ n / \alpha \} \in \big( (\alpha-1) / \alpha, 1 \big) \big\},\]
    where $\{ \cdot \}: \R \to \R / \Z$ is the quotient map.  Thus, Beatty sequences yield Bohr, and hence dynamically syndetic, sets.  What is not so clear is that, in fact, Beatty sequences yield dynamically central syndetic sets.  We will see this as a consequence of \cref{thm:simple_equivalent_dcS}.

    \item Using the notation from \eqref{item:irrational_rotation_dcS_definition}, the set
    \[\big\{ n \in \N \ \big| \  n^2 \alpha  \in U \big\}\]
    is a dynamically syndetic set.  It is, in fact, an example of a nil-Bohr set, one that arises from a rotation on a compact homogeneous space of a nilpotent Lie group.  Nil-Bohr (nil-Bohr$_0$) sets are dynamically (central) syndetic \cite{Host-Kra_nilbohr}.
    
    \item \label{item_dcs_two_adic_example} Denote by $\nu_2: \N \to \N_0$ the 2-adic valuation.  The set
    \[A \defeq \big\{ n \in \N \ \big | \ \nu_2(n) \text{ is even} \big \}\]
    is a dynamically syndetic set.  To see this, one can show that the point $1_A$ is uniformly recurrent in the full shift $(\{0,1\}^{\N_0},\sigma)$, whereby $A = R_\sigma(1_A,[1]_0)$ is dynamically syndetic.  In fact, the point $1_{A \cup \{0\}} \in [1]_0$ is uniformly recurrent, and so $A$ is dynamically central syndetic.  By \cref{lemma_dilates_of_dsyndetic_sets} and the fact that $\N \setminus A = A / 2$, the same can be said about the set of positive integers with odd 2-adic valuation.
    
    \item The \emph{Chacon word} $0010001010010 0010001010010 1 0010001010010\ldots$ is the unique, non-constant word in the symbols $0$ and $1$ that is fixed under the substitutions $0 \mapsto 0010$ and $1 \mapsto 1$.  The set of locations of the 1's in the Chacon word is a dynamically syndetic set, for the same reason as in the previous point: the word, interpreted as a point in the full shift $(\{0,1\}^{\N_0},\sigma)$, is known to be uniformly recurrent \cite[Sec. 5.6.6]{devries_2014}.\\
\end{enumerate}

Many of the results in this section demonstrate the robustness of the family of dynamically (central) syndetic sets under changes to the combinatorial, topological, or dynamical requirements in its definition.  The family of dynamically (central) syndetic subsets of $\N$ remains unchanged if \dots
\begin{enumerate}
    \item \dots \ the map $T$ is required to be a homeomorphism, that is, if the system $(X,T)$ is required to be invertible.  This is shown in \cref{lemma_may_assume_invertible}.

    \item \dots \ the space $X$ is allowed to be non-metrizable. This is a consequence of \cref{lemma_metrizability_doesnt_matter}.

    \item \dots \ the space $X$ is required to be zero dimensional and the set $U$ is required to be clopen. This is a consequence of \cref{thm:simple_equivalent_dcS}.

    \item \dots \ (for dynamically central syndetic sets) the point $x$ is allowed to be in $\overline{U}$ instead of being required to be in $U$.  Sets of this form are called \emph{very strongly central} in \cite[Def. 2.10]{bergelson_hindman_strauss_2012}, where a number of equivalent characterizations in terms of ultrafilters are given.  That the families of very strongly central sets and dynamically central syndetic sets are the same is shown in \cref{thm:simple_equivalent_dcS}.

    \item \dots \ sets of return times $R(x,U)$ are considered up to equivalence on thickly syndetic sets ($\PS^*$ sets). This is a consequence of \cref{cor:largest_family_satisfies_cap}: a set $A \subseteq \N$ is dynamically (central) syndetic if and only if for all thickly syndetic $B \subseteq \N$, the set $A \cap B$ is dynamically (central) syndetic.  We will have immediate need for a more basic version of this result for cofinite sets in \cref{lemma_ds_dcs_modifyable_on_a_finite_set}.\\
\end{enumerate}

There are other directions, however, in which this definitional robustness does not extend.
\begin{enumerate}
\setcounter{enumi}{5}
    \item Minimality is a key feature of the definition.  For all $A \subseteq \N$, there exists a system $(X,T)$, a point $x \in X$, and a nonempty, open set $U \subseteq X$ containing $x$ such that $A = R(x,U)$.  Indeed, note that considering $A$ as a subset of $\N_0$, we have that $A = R(1_{A},[1]_0)$ in the full shift $(\{0,1\}^{\N_0},\sigma)$.  Thus, dropping the minimality assumption trivializes the definition of dynamical syndeticity.

    \item \label{item_generalization_to_z} That we are considering actions of the semigroup $(\N,+)$ instead of the group $(\Z,+)$ matters.  For example, if we define dynamically central syndetic subsets of $\Z$ analogously to those in $\N$, we see that all dynamically central syndetic subsets of $\Z$ contain 0.  In this case, if $x \in \overline{U} \setminus U$, then the set $R(x,U) \subseteq \Z$ is, in contrast to \cref{thm:simple_equivalent_dcS}, not dynamically central syndetic because it does not contain 0.  Other results from this paper for $(\N,+)$-systems fail to generalize under this definition, too, most notably the fact that thickly syndetic ($\PS^*$) sets (which do not necessarily contain 0) are dynamically central syndetic. Along with other evidence, this suggests that allowing $x \in \overline{U}$ in the definition of the family of dynamically central syndetic sets (as in \cref{thm:simple_equivalent_dcS} \eqref{item:very_strongly_central}) may yield the ``right'' notion for more general (semi)group actions.
\end{enumerate}



\subsubsection{Definitional robustness of dynamically syndetic sets}
\label{sec_dsyndetic_robust}

Here we show that invertibility, changes on finite sets, and metrizability do not affect the definition of the family of dynamically syndetic sets.

\begin{lemma}
\label{lemma_may_assume_invertible}
    A set $A \subseteq \N$ is dynamically (central) syndetic if and only if there exists an invertible system $(X,T)$, a point $x \in X$, and a nonempty, open set $U \subseteq X$ (containing $x$) such that $R(x,U) \subseteq A$.
\end{lemma}

\begin{proof}
    The ``if'' statement is immediate.  To see the ``only if'' statement, suppose $A \subseteq \N$ is dynamically syndetic.  There exists a system $(X,T)$, a point $x \in X$, and a nonempty, open set $U \subseteq X$ such that $R(x,U) \subseteq A$.  Let $\pi: (W,T) \to (X,T)$ be the natural extension (see \cite[Definition 6.8.10]{downarowiczbook}) of $(X,T)$, and let $w \in W$ be such that $\pi w = x$. Since $(X,T)$ is minimal, so is $(W,T)$ (see, for example, \cite[Lemma 2.9]{glasscock_koutsogiannis_richter_2019}). The system $(W,T)$ is invertible, and $R(w,\pi^{-1} U) = R(x,U)$.  Thus, the set $A$ contains $R(w,\pi^{-1} U)$ from the invertible system $(W,T)$.  Note that if $A$ is dynamically central syndetic and $x \in U$, then $w \in \pi^{-1} U$, whereby the set $R(x,U) = R(w,\pi^{-1} U)$ is dynamically central syndetic from the invertible system $(W,T)$.
\end{proof}

The following lemma shows that dynamically (central) syndetic sets remain such after changes on a finite set.  This result is upgraded in \cref{cor:largest_family_satisfies_cap}, which reaches the same conclusion after changes are allowed on a non-piecewise syndetic set.

\begin{lemma}
    \label{lemma_ds_dcs_modifyable_on_a_finite_set}
    Let $A \subseteq \N$ be dynamically (central) syndetic.  If $B \subseteq \N$ is cofinite, then $A \cap B$ is dynamically (central) syndetic.
\end{lemma}

\begin{proof}
    By the definition of dynamically (central) syndetic, it suffices to show that for all minimal systems $(X,T)$, all points $x \in X$, all nonempty, open sets $U \subseteq X$ (containing $x$), and all $N \in \N$, the set $R(x,U) \setminus \{1, \ldots, N\}$ is dynamically (central) syndetic.  So, let $(X,T)$, $x \in X$, $U \subseteq X$ (contain $x$), and $N \in \N$ be as described.

    The set $V \defeq U \setminus \{Tx, T^2x, \ldots, T^N x\}$ is open and $R(x,U) \setminus \{1, \ldots, N\} \supseteq R(x, V)$.  If $V$ is nonempty, then we see that the set $R(x,U) \setminus \{1, \ldots, N\}$ is dynamically syndetic, as desired. Note that if $x \in \{Tx, T^2x, \ldots, T^N x\}$, then the system is periodic, $X = \{Tx, T^2x, \ldots, T^N x\}$, and $V$ is empty.  Therefore, if $V$ is nonempty and $x \in U$, then $x \in V$, and the set $R(x,U) \setminus \{1, \ldots, N\}$ is dynamically central syndetic, as desired.

    If, on the other hand, the set $V$ is empty, then $U \subseteq \{Tx, T^2x, \ldots, T^N x\}$.  Since $U$ is open, we see that some $T^i x$ is an isolated point of $X$.  By \cref{lemma_isolated_points_and_minimal_systems}, the system $(X,T)$ is finite and periodic.  In this case, the set $R(x,U) \setminus \{1, \ldots, N\}$ contains an infinite arithmetic progression (of the form $k \N$), which is dynamically (central) syndetic, as desired.
\end{proof}


\begin{theorem}
\label{lemma_metrizability_doesnt_matter}
Let $(X,T)$ be a not-necessarily-metrizable, minimal topological dynamical system.  For all $x \in X$ and all nonempty, open $U \subseteq X$ (containing $x$), the set $R(x,U)$ is dynamically (central) syndetic.
\end{theorem}

\begin{proof}
    Let $x \in X$ and $U \subseteq X$ be nonempty and open.  Fix $L \subseteq \beta \N$ a minimal left ideal.  As described in \cref{sec_topology_and_ultrafilters}, the map $\pi: (L, \lambda_1) \to (X,T)$ given by $\pi(p) = T^p x$ is a factor map of minimal systems.  By \cite[Thm. 19.23]{hindman_strauss_book_2012}, since $x$ is uniformly recurrent, there exists an idempotent $p \in L$ such that $\pi(p) = x$. It is quick to check that
    \[R_T(x, U) = R_{\lambda_1}(p, \pi^{-1} U).\]
    We will show that the set $R_{\lambda_1}(p, \pi^{-1} U)$ is dynamically syndetic.

    By the definition of the topology on $\beta \N$, since $\pi^{-1} U$ is nonempty and open, there exists $A \subseteq \N$ such that $\emptyset \neq \overline{A} \cap L \subseteq \pi^{-1} U$.  In the notation defined in \cref{sec_family_algebra}, we see that
    \begin{align}
    \label{eqn_requisite_containment_of_non_metrizable_returns}
        R_{\lambda_1}(p, \pi^{-1} U) \supseteq R_{\lambda_1}(p, \overline{A}) = A - p.
    \end{align}
    Since $\overline{A} \cap L \neq \emptyset$ and $(L,\lambda_1)$ is minimal, we see that the set $A-p$ is nonempty (indeed, syndetic).

    Considering $1_A$ as an element of $\{0,1\}^{\N_0}$, we see in the full shift $(\{0,1\}^{\N_0}, \sigma)$ that $A = R_{\sigma}(1_A,[1]_0)$.  It follows by \eqref{eqn_ultrafilter_shift_of_return_set} in \cref{sec_topology_and_ultrafilters} that
    \[A - p = R_{\sigma}(1_A,[1]_0) - p = R_{\sigma}(T^p 1_A, [1]_0).\]
    Since $p$ is a minimal idempotent ultrafilter, we have by \cite[Thm. 19.23]{hindman_strauss_book_2012} that the point $T^p 1_A$ is uniformly recurrent under the shift.  Since $R_{\sigma}(T^p 1_A, [1]_0)$ is nonempty, we see that the set $R_{\sigma}(T^p 1_A, [1]_0)$, and hence the set $A - p = R_{\lambda_1}(p, \overline{A})$, is dynamically syndetic.  Because these families are upward closed, by \eqref{eqn_requisite_containment_of_non_metrizable_returns}, the same can be said about the set $R_{\lambda_1}(p, \pi^{-1} U) = R_T(x, U)$, as was to be shown.

    In the case that $x \in U$, we wish to show that the set $R_T(x,U)$ is dynamically central syndetic.  The argument given above can be followed with some additional information.  First, since $x \in U$, the point $p \in \pi^{-1} U$.  The set $A \subseteq \N$ can then be assumed to be such that $p \in \overline{A} \cap L$. Considering $1_A$ as an element of $\{0,1\}^{\N_0}$, we see
    \[R_\sigma(1_A,[1]_0) = A \in p,\]
    whereby $T^p 1_A \in [1]_0$.  Thus, the set $A-p = R_{\sigma}(T^p 1_A, [1]_0)$ is dynamically central syndetic, as desired.
\end{proof}

\subsubsection{Translates and dilates of dynamically syndetic sets}

In this section, we describe what happens to dynamically (central) syndetic sets under translation and dilation.  Recall that $A + n$ and $A-n$ are computed as subsets of $\N$, as described in \cref{sec_prelims}.


\begin{lemma}
\label{lemma_translates_of_dsyndetic_sets}
    Let $A \subseteq \N$ be dynamically syndetic.
    \begin{enumerate}
        \item
        \label{item_all_translates_are_ds}
        For all $n \in \N$, the set $A-n$ is dynamically syndetic.
        
        \item
        \label{item_some_translates_are_dcs}
        The set
        \begin{align}
            \label{eqn_translates_of_ds_set}
            \big\{ n \in \N \ \big| \ A-n \text{ is dynamically central syndetic} \big\}
        \end{align}
        is dynamically syndetic.
        
        \item
        \label{item_dcs_translates_of_dcs_are_dcs}
        If $A$ is dynamically central syndetic, then the set in \eqref{eqn_translates_of_ds_set} is dynamically central syndetic.
    \end{enumerate}
    Moreover, \eqref{item_all_translates_are_ds}, \eqref{item_some_translates_are_dcs}, and \eqref{item_dcs_translates_of_dcs_are_dcs} hold with $A-n$ replaced by $A+n$. Finally,
    \begin{align}
        \label{eqn_relationship_between_ds_and_dcs_classes}
        \dsyndetic = \bigcup_{n \in \N} \big(\dcsyndetic - n \big) = \bigcup_{n \in \N} \big(\dcsyndetic + n \big).
    \end{align}
\end{lemma}


\begin{proof}
    Since $A$ is dynamically syndetic, there exists a minimal system $(X,T)$, a point $x \in X$, and a nonempty, open set $U \subseteq X$ such that $R(x,U) \subseteq A$. According to \cref{lemma_may_assume_invertible}, we can assume the system $(X, T)$ is invertible.
    
    For $n \in \N$, we have that
    \begin{align}
        \label{eqn_translate_of_rxu_in_aminusn}
        R(T^n x, U) = R(x,U) - n \subseteq A - n.
    \end{align}
    Thus, the set $A-n$ is dynamically syndetic, demonstrating \eqref{item_all_translates_are_ds}.  It also follows from \eqref{eqn_translate_of_rxu_in_aminusn} that $R(x,U)$ is a subset of the set in \eqref{eqn_translates_of_ds_set}.  Indeed, if $n \in R(x,U)$, then $T^n x \in U$, so the set $A-n$ is dynamically central syndetic.  The set $R(x,U)$ is dynamically syndetic, whereby the set in \eqref{eqn_translates_of_ds_set} is dynamically syndetic, demonstrating \eqref{item_some_translates_are_dcs}.  All of this holds in the case that $A$ is dynamically central syndetic, but we have additionally that $R(x,U)$ is dynamically central syndetic, demonstrating \eqref{item_dcs_translates_of_dcs_are_dcs}.

    Now we will show that \eqref{item_all_translates_are_ds}, \eqref{item_some_translates_are_dcs}, and \eqref{item_dcs_translates_of_dcs_are_dcs} hold with $A-n$ replaced by $A+n$.
    
    Let $n \in \N$.  Recall that $T: X \to X$ is a homeomorphism. We claim that $R(x,U) + n = R(T^{-n}x, U) \cap \{n+1,n+2,\ldots\}$.  Indeed, note that $m \in R(x,U) + n$ if and only if $m - n \in R(x,U)$ if and only if $T^{m-n} x \in U$ and $m \geq n+1$ if and only if $T^m T^{-n}x \in U$ and $m \geq n+1$ if and only if $m \in R(T^{-n}x,U) \cap \{n+1, n+2, \ldots\}$.  By \cref{lemma_ds_dcs_modifyable_on_a_finite_set}, the set $R(T^{-n}x,U) \cap \{n+1,n+2,\ldots\}$ is dynamically syndetic.  Therefore, the set $A + n$ is dynamically syndetic, demonstrating \eqref{item_all_translates_are_ds}.
    
    Since $(X,T)$ is minimal and invertible, the system $(X,T^{-1})$ is minimal (see, for example, \cite[Lemma 2.7]{glasscock_koutsogiannis_richter_2019}). We claim that the set $R_{T^{-1}}(x,U)$ is contained in the set in \eqref{eqn_translates_of_ds_set}.  This will finish the proof of \eqref{item_some_translates_are_dcs} and \eqref{item_dcs_translates_of_dcs_are_dcs} since $R_{T^{-1}}(x,U)$ is dynamically syndetic and dynamically central syndetic when $x \in U$.  Suppose $n \in R_{T^{-1}}(x,U)$ so that $T^{-n} x \in U$.  From the previous paragraph, we see that $R(T^{-n}x, U) \cap \{n+1, n+2, \ldots\} = R(x,U) + n \subseteq A +n$.  By \cref{lemma_ds_dcs_modifyable_on_a_finite_set}, we have that the set $A + n$ is dynamically central syndetic, whereby $n$ belongs to the set in \eqref{eqn_translates_of_ds_set}, as desired.

    Finally, that \eqref{eqn_relationship_between_ds_and_dcs_classes} holds follows immediately from \eqref{item_all_translates_are_ds} and \eqref{item_some_translates_are_dcs} for $A-n$ and $A+n$.
\end{proof}

It is interesting to formulate the conclusions of \cref{lemma_translates_of_dsyndetic_sets} in terms of the family algebra developed in \cref{sec_combinatorics}. Thus,
\begin{enumerate}
    \item the family $\dsyndetic$ is translation invariant, that is, $\dsyndetic \subseteq \{\N\} + \dsyndetic$;
    \item $\dsyndetic \subseteq \dsyndetic + \dcsyndetic$; and
    \item the family $\dcsyndetic$ is idempotent, that is, $\dcsyndetic \subseteq \dcsyndetic + \dcsyndetic$.\\
\end{enumerate}

The following lemma, well known to experts, describes the families of piecewise syndetic and central sets in dynamical terms.  Both equalities can be derived from the result of Auslander \cite{auslander_1960} and Ellis \cite{ellis_1957} that any point in any system is proximal to a uniformly recurrent point.  We opt for a short derivation of the first from the second, which is shown in \cite{Huang_Shao_Ye_2019}.


\begin{lemma}
\label{lem:central_characterization}
We have that $\PS = \dsyndetic \classcap \thick$ and $\central = \dcsyndetic \classcap \thick$.
\end{lemma}


\begin{proof}
    The second equality is proven in \cite[Thm. 3.7]{Huang_Shao_Ye_2019}.  We will derive the first from the second.  Suppose $A \subseteq \N$ is piecewise syndetic; it is not hard to see from the definition that there exists $N \in \N$ such that $A \cup (A+1) \cup \cdots \cup (A+N)$ is thick.  Since thick sets are central and the family of central sets is partition regular, some $A+i$ is central.  By the second equality, there exists $B \in \dcsyndetic$ and $H \in \thick$ such that $A+i \supseteq B \cap H$.  Therefore, we have that $A \supseteq (B-i) \cap (H-i)$.  By \cref{lemma_translates_of_dsyndetic_sets}, the set $B-i$ is dynamically syndetic and the set $H-i$ is thick, as was to be shown.
\end{proof}

For the following lemma, recall the dilation notation set out at the beginning of \cref{sec_prelims}.

\begin{lemma}
\label{lemma_dilates_of_dsyndetic_sets}
    Let $A \subseteq \N$ and $k \in \N$. If $A$ is dynamically central syndetic, then so are the sets $kA$ and $A/k$.  If $A$ is dynamically syndetic, then so is the set $kA$.
\end{lemma}

\begin{proof}
    Let $A \in \dcsyndetic$. Let $(X, T)$ be a minimal system, $x \in X$, and $U \subseteq X$ be an open neighborhood of $x$ such that $A \supseteq R(x, U)$. We will show that $k A \in \dcsyndetic$ by considering a suspension system with base $(X, T)$. More precisely, let $Y = X \times \{0, 1, \ldots, k - 1\}$, and define $S: Y \to Y$ by 
    \[
        S(x', \eta) = \begin{cases} (x', \eta + 1) \text{ if } \eta \leq k - 2, \\
        (Tx', 0) \text{ if } \eta = k - 1.
        \end{cases}
    \]
    It is easy to see that the orbit of every point in $Y$ is dense, whereby $(Y, S)$ is minimal.  Also, we see that $k R_T(x, U) = R_S ((x, 0), U \times \{0\}) \in \dcsyndetic$, as desired.

    To see that $A/k \in \dcsyndetic$, we appeal to a well-known fact (cf. \cite[Cor. 3.2]{Ye-D_function}): if a point in a system is uniformly recurrent under the transformation, then it is uniformly recurrent under every power of the transformation.  Since $(X,T)$ is minimal, the point $x$ is uniformly recurrent under $T$, and hence uniformly recurrent under $T^k$.  Therefore, the system $(\overline{T^{k \N_0} x}, T^k)$ is minimal, contains $x$, and $A/k \supseteq R_{T^k}(x, U \cap \overline{T^{k \N_0} x}) \in \dcsyndetic$, as desired.
    
    Suppose $A \in \dsyndetic$. By \cref{lemma_translates_of_dsyndetic_sets} \eqref{item_some_translates_are_dcs}, there exists $n \in \N$ such that $A - n \in \dcsyndetic$. By the first paragraph of this proof, we have $k(A - n) \in \dcsyndetic$. Therefore, by \cref{lemma_translates_of_dsyndetic_sets} \eqref{item_all_translates_are_ds}, we see $k A \supseteq k(A - n) + kn \in \dsyndetic$, whereby $kA \in \dsyndetic$, as desired.
\end{proof}

\subsubsection{Symbolic characterizations of dynamically syndetic sets}

We show in this section that the family of dynamically (central) syndetic sets is not affected by requiring return time sets to come from zero-dimensional systems.  This is accomplished by considering the indicator functions $1_A$ of sets $A \subseteq \N$ in the symbolic space $\{0,1\}^{\N}$.

As a warm-up to the challenges that arise in \cref{thm:simple_equivalent_dcS}, it is worth noting that when $(X,T)$ is a minimal system, $x \in X$, and $U \subseteq X$ is nonempty and open, the point $1_{R(x,U)} \in \{0,1\}^{\N}$ is not always uniformly recurrent under the shift map $\sigma$.  Indeed, here is an example.

\begin{example}
\label{ex:never_repeat}
    Let $\alpha \in (0, 1)$ be irrational.  Let $X \defeq \R / \Z$ and $T: X \to X$ be defined by $T(x) = x + \alpha$.  Let $U \defeq X \setminus \{\alpha\}$ and $x \defeq 0$. We see that
    \[
        1_{R(0,U)} = (1_{U}(T^n x))_{n \in \N} = (1_U( \alpha), 1_U(2 \alpha), 1_U(3 \alpha), \ldots) = (0, 1, 1, \ldots).
    \]
    Since $\alpha$ is irrational, for all $n \geq 2$, $T^n \alpha \neq \alpha$. Therefore, the function $1_{R(0,U)}$ takes the value $0$ only at $1$.  The point $1_{R(0,U)} \in [0]_1$ is not uniformly recurrent because it never returns to the neighborhood $[0]_1$.
\end{example}



In the following theorem, point \eqref{item:very_strongly_central} says that the set $A$ is \emph{very strongly central}, according to the terminology in \cite[Def. 2.10]{bergelson_hindman_strauss_2012}.

\begin{theorem}
\label{thm:simple_equivalent_dcS}
    Let $A \subseteq \N$. The following statements are equivalent.
    \begin{enumerate}
        \item \label{item:dcsyndetic} The set $A$ is dynamically central syndetic.  

        \item \label{item:very_strongly_central}
        There exists a minimal system $(X,T)$, a nonempty, open set $U \subseteq X$, and a point $x \in \overline{U}$ such that $R(x,U) \subseteq A$.
        
        
        \item \label{item:1_A_0_uniformly_recurrent} There exists a subset $B \subseteq A$ such that the point $1_{B \cup \{0\}}$ is uniformly recurrent in the full shift $(\{0, 1\}^{\N_0}, \sigma)$. In particular, the system $(X \defeq \overline{\sigma^{\N_0} 1_{B \cup \{0\}}}, \sigma)$ is minimal, the set $X \cap [1]_0$ is a clopen neighborhood of $1_{B \cup \{0\}}$, and $B = R(1_{B \cup \{0\}}, X \cap [1]_0) \subseteq A$.
    \end{enumerate}
\end{theorem}

\begin{proof}
    (\ref{item:dcsyndetic}) $\Rightarrow$ (\ref{item:very_strongly_central}) \ This follows by definition.

    (\ref{item:very_strongly_central}) $\Rightarrow$ (\ref{item:1_A_0_uniformly_recurrent}) \ Assume that $A \supseteq R(x, U)$, where $(X, T)$ is a minimal system, $U \subseteq X$ is nonempty and open, and $x \in \overline{U}$.
    
    If $x$ is a periodic point of $(X, T)$, let $n_0 \in \N$ be least such that $T^{n_0} x = x$. Since $(X, T)$ is minimal, $X = \{x, Tx, \ldots, T^{n_0 - 1} x\}$. As $X$ is finite and $x \in \overline{U}$, it must be that $x \in U$. Letting $B = n_0 \N$ be the set of positive multiples of $n_0$, we have that $A \supseteq R(x, U) \supseteq B$ and that $1_{B \cup \{0\}}$ is a periodic, hence uniformly recurrent, point in $\{0, 1\}^{\N_0}$.

    Suppose that $x$ is not a periodic point of $(X,T)$. We claim that there exists a nonempty, open set $V \subseteq U$ such that $x \in \overline{V}$ and the boundary $\partial V$ is disjoint from $T^{\N} x \defeq \{T^n x \ | \ n \in \N\}$.

    If $x \in U$, choose $V = B_\delta(x)$ for some small $\delta > 0$. Note that the boundary of $B_\delta(x)$ is the set $\{z \in X \ | \ d_X(x, z) = \delta\}$, which may be empty. Since there are uncountably many choices for $\delta$, there exists a choice so that the boundary is disjoint from $T^{\N} x$. 


    Suppose that $x \in \partial U = \overline{U} \setminus U$.  There exists a sequence of pairwise distinct points $(x_k)_{k \in \N}$ in $U \setminus \{x\}$ such that $\lim_{k \to \infty }x_k = x$.
    Inductively, choose a sequence of disjoint balls $B_{\delta_k}(x_k) \subseteq U$ with $\lim_{k \to \infty} \delta_k = 0$ such that the boundary of each $B_{\delta_k}(x_k)$ has empty intersection with $T^{\N} x$.
    Let $V = \bigcup_{k = 1}^{\infty} B_{\delta_k}(x_k)$. It is clear that $V$ is an open subset of $U$. Since $\lim_{k \to \infty }x_k = x$, we have that $\partial V = \{x\} \cup \bigcup_{k=1}^{\infty} \partial V_k$, and so the boundary of $V$ is disjoint from $T^{\N} x$. Our claim is proved.
    
    Let $B = R(x, V)$. Since $V$ is nonempty and $V \subseteq U$, we have that $B$ is nonempty and that $B \subseteq A$. It remains to show that $1_{B \cup \{0\}}$ is a uniformly recurrent point in $(\{0, 1\}^{\N_0}, \sigma)$.  By the topology on the space $\{0, 1\}^{\N_0}$, we must show that for all $L \in \N$, there exists a syndetic set of positive integers $m$ such that
    \begin{align}
    \label{eq:syndetic_return_condition}
        \text{for all $i \in \{0, \ldots, L\}$}, \ 1_{B \cup \{0\}}(m+i) = 1_{B \cup \{0\}}(i).
    \end{align}
    Let $L \in \N$.  As $\partial V \cap T^{\N} x = \emptyset$, for every $n \in \N$, the point $T^n x$ is either in $V$ or in $X \setminus \overline{V}$.
    For $n \in \{1, \ldots, L\}$, define
    \[
        W_n = \begin{cases} V & \text{ if } T^n x \in V \\ X \setminus \overline{V} & \text{ if } T^n x \in X \setminus \overline{V}. \end{cases}
    \]
    It follows that $x \in \bigcap_{n = 1}^L T^{-n} W_n$. Since each $W_n$ is open, the set $\bigcap_{n = 1}^L T^{-n} W_n$ is an open neighborhood of $x$. Because $x \in \overline{V}$, the set $V \cap \bigcap_{n = 1}^L T^{-n} W_n$ is nonempty and open. Since $(X, T)$ is minimal, the set
    \[R \left( x, V \cap \bigcap_{n = 1}^L T^{-n} W_n \right)\]
    is syndetic, and all integers $m$ in this set satisfy the condition in \eqref{eq:syndetic_return_condition}. This demonstrates that $1_{B \cup \{0\}}$ is uniformly recurrent under the shift, as desired.
    
    (\ref{item:1_A_0_uniformly_recurrent}) $\Rightarrow$ (\ref{item:dcsyndetic}) \ Suppose $A$ contains a subset $B$ for which $1_{B \cup \{0\}}$ is uniformly recurrent in $(\{0, 1\}^{\N_0}, \sigma)$. 
    The system $(X \defeq \overline{\sigma^{\N_0} 1_{B \cup \{0\}}}, \sigma)$ is minimal, and the cylinder set $X \cap [1]_0$ is open and contains $1_{B \cup \{0\}}$.  Since $A \supseteq B = R(1_{B \cup \{0\}}, X \cap [1]_0)$, we have that $A$ is dynamically central syndetic, as desired.
\end{proof}

A version of the following theorem is given in \cite[Prop. 2.3]{huang_ye_2005}.  The derivation from \cref{thm:simple_equivalent_dcS} is short, so we give it here.

\begin{theorem}
\label{thm:simple_equivalent_dS}
    Let $A \subseteq \N$.  The following statements are equivalent.
    \begin{enumerate}
        \item \label{item:dsyndetic_one} The set $A$ is dynamically syndetic.
        
        \item \label{item:1_A_uniform_recurrent_one} There exists a nonempty subset $B \subseteq A$ for which the point $1_B$ is uniformly recurrent in the full shift $(\{0, 1\}^{\N},\sigma)$.
    \end{enumerate}
\end{theorem}


\begin{proof}
    (\ref{item:dsyndetic_one}) $\Rightarrow$ (\ref{item:1_A_uniform_recurrent_one}) \ Suppose $A$ is dynamically syndetic. By \cref{lemma_translates_of_dsyndetic_sets}, there exists $k \in \N$ such that $A + k$ is dynamically central syndetic. By \cref{thm:simple_equivalent_dcS}, there exists a set $B \subseteq A + k$ such that $1_{B \cup \{0\}}$ is uniformly recurrent in $(\{0, 1\}^{\N_0},\sigma)$. Note that $B - k \subseteq A$.  Since $1_{B \cup \{0\}}$ is uniformly recurrent in $(\{0, 1\}^{\N_0},\sigma)$, we have that $1_{B}$ is a uniformly recurrent point in $(\{0, 1\}^{\N},\sigma)$.  Applying $\sigma^k$, we have that $\sigma^k 1_B = 1_{B-k}$ is uniformly recurrent, as was claimed.  

    (\ref{item:1_A_uniform_recurrent_one}) $\Rightarrow$ (\ref{item:dsyndetic_one}) \ 
    Suppose $A$ contains a nonempty subset $B$ for which $1_{B}$ is uniformly recurrent in $(\{0, 1\}^{\N},\sigma)$. Let $k = \min B$. Since $1_{B}$ is uniformly recurrent, so is $\sigma^k 1_B = 1_{B-k}$.  Since $k \in B$, we have that $1_{(B-k) \cup \{0\}}$ is uniformly recurrent in $\{0, 1\}^{\N_0}$. By \cref{thm:simple_equivalent_dcS}, the set $B - k \in \dcsyndetic$. Since $A - k \supseteq B - k$, the set $A - k$ is dynamically central syndetic.  By \cref{lemma_translates_of_dsyndetic_sets}, the set $A \supseteq (A-k) + k$ is dynamically syndetic, as desired.
\end{proof}


\subsubsection{Partitioning dynamically central syndetic sets: a proof of \cref{mainthm_partition_of_dcs_sets}}
\label{sec_partitioning_dcs_sets}


Considering some first examples of dynamically central syndetic sets, one may be led to believe that the intersection of two $\dcsyndetic$ sets is again a $\dcsyndetic$ set.  This is true in some special cases: a result of Furstenberg \cite[Thm. 9.11]{furstenberg_book_1981} implies, for example, that the intersection of the set of times of return of a distal point to a neighborhood of itself with any other $\dcsyndetic$ set is $\dcsyndetic$.  It is not, however, true in general, as we show in \cref{mainthm_partition_of_dcs_sets}.  This theorem is best contextualized as a strengthening of \cite[Thm. 2.13]{bergelson_hindman_strauss_2012}, which (when combined with \cref{thm:simple_equivalent_dcS}) says that $\N$ can be partitioned into infinitely many disjoint $\dcsyndetic$ sets.


\begin{proof}[Proof of \cref{mainthm_partition_of_dcs_sets}]
    Let $A \in \dcsyndetic$. We will prove that $A$ can be partitioned into two disjoint $\dcsyndetic$ sets. 
    Then, by repeatedly partitioning the second set into disjoint $\dcsyndetic$ sets,
    \[A = A_0 \cup A_1 = A_{0} \cup A_{10} \cup A_{11} = A_{0} \cup A_{10} \cup A_{110} \cup A_{111} = \cdots,\]
    we see that $A$ can be partitioned into infinitely many disjoint $\dcsyndetic$ sets.
    Let $(X, T)$ be a minimal system and $U \subseteq X$ be an open neighborhood of $x$ such that $A \supseteq R(x, U)$. We consider two cases.

    Case 1: the point $x$ is an isolated point of $X$. By \cref{lemma_isolated_points_and_minimal_systems}, the system $(X,T)$ is finite and periodic. Let $(\T \defeq \R / \Z, S)$ be the rotation on the one-dimensional torus by an irrational $\alpha$. Let $V_1 = (0, 1/2), V_2 = (1/2, 1)$, which are two disjoint, open subsets of $\T$ whose closures contain $0$. Then $\{x\} \times V_1$ and $\{x\} \times V_2$ are two disjoint, open subsets of $X \times \T$ whose closures contain $(x, 0)$. Note that the system $(X \times \T, T \times S)$ is minimal, and so by \cref{thm:simple_equivalent_dcS}, the sets $R((x, 0), \{x\} \times V_1)$ and $R((x, 0), \{x\} \times V_2)$ are $\dcsyndetic$. Moreover,
    \[
        A \supseteq R(x, U) \supseteq R(x, \{x\}) \supseteq R((x, 0), \{x\} \times V_1) \cup R((x, 0), \{x\} \times V_2).
    \]

    Case 2: the point $x$ is a limit point of $U$. Arguing as in the proof of \cref{thm:simple_equivalent_dcS}, there exists a sequence of disjoint open balls $B_{\delta_k}(x_k)$ in $U$ such that $x_k \to x$ and $\delta_k \to 0$. Take 
    \[
        V_0 = \bigcup_{k=1}^{\infty} B_{\delta_{2k}}(x_{2k}) \text{ and } 
        V_1 = \bigcup_{k=1}^{\infty} B_{\delta_{2k - 1}}(x_{2k - 1}).
    \]
    Then $V_0, V_1$ are disjoint open subsets of $U$ satisfying $x \in \overline{V_0} \cap \overline{V_1}$. By \cref{thm:simple_equivalent_dcS}, the sets $R(x, V_0)$ and $R(x, V_1)$ are $\dcsyndetic$ sets and, by construction, they are disjoint subsets of $A$. Put $A_0 \defeq R(x,V_0)$ and $A_1 \defeq A \setminus R(x,V_0)$ to see that $A = A_0 \cup A_1$ is a disjoint partition of $A$ into $\dcsyndetic$ sets, as desired.
\end{proof}


\subsection{Central syndetic sets}

In this section, we define the family of central syndetic sets and give a combinatorial characterization of them.


\begin{definition}
\label{def_central_syndetic}
    A set $A \subseteq \N$ is a \emph{central syndetic} if it is a member of a syndetic, idempotent filter on $\N$.
\end{definition}

Canonical examples of central syndetic sets are provided by dynamically central syndetic sets: return time sets $R(x,U)$ where $(X,T)$ is a minimal system, $x \in X$, and $U \subseteq X$ is an open neighborhood of $x$.  That sets of this form are central syndetic is proven in the ``only if'' direction of \cref{thm_dcs_iff_cs} below.  The ``if'' direction of that theorem shows that these are essentially the only examples of central syndetic sets.

Let us inject a brief historical remark to motivate the nomenclature. Central sets -- described first by Furstenberg  \cite[Def. 8.3]{furstenberg_book_1981} in a dynamical context -- have since been characterized as members of minimal, idempotent ultrafilters \cite{glasner_1980,bergelson_hindman_1990,shi_yang_1996}. To discuss the dynamical and combinatorial families separately, it has become customary (see, eg., \cite[Sec. 19.3]{hindman_strauss_book_2012}) to call Furstenberg's sets ``dynamically central'' and members of minimal, idempotent ultrafilters ``central.''  We take a similar approach here: ``dynamically central syndetic'' sets are defined dynamically as above, while ``central syndetic'' sets are defined combinatorially as members of syndetic, idempotent filters.  We will demonstrate in \cref{thm_dcs_iff_cs} that these families are the same.


\subsubsection{Central syndetic sets are central}

In this section, we justify the terminology ``central syndetic'' by showing that central syndetic sets are, in fact, both central and syndetic.  That they are syndetic follows immediately from the definition -- they are members of syndetic filters -- so we need only to show that they are central. In what follows, for a family $\family$ on $\N$, we denote by $\overline{\family}$ the set $\bigcap_{A \in \family} \overline{A}$ in $\beta \N$.

\begin{lemma}
\label{lemma_syndetic_contained_in_minimal}
    For all syndetic filters $\filter$ on $\N$ and all minimal left ideals $L \subseteq \beta \N$, we have $L \cap \overline{\filter} \neq \emptyset$.
\end{lemma}

\begin{proof}
    We will use the following well-known fact (\cite[Thm. 4.48]{hindman_strauss_book_2012}): for all syndetic sets $A \subseteq \N$ and all minimal left ideals $L \subseteq \beta \N$, we have $\overline{A} \cap L \neq \emptyset$.

    Let $\filter$ be a syndetic filter on $\N$, and let $L \subseteq \beta \N$ be a minimal left ideal.  Consider the set $\{\overline{A} \cap L \ | \ A \in \filter\}$.  By the fact from the previous paragraph, the set consists of closed, nonempty subsets of $L$.  Since $\filter$ is a filter, the set has the finite intersection property.  Since $L$ is compact, we have that $\overline{\filter} \cap L$ is nonempty, as desired.
\end{proof}


\begin{theorem}
\label{thm_sifs_are_contained_in_idempotents}
    For all syndetic, idempotent filters $\filter$ on $\N$ and all minimal left ideals $L \subseteq \beta \N$, there exists an idempotent ultrafilter in $L$ containing $\filter$.
\end{theorem}

\begin{proof}
    Let $\filter$ be a syndetic, idempotent filter on $\N$. By \cref{lemma_syndetic_contained_in_minimal}, there exists a minimal ultrafilter $p \in \beta \N$ such that $\filter \subseteq p$.  Define $L = \beta \N + p$.  Since $p$ is minimal, the set $L$ is a minimal left ideal, and hence subsemigroup, of $(\beta \N, +)$.  By \cite[Thm. 2.6]{hindman_strauss_book_2012}, the set $L$ is compact.  Therefore, the set $L$ is a compact subsemigroup of $(\beta \N, +)$ containing $p$.

    It follows from the proof of \cref{thm_equiv_forms_of_max_idempot_filter} that the set $\overline{\filter}$ is a compact subsemigroup of $(\beta \N, +)$.  It clearly contains $p$.

    Define $X = L \cap \overline{\filter}$.  Since both $L$ and $\overline{\filter}$ are compact subsemigroups of $(\beta \N, +)$ containing $p$, so is $X$.  By \cite[Thm. 2.5]{hindman_strauss_book_2012}, there exists an idempotent ultrafilter $q \in X$.  Since $X \subseteq L$, we see that $q$ is minimal.  Since $X \subseteq \overline{\filter}$, we see that $\filter \subseteq q$.  Thus, we have shown that $\filter$ is contained in a minimal, idempotent ultrafilter in $L$, as desired.
\end{proof}

\begin{theorem}
\label{cor_central_syndetic_is_strongly_central}
    Let $A \subseteq \N$ be central syndetic.  For all minimal left ideals $L \subseteq \beta \N$, there exists an idempotent $p \in L$ such that $A \in p$.  In particular, for all thick sets $H \subseteq \N$, the set $A \cap H$ is central.
\end{theorem}

\begin{proof}
    Since $A$ is central syndetic, there exists a syndetic, idempotent filter $\family$ containing $A$. Let $L \subseteq \N$ be a minimal left ideal. By \cref{thm_sifs_are_contained_in_idempotents}, there exists an idempotent $p \in L$ such that $\family \subseteq p$.  In particular, we have that $A \in p$.

    If $H \subseteq \N$ is thick, by \cite[Thm. 4.48]{hindman_strauss_book_2012}, there exists a minimal left ideal $L \subseteq \overline{H}$.  By the reasoning in the previous paragraph, there exists an idempotent $p \in L$ containing $A$.  Therefore, $A \cap H \in p$, whereby the set $A \cap H$ is central.
\end{proof}

\begin{remark}
    In \cite{bergelson_hindman_strauss_2012}, a set $A \subseteq \N$ is called \emph{strongly central} if it has the property that for all minimal left ideals $L \subseteq \beta \N$, there exists an idempotent $p \in L$ containing $A$.  Thus, \cref{cor_central_syndetic_is_strongly_central} shows that central syndetic sets are strongly central.  It will follow, however, from our main result in this section, \cref{thm_dcs_iff_cs}, that strongly central sets need not be central syndetic.  Indeed, a corollary of \cref{thm:simple_equivalent_dcS,thm_dcs_iff_cs} is that the family of central syndetic sets is the same as the family of \emph{very strongly central sets}, in the terminology of \cite{bergelson_hindman_strauss_2012}.  An example of a strongly central, but not very strongly central, set is given in \cite[Thm. 2.17]{bergelson_hindman_strauss_2012}.
\end{remark}

   

\subsubsection{A combinatorial characterization of central syndetic sets}
\label{sec_combo_condition_for_central_syndetic}

In this section and the next, we give a combinatorial characterization of central syndetic sets.  Consider the following statement for a set $B \subseteq \N$:
\begin{align}
    \label{eqn_def_of_cssd}
    \text{for all finite $F \subseteq B$, \ the set } B \cap \bigcap_{f \in F} \big( B - f \big) \text{ is syndetic}.
\end{align}
The following theorem shows that the condition in \eqref{eqn_def_of_cssd} captures membership in a syndetic, idempotent filter.


\begin{theorem}
\label{thm_conditions_for_central_syndetic}
    Let $A \subseteq \N$.
    \begin{enumerate}
        \item \label{item_central_syndetic_implies_comb_condition}
        If $A$ is central syndetic and $\filter$ is a syndetic, idempotent filter to which $A$ belongs, then the set $A \cap (A-\filter)$ satisfies the condition in \eqref{eqn_def_of_cssd}.

        \item \label{item_comb_condition_implies_central_syndetic}
        If $A$ satisfies the condition in \eqref{eqn_def_of_cssd}, then the family
        \[\filter \defeq \bigupclose \bigg\{ A \cap \bigcap_{f \in F} (A - f) \ \bigg | \ F \subseteq A \text{ is finite} \bigg\}\]
        is a syndetic, idempotent filter that contains $A$, whereby $A$ is central syndetic.
    \end{enumerate}
    In particular, the set $A$ is central syndetic if and only if some subset of it satisfies the condition in \eqref{eqn_def_of_cssd}.
\end{theorem}

\begin{proof}
    \eqref{item_central_syndetic_implies_comb_condition} \ Suppose $A$ belongs to a syndetic, idempotent filter $\filter$.  Define $A_1 \defeq A \cap (A-\filter)$.  We claim that this set satisfies the condition in \eqref{eqn_def_of_cssd}.

    First we will show that $A_1 \in \filter$ and $A_1 \subseteq A_1 - \filter$.  That $A_1 \in \filter$ follows because $A \in \filter$ and $A-\filter \in \filter$ (since $\filter$ is idempotent).  Because $\filter \subseteq \filter + \filter$, using the algebra in \cref{lemma_filter_algebra}, we see
    \begin{align*}
        A_1 = A \cap (A-\filter) \subseteq A - \filter &\subseteq (A - \filter) \cap \big(A - (\filter + \filter) \big) \\
        &= (A - \filter) \cap \big((A - \filter) - \filter \big) \\
        &= \big( A \cap (A - \filter) \big) - \filter \\
        &= A_1 - \filter.
    \end{align*}

    To see that the condition in \eqref{eqn_def_of_cssd} is satisfied, let $F \subseteq A_1$ be finite.  Since $F \subseteq A_1 - \filter$, the set $\bigcap_{f \in F} (A_1 - f) \in \filter$.  Since $A_1 \in \filter$, we have that $A_1 \cap \bigcap_{f \in F} (A_1 - f) \in \filter$, whereby it is syndetic, as desired.

    \eqref{item_comb_condition_implies_central_syndetic} \ That $\filter$ is a filter is immediate from its definition.  That $A \in \filter$ follows immediately from the fact that $\filter$ is upward closed.

    To see that $\filter$ is syndetic, it suffices to show that for all $F \subseteq A$ finite, the set $B \defeq A \cap \bigcap_{f \in F} (A - f)$ is syndetic.  But this is immediate from \eqref{eqn_def_of_cssd}.

    To see that $\filter$ is idempotent, we must show that $\filter \subseteq \filter + \filter$.  Since $\filter + \filter$ is a filter, it suffices to show that for all $a \in A \cup\{0\}$, the set $A - a \in \filter + \filter$.
    
    Let $a \in A \cup\{0\}$.  We will show that $(A - a) - \filter \in \filter$.  Since $A - a \in \filter$, it suffices to show that $A - a \subseteq (A - a) - \filter$.  This holds since for any $n \in A - a$, we have that $(A - a) - n = A - (n + a)$, which is a member of $\filter$ by the definition of $\filter$ since $n+a \in A$.
\end{proof}

\begin{remark}
    It is a short exercise to show that the condition in \eqref{eqn_def_of_cssd} is equivalent to the following: \emph{there exists a chain of syndetic sets $A \supseteq A_1 \supseteq A_2 \supseteq A_3 \supseteq \cdots $ with the property that for all $n \in \N$ and all $a \in A_n$, there exists $m \in \N$ such that $A_n - a \supseteq A_m$.}  The combinatorial characterization of central sets in \cite[Thm. 14.25]{hindman_strauss_book_2012} says the same thing with with ``syndetic'' replaced by ``collectionwise piecewise syndetic.''
\end{remark}

We finish this section by showing that sets $B \subseteq \N$ satisfying the condition in \eqref{eqn_def_of_cssd} satisfy the ostensibly stronger condition:
\begin{align}
    \label{eqn_stronger_than_cssd}
    \text{for all finite $F \subseteq B$ and all $m \in \N$, the set } B \cap \bigcap_{f \in F} \big( B - f \big) \cap m \N \text{ is syndetic}.
\end{align}
This upgrade will help us with the proof of \cref{thm_dcs_iff_cs} in the next section.



\begin{lemma}
\label{lemma_upgraded_cssd_property}
    A set $A \subseteq \N$ satisfies the condition in \eqref{eqn_def_of_cssd} if and only if it satisfies the condition in \eqref{eqn_stronger_than_cssd}.
\end{lemma}

\begin{proof}
    If $A$ satisfies the condition in \eqref{eqn_stronger_than_cssd}, then it clearly satisfies the one in \eqref{eqn_def_of_cssd}.

    Conversely, suppose that $A \subseteq \N$ satisfies the condition in \eqref{eqn_def_of_cssd}. To see that it satisfies \eqref{eqn_stronger_than_cssd}, let $F \subseteq A$ be finite and $m \in \N$. To see that the set $A \cap \bigcap_{f \in F} (A-f) \cap m\N$ is syndetic, we will show that it has nonempty intersection with all thick sets.
    
    Let $H \subseteq \N$ be thick. By \cref{thm_conditions_for_central_syndetic}, we see that the set $A \cap \bigcap_{f \in F} (A-f)$ belongs to a syndetic, idempotent filter, and hence is central syndetic.  It follows from \cref{cor_central_syndetic_is_strongly_central} that the set $A \cap \bigcap_{f \in F} (A-f) \cap H$ is central.  The set $m \N$ is $\IP^*$, hence $\central^*$, and so the set $A \cap \bigcap_{f \in F} (A-f) \cap m \N \cap H$ is nonempty, as was to be shown.
\end{proof}


\subsection{Central and dynamically central syndetic sets}
\label{sec_char_of_dy_central_synd_sets}

We show in this section that the family of dynamically central syndetic sets and central syndetic sets are the same.  Combined with \cref{thm_conditions_for_central_syndetic}, we get a combinatorial characterization of the family $\dcsyndetic$.

\begin{theorem}
\label{thm_dcs_iff_cs}
    A set $A \subseteq \N$ is dynamically central syndetic if and only if it is central syndetic, that is, a member of a syndetic, idempotent filter on $\N$.
\end{theorem}

The proof of \cref{thm_dcs_iff_cs} is broken into two parts, the first of which we give now.

\begin{proof}[Proof of the ``only if'' statement in \cref{thm_dcs_iff_cs}]
    By assumption, there exists a minimal system $(X,T)$, a point $x \in X$, and an open set $U \subseteq X$ containing $x$ such that $R(x,U) \subseteq A$. Define
    \[
        \filter \defeq \upclose \big\{ R(x,V) \ \big| \ V \subseteq X \text{ is an open set containing $x$} \big\}.
    \]
    We claim that $\filter$ is a syndetic, idempotent filter containing $A$.

    That $\filter$ is a filter, that $A \in \filter$, and that $\filter$ is syndetic are all immediate.  To see that $\filter$ is idempotent, it suffices by \cref{rmk_showing_idempotency} to show that for all open $V \subseteq X$ with $x \in V$, we have that $R(x,V) \subseteq R(x,V) - \filter$.  Since $R(x,V) \in \filter$, this shows that $R(x,V) - \filter \in \filter$, as desired.

    Let $V \subseteq X$ be open and contain $x$.  To see that $R(x,V) \subseteq R(x,V) - \filter$, let $n \in R(x,V)$.  Thus, $x \in T^{-n} V$.  Since $T^{-n} V$ is an open neighborhood of $x$, we have that $R(x,V) - n = R(x,T^{-n} V) \in \filter$, as desired.
\end{proof}

The rest of this section is devoted to a proof of the converse, namely that central syndetic sets are dynamically central syndetic.  The argument requires a rather lengthy setup.

\subsubsection{The main idea}
\label{sec_main_idea_of_shift_punch}

Here is the main idea behind the proof of the ``if'' statement in \cref{thm_dcs_iff_cs}.  By \cref{thm_conditions_for_central_syndetic}, we may assume that the set $A$ satisfies the condition in \eqref{eqn_def_of_cssd}.

Our aim is to show that the set $A$ contains a set of the form $R(x,U)$, where the point $x$ belongs to the open set $U$ in a minimal system.  In view of \cref{thm:simple_equivalent_dcS}, it suffices to find a subset $B \subseteq A$ for which $1_{B \cup \{0\}} \in \{0, 1\}^{\N_0}$ is uniformly recurrent under the shift map $\sigma$.  Note that the condition in \eqref{eqn_def_of_cssd} is not sufficient for $1_{A \cup \{0\}}$ to be uniformly recurrent under $\sigma$.  (Indeed, the set $\N_0 \setminus \{n^2 \ | \ n \in \N\}$ satisfies the condition in \eqref{eqn_def_of_cssd} but has an indicator function that is not uniformly recurrent under $\sigma$.)  To find the set $B$, the idea is to modify the shift map $\sigma$ into a ``shift-punch'' map $\kappa$ under which the point $1_A$ becomes uniformly recurrent.


We append the integer $0$ to $A$ and consider it as a subset of $\N_0$. The ``punch map'' $\pi: \{0,1\}^{\N_0} \to \{0,1\}^{\N_0}$, which depends on the set $A$, is described roughly as follows.  At $1_C \in \{0,1\}^{\N_0}$, we find the maximal $N \in \N$ such that $A \cap [N] \subseteq C$.  Then, $\pi(1_C) = 1_D$, where $D \cap [N] = A \cap [N]$ and $D \cap (\N_0 \setminus [N]) = C \cap (\N_0 \setminus [N])$.  Roughly, the map $\pi$ ``punches out'' or ``deletes'' the elements of $C$ on the interval $[N]$ that are not in $A \cap [N]$.  The map $\pi$ continuously moves points of $\{0,1\}^{\N_0}$ closer to $1_A$.

The shift-punch map $\kappa$ is defined to be $\pi \circ \sigma$.  The dynamics of the shift-punch system $(\{0,1\}^{\N_0}, \kappa)$ is thus a shift followed by a punch.  The idea is to show that in this system, the point $1_A$ is uniformly recurrent.  This happens because the set $A$ satisfies the condition in \eqref{eqn_def_of_cssd} and thus syndetically often contains an initial interval of itself.  Keeping track of which elements of $A$ are deleted as the shift-punch map is repeatedly applied to $1_A$, we end up with a subset $B \subseteq A$. We will show that the point $1_B \in \{0,1\}^{\N_0}$ is uniformly recurrent under the usual shift map, as desired.


The outline just given is difficult to control in practice because the punches can seemingly overlap in complicated ways.  To overcome this, we find it useful to have more control over the length of intervals on which the punches occurs.  Thus, the actual shift-punch system is defined as a skew product over a 2-adic odometer.  The odometer functions only to control the length of the punch intervals, which are restricted to occur only on dyadic intervals.  Dyadic intervals are convenient because a nonempty intersection implies containment.  This control on the punch intervals allows us to realize the outline above.

\subsubsection{Notation and spaces}

In the rest of this section, we will alternate between working in $\N$ and $\N_0 \defeq \{0,1,2,\ldots\}$.  Recall that for $N \in \N_0$, we write $[N] \defeq \{0, \ldots, N-1\}$. When $A \subseteq \N_0$, we compute the translate $A-n$ as a subset of $\N_0$.

It will be convenient to append some elements to $\N_0$.  Thus, we define $\N_{0,\pm \infty} \defeq \N_0 \cup \{-\infty, \infty\}$.  Subspaces, such as $\N_{0,\infty} \defeq \N_0 \cup \{\infty\}$, are defined analogously.  We give $\N_{0,\pm \infty}$ the topology that makes the bijection $\N_{0,\pm \infty} \to \{2\} \cup \{1/n \ | \ n \in \N\} \cup \{0\}$, where $-\infty \mapsto 2$, $x \mapsto 1/(x+1)$, and $\infty \mapsto 0$ a homeomorphism.  Thus, every point of $\N_{0,\pm \infty}$, except for $\infty$, is isolated.  Subspaces of $\N_{0,\pm \infty}$ are given the subspace topology.

\subsubsection{Dyadic valuation and intervals}

Denote by $\nu_2: \Z \to \N_{0,\infty}$ the two-adic valuation on $\Z$, defined by $\nu_2(n) \defeq \sup \{k \in \N_0 \ | \ 2^k \text{ divides } n \}$. A finite interval $\{a, a+1, \ldots, a+\ell - 1\}$ in $\N_0$ is a \emph{dyadic interval} if $\ell = 2^k$ for some $k \in \N_0$ and $\ell \mid a$.  Denote by $\cubes_{2^n}$ the set of dyadic intervals of cardinality $2^n$.  We will say that $\N_0$ is the dyadic interval of infinite cardinality.


The following lemma follows from standard facts about dyadic intervals.  We provide a proof for completeness and ease of reference for later.

\begin{lemma}
\label{lemma_basic_dyadic_cubes_facts}
Let $Q$ and $P$ be finite dyadic intervals in $\N_0$.
\begin{enumerate}
    \item \label{item_dyadic_prop_two}
    $Q \cap 2^{\nu_2(\min Q)}\N = \{\min Q\}$.
    \item \label{item_dyadic_prop_three}
    For all $\ell \in \N_0$, if $\big(Q \setminus \{\min Q\} \big) \cap 2^{\ell} \N \neq \emptyset$, then $|Q| \geq 2^{\ell + 1}$.
    \item \label{item_dyadic_prop_one}
    If $P \cap Q \neq \emptyset$, then either $P \subseteq Q$ or $Q \subseteq P$.
\end{enumerate}
\end{lemma}

\begin{proof}
    
    \eqref{item_dyadic_prop_two} \ There exists $n \in \N$ and $0 \leq k \leq \nu_2(n)$ such that $Q = n + [2^k]$.  Since $Q$ is an interval beginning at $n \in 2^{\nu_2(n)} \N$ with $|Q| = 2^k \leq 2^{\nu_2(n)}$, we have that $Q \cap 2^{\nu_2(n)}\N = \{n\}$.

    \eqref{item_dyadic_prop_three} \ Write $Q = n + [2^k]$ as in \eqref{item_dyadic_prop_two}. Let $\ell \in \N_0$, and suppose that $\big(Q \setminus \{n\} \big) \cap 2^{\ell} \N$ is nonempty.  It follows from the previous paragraph that $\ell < \nu_2(n)$.  Let $a \in \big(Q \setminus \{n\} \big) \cap 2^{\ell} \N$.  Since $|Q| = 2^k > a-n$ and $\nu_2(a-n) \geq \ell$, we see that $k \geq \ell + 1$.  Therefore, $|Q| \geq 2^{\ell + 1}$, as desired.

    \eqref{item_dyadic_prop_one} \ Let $P$ and $Q$ be dyadic intervals.  Note that if $\min P = \min Q$, then $P \subseteq Q$ or $Q \subseteq P$, and we are done.  Otherwise, we will show that if $\min P \in Q \setminus \{\min Q\}$, then $P \subseteq Q$.  It will follow that either $P$ and $Q$ are disjoint or that one is contained in the other.  
    
    Suppose for a contradiction that $\min P \in Q \setminus \{\min Q\}$ and $P \not\subseteq Q$.  Then $\max Q + 1 = \min Q + |Q| \in P$ and is divisible by $|Q|$.  It follows from \eqref{item_dyadic_prop_three} that $|P| \geq 2 |Q|$.  But since $P \not\subseteq Q$ and $|P|$ divides $\min P$, we see by \eqref{item_dyadic_prop_three} that $|Q| \geq 2 |P|$, a contradiction.
\end{proof}

We denote by $\cubes_{2^n}^*$ the collection of subsets of $\N$ that have nonempty intersection with all members of $\cubes_{2^n}$.  (This notation matches notation for the dual family defined in \cref{sec_dual_families}, but note that $\cubes_{2^n}$ is not a family since it is not upward closed.)


\begin{lemma}
\label{lemma_syndetic_characterization}
A set $A \subseteq \N_0$ is syndetic if and only if there exists $k \in \N$ such that $A \in \cubes_{2^k}^*$.
\end{lemma}

\begin{proof}
    This is left to the reader as an easy exercise.
\end{proof}

\subsubsection{The 2-adic odometer}
\label{sec_two_adic_odometer}

We denote by $(\odometer,\orotation)$ the 2-adic odometer, the inverse limit of the family of rotations $(\Z / 2^k \Z, \allowbreak +1)$, $k \in \N_0$, as topological dynamical systems.  We do not have need to specify a metric on $\odometer$ explicitly; it suffices for our purposes to know that two elements of $\odometer$ are near if and only if their projections to $\Z / 2^k \Z$ agree for some large value of $k$.

For $n \in \Z$, we denote by $n$ the element of $Z$ that projects, for all $k \in \N_0$, to $n \in \Z / 2^k \Z$. This association gives a dense embedding of $\Z$ into $\odometer$.  By a convenient abuse of notation, we write $\Z \subseteq \odometer$, identifying $\Z$ with this copy in $\odometer$.

The two-adic valuation $\nu_2: \Z \to \N_{0,\infty}$ is uniformly continuous in the subspace topology that $\Z$ inherits from the space $\odometer$.  Indeed, if the images of integers $n$ and $m$ are the same in $\Z / 2^k \Z$ for a large value of $k$, then $2^k$ divides $n-m$, so $\nu_2(n)$ is large if and only if $\nu_2(m)$ is large.  Therefore, the map $\nu_2: \Z \to \N_{0,\infty}$ extends uniquely to a continuous map $\nu_2: \odometer \to \N_{0,\infty}$.


\subsubsection{The punch maps}
\label{sec_punch_maps}

Let $X = \{0, 1\}^{\N_0}$, and fix $A \subseteq \N_0$.  We will define a family of continuous maps $\pi_n: X \to X$, $n \in \N_{0,\infty}$, depending on $A$ that drive the shift-punch dynamics described in the next section.


First, we define a map $\nu_A: X \to \N_{0,\pm\infty}$ in the following way.  For $\omega \in X$,
\begin{align}
    \label{eqn_def_of_M_set}
    \nu_A(\omega) \defeq \sup \big\{n \in \N_0 \ \big| \ A \cap [2^{n}] \subseteq \supp(\omega) \big\},
\end{align}
where the supremum of the empty set is $-\infty$.  Note that $\nu_A(\omega) = -\infty$ if and only if $0 \in A$ and $\omega(0) = 0$.  Also, note that $\nu_A(\omega) = \infty$ if and only if $A \subseteq \supp(\omega)$.

For $n \in \N_{0,\infty}$ and $\omega \in X$, we define the \emph{punch interval}
\[I_n(\omega) \defeq \big[2^{\min(\nu_A(\omega),n)} \big],\]
where $[2^{-\infty}]$ and $[2^{\infty}]$ are interpreted to be the empty set and $\N_0$, respectively. Then, we define the \emph{punch map} $\pi_n: X \to X$ by
\[\big(\pi_n (\omega) \big)(i) = \begin{cases} 1_A(i) & \text{ if $i \in I_n(\omega)$} \\ \omega(i) & \text{ if $i \in \N_0 \setminus I_n(\omega)$} \end{cases}.\]
Informally, if $A \cap [2^m] \subseteq \supp(\omega)$, where $m \in \{-\infty, 0, 1, \ldots, n\}$ is maximal as such, then $\pi_n(\omega)$ is defined so that $\supp(\pi_n(\omega)) \cap [2^m] = A \cap [2^m]$ and $\supp(\pi_n(\omega)) \cap (\N_0 \setminus [2^m]) = \supp(\omega) \cap (\N_0 \setminus [2^m])$. The word ``punch'' was chosen since the punch maps ``punch out'' ones by changing them to zeroes.


\begin{lemma}
\label{lemma_continuity_of_family_of_punches}
    Each of the maps $\pi_n: X \to X$, $n \in \N_{0,\infty}$, is continuous, and the map $\N_{0,\infty} \to C(X,X)$ given by $n \mapsto \pi_n$ is continuous, where $C(X,X)$ has the supremum topology.
\end{lemma}

\begin{proof}
    Let $n \in \N_{0,\infty}$.  To see that $\pi_n: X \to X$ is continuous, let $\omega \in X$ and $N \in \N$.  If $\nu_A(\omega) \in \N_{0,-\infty}$, then for all $\omega' \in X$ sufficiently close to $\omega$, the points $\omega'$ and $\omega$ agree on $[N]$ and $\nu_A(\omega') = \nu_A(\omega)$.  If, on the other hand, $\nu_A(\omega) = \infty$, then for all $\omega' \in X$ sufficiently close to $\omega$, the points $\omega'$ and $\omega$ agree on $[N]$ and $\nu_A(\omega') \geq N$.  In both cases, we see that $\pi_n(\omega')$ and $\pi_n(\omega)$ agree on $[N]$.  By the topology of the space $X$, this shows that $\pi_n$ is continuous.

    To show that the map $\N_{0,\infty} \to C(X,X)$ given by $n \mapsto \pi_n$ is continuous, it suffices by the topology on $\N_{0,\infty}$ to show that $\lim_{n \to \infty} \pi_n = \pi_\infty$, that is, for all $\eps > 0$ and all sufficiently large $n \in \N$,
    \[\sup_{\omega \in X} d_X\big(\pi_n(\omega), \pi_\infty(\omega) \big) < \eps.\]
    Let $\eps > 0$ and $\omega \in X$. If $\nu_A(\omega) = -\infty$, then for all $n \in \N_0$, $\pi_\infty(\omega) = \pi_n(\omega)$. If $\nu_A(\omega) \in \N_0$, then for all $n \in \N_0$ with $n \geq \nu_A(\omega)$, $\pi_\infty(\omega) = \pi_n(\omega)$. If $\nu_A(\omega) = \infty$, then $\pi_n(\omega)$ and $\pi_\infty(\omega)$ match on $[2^n]$, because they both match $1_A$ on $[2^n]$.  Thus, in all cases, for $n \in \N_0$ sufficiently large, $d_X(\pi_n(\omega), \pi_\infty(\omega)) < \eps$, as was to be shown.
\end{proof}

The following lemma gives finer information on the continuity of $\pi_n$, $n \in \N_0$.

\begin{lemma}
\label{lemma_continuity_of_pins}
    Let $n, \ell \in \N_0$ with $\ell \geq 2^n$.  If $\omega, \xi \in X$ agree on $[\ell]$, then $I_n(\omega) = I_n(\xi)$, and $\pi_n(\omega)$ and $\pi_n(\xi)$ agree on $[\ell]$.
\end{lemma}


\begin{proof}
    Suppose that $\omega, \xi \in X$ agree on $[\ell]$.  We consider the following cases.
    \begin{enumerate}
        \item Case 1: $\nu_A(\omega) = -\infty$ or $\nu_A(\xi) = -\infty$.  In this case, since $\omega$ and $\xi$ agree on $[\ell]$, we have that $\nu_A(\omega) = \nu_A(\xi) = -\infty$ and $I_n(\omega) = I_n(\xi) = \emptyset$.  Since $\omega$ and $\xi$ are unchanged by $\pi_{n}$, we have that $\pi_{n} (\omega)$ and $\pi_{n} (\xi)$ agree on $[\ell]$.

        \item Case 2: $0 \leq \nu_A(\omega) < n$.  By the definition of $\nu_A$, we have that $A \cap [2^{\nu_A(\omega)}] \subseteq \supp(\omega)$ but $A \cap [2^{\nu_A(\omega)+1}] \not\subseteq \supp(\omega)$. Since $2^{\nu_A(\omega)} < 2^{n} \leq \ell$, we have that $[2^{\nu_A(\omega)+1}] \subseteq [\ell]$. Since $\omega$ and $\xi$ agree on $[\ell]$, we see that $\nu_A(\omega) = \nu_A(\xi)$, and hence that $I_n(\omega) = I_n(\xi)$.  Since $\pi_{n}$ changes $\omega$ and $\xi$ on the same interval contained in $[\ell]$, we see that $\pi_{n} (\omega)$ and $\pi_{n} (\xi)$ agree on $[\ell]$, as desired.

        \item Case 3: $0 \leq \nu_A(\xi) < n$. We argue just as in Case 2 to arrive at the same conclusion.

        \item Case 4: $\nu_A(\omega) \geq n$ and $\nu_A(\xi) \geq n$. In this case, we see that $I_n(\omega) = I_n(\xi) = [2^{n}]$.  Since $2^{n} \leq \ell$ and $\omega$ matches $\xi$ on $[\ell]$, we see that $\pi_{n} (\omega)$ and $\pi_{n} (\xi)$ agree on $[\ell]$, as desired.
    \end{enumerate}
    This concludes the casework and the proof of the lemma.
\end{proof}

\subsubsection{The shift-punch system}
\label{sec_punch_system}


The \emph{shift-punch system} $(\odometer \times X, \kappa)$ is the topological skew product system defined by the \emph{shift-punch map} $\kappa: \odometer \times X \to \odometer \times X$,
\[\kappa(z,\omega) \defeq \big(\orotation z, \pi_{\nu_2( \orotation z)}(\sigma \omega) \big).\]
Though not apparent from the notation, the punch maps $\pi_n$, the shift-punch map $\kappa$, and the shift-punch system $(\odometer \times X, \kappa)$ all depend highly on the set $A$.

\begin{lemma}
    The shift-punch system $(\odometer \times X, \kappa)$ is a topological dynamical system.
\end{lemma}

\begin{proof}
    The space $Z \times X$ is compact, so we need only to show that the map $\kappa$ is continuous.  Because $\kappa$ is a skew product, it suffices to check that the map $\odometer \to C(X,X)$ defined by $z \mapsto \pi_{\nu_2(\orotation z)} \circ \sigma$ is continuous.  This follows from the fact that the shift $\sigma: X \to X$ is continuous and that the map $z \mapsto \pi_{\nu_2(\orotation z)}$ is a composition of three continuous maps: $\orotation: \odometer \to \odometer$, $\nu_2: \odometer \to \N_{0,\infty}$, and, by \cref{lemma_continuity_of_family_of_punches}, $\pi: \N_{0,\infty} \to C(X,X)$.
\end{proof}


We are interested primarily in the orbit of the point $(0,1_A)$ under the shift-punch map.  In the shift-punch dynamics, it will be important to keep track of where the punches (changes from 1 to 0) occur.  The following notation will assist with that.  Define $\bpunch^{(0)} = 1_{\N_0}$ and $\apunch^{(0)} = 1_A$, and for $n \in \N$, define $\bpunch^{(n)}, \apunch^{(n)} \in X$ to be such that
\begin{align*}
    (\orotation \times \sigma) \kappa^{n-1}(0,1_A) &= \big(n, \bpunch^{(n)}\big),\\
    \kappa^n (0,1_A) &= \big(n, \apunch^{(n)} \big).
\end{align*}
By the definition of $\kappa$, we see that for all $n \in \N_0$, $\apunch^{(n)} = \pi_{\nu_2(n)} \bpunch^{(n)}$ and $\bpunch^{(n+1)} = \sigma \apunch^{(n)}$.
Thus, the points $\bpunch^{(n)}$ and $\apunch^{(n)}$ are the second coordinates of $\kappa^{n}(0, 1_A)$ before and after, respectively, the punch by $\pi_{\nu_2(n)}$.

The \emph{punch window} at $n \in \N_0$ is defined to be
\[W(n) \defeq n + I_{\nu_2(n)}\big( \bpunch^{(n)} \big),\]
where $W(n)$ is understood to be the empty set if $I_{\nu_2(n)}\big( \bpunch^{(n)} \big)$ is the empty set.  Note that $W(0) = \N_0$; that $W(n)$ is an interval beginning at $n$ of length a divisor of $2^{\nu_2(n)}$ (and, hence, is a dyadic interval); and that $\apunch^{(n)}$ and $1_A$ agree on the interval $W(n) - n$.  Informally, we note that $\supp(\bpunch^{(n)}) \subseteq A-n$ (as $\bpunch^{(n)}$ is the result of the first $n-1$ shift-punches and a single shift) just before $\pi_{\nu_2(n)}$ is applied to change $\bpunch^{(n)}$ on the punch interval $I_{\nu_2(n)}\big( \bpunch^{(n)} \big)$ to yield $\apunch^{(n)}$, the second coordinate of $\kappa^n (0,1_A)$.  The interval $W(n)$ is the interval in the original set $A$ at which that punch occurs.

\begin{lemma}
\label{lemma_explicit_description_of_omega_and_kappa}
    For all $n \in \N_0$,
    \begin{enumerate}
        \item \label{item_explicit_description_of_omega_and_kappa_one}
    \begin{align}
        \label{eqn_set_indicated_by_omega_n}
        \supp(\bpunch^{(n)}) = \bigcap_{m=0}^{n-1} \Big( \big( A - (n-m) \big) \cup \big( \N_0 \setminus \big( W(m)-n \big) \big) \Big),
    \end{align}
        where the empty intersection is interpreted as $\N_0$;

        \item \label{item_explicit_description_of_omega_and_kappa_two}
    \begin{align}
        \label{eqn_set_indicated_by_kappa_n}
        \supp(\apunch^{(n)}) = \bigcap_{m=0}^{n} \Big( \big( A - (n-m) \big) \cup \big( \N_0 \setminus \big( W(m)-n \big) \big) \Big).
    \end{align}
    \end{enumerate}
\end{lemma}

\begin{proof}
    We will prove both \eqref{item_explicit_description_of_omega_and_kappa_one} and \eqref{item_explicit_description_of_omega_and_kappa_two} simultaneously by induction on $n \in \N_0$. The base case $n=0$ follows by set algebra, recalling that, by definition, $\bpunch^{(0)} = 1_{\N_0}$ and $\apunch^{(0)} = 1_A$.
    
    Suppose both \eqref{item_explicit_description_of_omega_and_kappa_one} and \eqref{item_explicit_description_of_omega_and_kappa_two} hold for some $n \in \N_0$. Let $i \in \N_0$.  Since $\bpunch^{(n+1)} = \sigma \apunch^{(n)}$, we see that $\big(\bpunch^{(n+1)} \big)(i) = 1$ if and only if $\big(\apunch^{(n)} \big)(i+1) = 1$. By the inductive hypothesis, we have that $\big(\apunch^{(n)} \big)(i+1) = 1$ if and only if
    \[i+1 \in \bigcap_{m=0}^{n} \Big( \big( A - (n-m) \big) \cup \big( \N_0 \setminus \big( W(m)-n \big) \big) \Big).\]
    The previous line rearranges to
    \begin{align}
    \label{eqn_big_set_description_for_inductive_step}
        i \in \bigcap_{m=0}^{n} \Big( \big( A - (n+1-m) \big) \cup \big( \N_0 \setminus \big( W(m)-(n+1) \big) \big) \Big),
    \end{align}
    which shows \eqref{item_explicit_description_of_omega_and_kappa_one} for $n+1$.

    To see \eqref{item_explicit_description_of_omega_and_kappa_two} for $n+1$, let $i \in \N_0$. Note that by the definition of the map $\pi_{\nu_2(n+1)}$, $\big(\apunch^{(n+1)} \big)(i) = \big(\pi_{\nu_2(n+1)} \bpunch^{(n+1)} \big)(i) = 1$ if and only if $\big(\bpunch^{(n+1)}\big)(i) = 1$ and either
    \begin{itemize}
        \item $i \not\in I_{\nu_2(n+1)}(\bpunch^{(n+1)}) = W(n+1) - (n+1)$, so that $\pi_{\nu_2(n+1)}$ does not modify $\bpunch^{(n+1)}$ at $i$; or
        \item $i \in A$, so that, if $\pi_{\nu_2(n+1)}$ modifies $\bpunch^{(n+1)}$ at $i$, it does not change 1 to 0.
    \end{itemize}
    By the previous paragraph, we have that $\big(\bpunch^{(n+1)}\big)(i) = 1$ if and only if \eqref{eqn_big_set_description_for_inductive_step} holds, and so we see by the previous sentence that $\big(\apunch^{(n+1)} \big)(i) = 1$ if and only if $i$ belongs to the set
    \[\bigcap_{m=0}^{n} \Big( \big( A - (n+1-m) \big) \cup \big( \N_0 \setminus \big( W(m)-(n+1) \big) \big) \Big) \cap \big(A \cup \big(\N_0 \setminus (W(n+1)-(n+1))\big) \big).\]
    The set on the previous line simplifies to the one in \eqref{eqn_set_indicated_by_kappa_n} for $n+1$, demonstrating \eqref{item_explicit_description_of_omega_and_kappa_two} for $n+1$, as desired.
\end{proof}



\begin{lemma}
\label{lemma_second_coords_agree}
    Let $m \in \N_0$.  For all $i \in W(m) - m$, the points $\apunch^{(i)}$ and $\apunch^{(m+i)}$ agree on the interval $W(m) - m - i$.
\end{lemma}

\begin{proof}
The conclusion is trivial for $m = 0$, so suppose $m \in \N$. Let $\ell = |W(m)|$. We must show that for all $i \in [\ell]$, the points $\apunch^{(i)}$ and and $\apunch^{(m+i)}$ agree on $W(m) - m - i = [\ell - i]$.  We will prove this by induction on $i$.  The base case $i = 0$ follows from the definition of $W(m)$: the points $\apunch^{(0)} = 1_A$ and $\apunch^{(m)}$ agree on $W(m) - m = [\ell]$.

Let $i \in [\ell - 1]$, and suppose that $\apunch^{(i)}$ and $\apunch^{(m+i)}$ agree on $[\ell - i]$.  We will show that $\apunch^{(i+1)}$ and $\apunch^{(m+i+1)}$ agree on $[\ell - i - 1]$.  Note that $\apunch^{(i+1)} = \pi_{\nu_2(i+1)} \sigma \apunch^{(i)}$ and $\apunch^{(m+i+1)} = \pi_{\nu_2(m+i+1)} \sigma \apunch^{(m+i)}$.



Since $\apunch^{(i)}$ and $\apunch^{(m+i)}$ agree on $[\ell - i]$, we have that $\sigma \apunch^{(i)}$ and $\sigma \apunch^{(m+i)}$ agree on $[\ell - i - 1]$.  If we show that $\nu_2(i+1) = \nu_2(m+i+1)$ and that $2^{\nu_2(m+i+1)} \leq \ell - i - 1$, then it will follow from \cref{lemma_continuity_of_pins} that $\pi_{\nu_2(i+1)} \sigma \apunch^{(i)}$ and $\pi_{\nu_2(m+i+1)} \sigma \apunch^{(m+i)}$ agree on $[\ell - i - 1]$, concluding the proof of the inductive step.

Since $W(m)$ is a dyadic interval, $m+i+1 \in (W(m) \setminus \{m\}) \cap 2^{\nu_2(m+i+1)} \N$, and $|W(m)| \leq 2^{\nu_2(m)}$, we get from \cref{lemma_basic_dyadic_cubes_facts} that $\nu_2(m+i+1) < \nu_2(m)$.  By properties of the 2-adic valuation, it follows that $\nu_2(m+i+1) = \nu_2(i+1)$. To see that $2^{\nu_2(m+i+1)} \leq \ell - i - 1$, note that $m+i+1+[2^{\nu_2(m+i+1)}]$ is a dyadic interval, intersecting, hence contained in, $W(m)$.  Therefore, $m+i+1+2^{\nu_2(m+i+1)} \leq m + |W(m)| = m + \ell$, which rearranges to the desired inequality, finishing the proof of the inductive step.
\end{proof}

The following lemma shows the ``fractal'' structure of containment amongst the punch windows.

\begin{lemma}
\label{lemma_fractal_windows}
    For all $n, m \in \N_0$, if $m < n$ and $n \in W(m)$, then $W(n)-m = W(n-m)$.
\end{lemma}

\begin{proof}
Let $n, m \in \N_0$ with $m < n$ and $n \in W(m)$.  If $m = 0$, then conclusion is immediate, so suppose $m \geq 1$.  Define $\ell = |W(m)|$.

We claim that $\nu_2(n) = \nu_2(n-m)$ and that $2^{\nu_2(n)} \leq \ell - (n-m)$.  Indeed, since $W(m)$ is a dyadic interval, $n \in (W(m) \setminus\{m\}) \cap 2^{\nu_2(n)} \N$, and $2^{\nu_2(m)} \geq |W(m)|$, we see from \cref{lemma_basic_dyadic_cubes_facts} that $\nu_2(n) < \nu_2(m)$.  By properties of the 2-adic valuation, this implies that $\nu_2(n) = \nu_2(n-m)$.  Since $n + [2^{\nu_2(n)}]$ is a dyadic interval intersecting, and hence contained in, $W(m)$, we see that $n + 2^{\nu_2(n)} \leq m + |W(m)| = m + \ell$.  Rearranging, we have that $2^{\nu_2(n)} \leq \ell - (n-m)$.



Next, we claim that $\bpunch^{(n)}$ and $\bpunch^{(n-m)}$ agree on $[\ell - n+m]$.  Indeed, it follows from \cref{lemma_second_coords_agree} with $i = n-m-1 \in [\ell]$ that the points $\apunch^{(n-m-1)}$ and $\apunch^{(n-1)}$ agree on $[\ell - (n-m-1)]$.  Applying the map $\sigma$, we see that $\bpunch^{(n)} = \sigma \apunch^{(n-1)}$ and $\bpunch^{(n-m)} = \sigma \apunch^{(n-m-1)}$ agree on $[\ell - n+m]$.



To reach the conclusion of the lemma, since $W(n)-m$ and $W(n-m)$ are both (possibly empty) intervals starting at $n-m$, it suffices to show that $|W(n)| = |W(n-m)|$. Recall that $|W(n)|$ is the length of the interval on which $\pi_{\nu_2(n)}$ changes $\bpunch^{(n)}$, and similarly for $|W(n-m)|$ and $\bpunch^{(n-m)}$.  Since $\bpunch^{(n)}$ and $\bpunch^{(n-m)}$ agree on $[\ell - n+m]$, and since $\nu_2(n) = \nu_2(n-m)$ and $2^{\nu_2(n)} \leq \ell - (n-m)$, it follows from \cref{lemma_continuity_of_pins} that $|W(n)| = |W(n-m)|$, as desired.
\end{proof}



\begin{lemma}
\label{lemma_punches_at_new_elements}
    Let $n, \ell \in \N_0$.  If
    \begin{enumerate}
        \item \label{item_n_is_new_position}
        the greatest integer $m \in [n]$ such that $n \in W(m)$ is $0$, and

        \item \label{item_original_A_is_punchable_at_n}
        $n + \big(A \cap [2^{\ell}]\big) \subseteq A$,
    \end{enumerate}
    then $|W(n)| \geq 2^{\min(\ell,\nu_2(n))}$.
\end{lemma}

\begin{proof}
    If $n = 0$, then $W(0) = \N$ and the conclusion holds.  Suppose $n \geq 1$.  It follows from \cref{lemma_explicit_description_of_omega_and_kappa} that $\bpunch^{(n)}$ is equal to $1_{A-n}$ on the set $\N_0 \setminus \bigcup_{m=1}^{n-1} (W(m) - n)$. Indeed, if $i \in \N_0 \setminus \bigcup_{m=1}^{n-1} (W(m) - n) = \bigcap_{m=1}^{n-1} \big( \N_0 \setminus (W(m) - n) \big)$, then $i$ belongs to the set in \eqref{eqn_set_indicated_by_omega_n} if and only if $i \in A-n$. If \eqref{item_n_is_new_position} holds, then the set $\bigcup_{m=1}^{n-1} (W(m) - n)$ contains no non-negative integers.  Therefore, the points $\bpunch^{(n)}$ and $1_{A-n}$ are equal on the set $\N_0$.

    Suppose \eqref{item_n_is_new_position} and \eqref{item_original_A_is_punchable_at_n} hold. Since $\bpunch^{(n)} = 1_{A-n}$ and $n + \big(A \cap [2^{\ell}]\big) \subseteq A$, we see that $\nu_A(\bpunch^{(n)}) \geq \ell$.  It follows from the definition of $W(n)$ that
    \[W(n) = n + [2^{\min(\nu_A(\bpunch^{(n)}),\nu_2(n))}] \supseteq n + [2^{\min(\ell,\nu_2(n))}],\]
    whereby $|W(n)| \geq 2^{\min(\ell,\nu_2(n))}$, as desired.
\end{proof}


For $\ell \in \N_0$, we define
\[L(\ell) \defeq \big\{n \in \N \ \big| \ |W(n)| \geq 2^\ell \big\}\]
to be the set of those positive integers $n$ at which the punch window $W(n)$ is of length at least $2^\ell$.  The following lemma shows that the collection of return time sets of the point $(0,1_A)$ to neighborhoods of itself under the shift-punch map is essentially the same as the collection of sets $L(\ell)$, $\ell \in \N_0$.


\begin{lemma} \leavevmode
\label{lemma_punch_windows_and_return_times}
    \begin{enumerate}
        \item \label{item_Lell_contained_in_return_times}
        For all $\eps > 0$, there exists $\ell \in \N_0$ such that
        \[L(\ell) \subseteq R_{\kappa}\big((0,1_A),B_{\eps}((0,1_A))\big).\]

        \item \label{item_return_times_contained_in_Lell}
        For all $\ell \in \N_0$, there exists $\eps > 0$ such that
        \[R_{\kappa}\big((0,1_A),B_{\eps}((0,1_A))\big) \subseteq L(\ell).\]
    \end{enumerate}
\end{lemma}

\begin{proof}
    \eqref{item_Lell_contained_in_return_times} \
    Let $\eps > 0$.  Choose $\ell \in \N_0$ sufficiently large so that for all $n \in \N$ and all $\omega \in X$, if $\nu_2(n) \geq \ell$ and $\omega$ and $1_A$ agree on $[2^\ell]$, then $d_{\odometer \times X}\big( (n,\omega), (0, 1_A) \big) < \eps$.

    Let $n \in L(\ell)$.  We see from the definition of $W(n)$ that $\nu_2(n) \geq |W(n)| \geq \ell$ and that the points $\apunch^{(n)}$ and $1_A$ agree on $W(n) - n \supseteq [2^\ell]$.  By the choice of $\ell$, we have that $d\big( \kappa^n (0,1_A), (0, 1_A)\big) < \eps$, whereby $n \in R_{\kappa}\big((0,1_A),B_{\eps}((0,1_A))\big)$, as was to be shown.

    \eqref{item_return_times_contained_in_Lell} \
    Let $\ell \in \N_0$.  Choose $\eps > 0$ so that for all pairs of points $(n,\omega) \in \odometer \times X$, if $(n,\omega)$ and $(0,1_A)$ are at a distance of at most $\eps$, then $\nu_2(n) \geq \ell$ and $\omega$ and $1_A$ agree on $[2^{\ell}]$.

    Let $n \in R_{\kappa}\big((0,1_A),B_{\eps}((0,1_A))\big)$.  Since $\kappa^n(0,1_A) = (n, \apunch^{(n)})$ and $(0,1_A)$ are at a distance of at most $\eps$, we have that $\nu_2(n) \geq \ell$ and $\apunch^{(n)}$ and $1_A$ agree on $[2^{\ell}]$. By the definition of $\pi_{\nu_2(n)}$, we have that $\apunch^{(n)} \leq \bpunch^{(n)}$ pointwise.  Since $\apunch^{(n)}$ and $1_A$ agree on $[2^{\ell}]$, we have that $A \cap [2^\ell] \subseteq \supp(\bpunch^{(n)})$.  Thus, $\nu_A(\bpunch^{(n)}) \geq \ell$. It follows that $|W(n)| = 2^{\min(\nu_A(\bpunch^{(n)}),\nu_2(n))} \geq 2^{\ell}$, whereby $n \in L(\ell)$, as was to be shown.
\end{proof}

We are ultimately interested in the result of the punches made on the set $A$.  Thus, we define
\begin{align}
\label{eqn_result_of_set_after_punches}
    B \defeq \bigcap_{n=0}^\infty \Big( \big( A + n \big) \cup \big(\N_0 \setminus W(n) \big) \Big).
\end{align}
Informally, the set $B$ is the set $A$ after all punches have occurred; this is made precise by the equivalent dynamical characterization of the set $B$ given in \cref{lemma_dynamical_description_of_B}. Since $W(0) = \N_0$, we see from the definition that $B \subseteq A$.  Moreover, we see that if $0 \in A$, then $0 \in B$.

\begin{lemma}
\label{lemma_dynamical_description_of_B}
    For all $n \in \N_0$, we have that $\big(\apunch^{(n)}\big)(0) = 1_B(n)$.
\end{lemma}

\begin{proof}
    Let $n \in \N_0$.  It follows by \cref{lemma_explicit_description_of_omega_and_kappa} that $\big(\apunch^{(n)}\big)(0) = 1$ if and only if
    \[0 \in \bigcap_{m=0}^{n} \Big( \big( A - (n-m) \big) \cup \big( \N_0 \setminus \big( W(m)-n \big) \big) \Big).\]
    By set algebra, the previous line holds if and only if
    \[n \in \bigcap_{m=0}^{n} \Big( \big( A +m \big) \cup \big( \N_0 \setminus W(m) \big) \Big).\]
    Since for all $m \in \N_0$ greater than $n$, $W(m) \subseteq \{n+1, n+2, \ldots \}$, we see that the previous line holds if and only if
    \[n \in \bigcap_{m=0}^{\infty} \Big( \big( A +m \big) \cup \big( \N_0 \setminus W(m) \big) \Big) = B,\]
    as was to be shown.
\end{proof}

The following lemma connects the times of return of the point $1_B$ to neighborhoods of itself under the usual shift to the family of sets $L(\ell)$ defined above.

\begin{lemma}
\label{lemma_punch_returns_vs_shift_returns}
    For all $\eps > 0$, there exists $\ell \in \N_0$ such that
    \[L(\ell) \subseteq R_\sigma\big(1_B, B_\eps\big(1_B \big) \big).\]
\end{lemma}

\begin{proof}
Let $\eps > 0$.  Choose $\ell \in \N_0$ so that if two points in $X$ agree on $[2^\ell]$, then they are at a distance of no more than $\eps$.  

Let $n \in L(\ell)$.  By \cref{lemma_second_coords_agree}, for all $i \in [2^\ell] \subseteq W(n) - n$, the points $\apunch^{(i)}$ and $\apunch^{(n+i)}$ agree on $W(n) - n - i$.  In particular, they agree at 0.  It follows from \cref{lemma_dynamical_description_of_B} that for all $i \in [2^\ell]$, $1_B(i) = 1_B(i+n) = 1_{B-n}(i) = (\sigma^{n} 1_B)(i)$.  Since $1_B$ and $\sigma^{n} 1_B$ agree on $[2^\ell]$, we see that $n \in R_\sigma(1_B,B_\eps(1_B))$, as was to be shown.
\end{proof}

\begin{remark}
    If $0 \not\in A$, then $B = \emptyset$.  Indeed, since $0 \not\in A$, for all $n \in \N_0$, the map $\pi_{\nu_2(n)}$ will punch $\bpunch^{(n)}$ at least at the zero coordinate (that is, $n \in W(n)$), changing that coordinate to 0 if it is 1.  Thus, $\apunch^{(n)}(0) = 0$, and according to \cref{lemma_dynamical_description_of_B}, the set $B$ is empty. In this case, we see that $R_\sigma\big(1_B, B_\eps\big(1_B \big) \big) = \N$.  This shows that a result analogous to \cref{lemma_punch_windows_and_return_times} for the point $1_B$ under the shift $\sigma$ does not necessarily hold: it may be that there exists $\ell \in \N_0$ such that for all $\eps > 0$, the set $R_\sigma\big(1_B, B_\eps(1_B ) \big)$ is not contained in $L(\ell)$.
\end{remark}









\subsubsection{Central syndetic sets are dynamically central syndetic}
\label{sec_proof_that_member_of_synd_idem_filter_is_dcsyndetic}

To finish the proof of \cref{thm_dcs_iff_cs}, we must show that central syndetic sets are dynamically central syndetic.  It suffices by \cref{thm_conditions_for_central_syndetic} and \cref{lemma_upgraded_cssd_property} to show: if $A \subseteq \N$ satisfies the condition in \eqref{eqn_stronger_than_cssd}, then $A$ is dynamically central syndetic. We carry forward all of the notation from the previous sections.

Let $A \subseteq \N$ satisfy the condition in \eqref{eqn_stronger_than_cssd}.  We append $0$ to $A$ and, by a slight abuse of notation, consider $A$ as a subset of $\N_0$.  The key step is to show that the condition in \eqref{eqn_stronger_than_cssd} implies that the sets $L(\ell)$ defined in the previous section are syndetic.


\begin{theorem}
\label{prop_long_punch_times_are_syndetic}
    For all $\ell \in \N_0$, the set $L(\ell)$ is syndetic.
\end{theorem}

\begin{proof}
    Let $\ell \in \N_0$. We will show first that the set $L(\ell)$ is nonempty.  It follows from \eqref{eqn_stronger_than_cssd} that there exists $n \in 2^\ell \N$ such that $n + \big(A \cap [2^\ell]\big) \subseteq A$.  Let $m \in [n]$ be the greatest integer such that $n \in W(m)$.  If $m=0$, then by \cref{lemma_punches_at_new_elements}, we have that $|W(n)| \geq 2^\ell$, whereby $n \in L(\ell)$.  Otherwise, we have that $1 \leq m < n$ and $n \in W(m)$, whereby $\big( W(m) \setminus \{m\} \big) \cap 2^\ell \N \neq \emptyset$.  It follows from \cref{lemma_basic_dyadic_cubes_facts} that $|W(m)| \geq 2^{\ell + 1}$, so $m \in L(\ell)$.

    By \eqref{eqn_stronger_than_cssd}, \cref{lemma_syndetic_characterization}, and the previous paragraph, there exists $k \in \N$ so that
    \begin{align}
    \label{eqn_cube_syndeticity_from_upgraded_membership_property}
        2^\ell \N_0 \cap \bigcap_{a \in A \cap [2^{\ell}]} (A-a) \in \cubes_{2^{k}}^*
    \end{align}
    and so that $L(\ell) \cap [2^k] \neq \emptyset$.  We will show that $L(\ell)$ is syndetic by showing that $L(\ell) \in \cubes_{2^{k+1}}^*$ and appealing again to \cref{lemma_syndetic_characterization}.

    To show that $L(\ell) \in \cubes_{2^{k+1}}^*$, we will show that for all $a \in 2^{k+1}\N_0$, the set $L(\ell) \cap \big(a + ([2^{k+1}]\setminus \{0\})\big)$ is nonempty.  We proceed by induction on $a$.  We have already shown the base case $a=0$. (Recall that, by definition, $L(\ell) \subseteq \N$, so that $L(\ell) \cap [2^k] \neq \emptyset$ implies $L(\ell) \cap ([2^k] \setminus \{0\}) \neq \emptyset$.)

    Let $a \in 2^{k+1}\N$, and define $Q = a + [2^{k+1}]$.  From \eqref{eqn_cube_syndeticity_from_upgraded_membership_property}, there exists
    \[n \in \big( Q \setminus \{a\} \big) \cap 2^\ell \N_0 \cap \bigcap_{a \in A \cap [2^{\ell}]} (A-a).\]
    Let $m \in [n]$ be the greatest integer such that $n \in W(m)$.  If $m=0$, then by \cref{lemma_punches_at_new_elements}, we have that $|W(n)| \geq 2^\ell$, whereby $n \in L(\ell) \cap \big( Q \setminus \{a\} \big)$, as desired. Otherwise, we have that $1 \leq m < n$.

    If $m > a$, then $m \in Q \setminus \{a\}$.  Since $m < n$ and $n \in W(m)$, we have that $\big( W(m) \setminus \{m\} \big) \cap 2^\ell \N \neq \emptyset$.  It follows from \cref{lemma_basic_dyadic_cubes_facts} that $|W(m)| \geq 2^{\ell + 1}$, so $m \in L(\ell) \cap \big( Q \setminus \{a\} \big)$, as desired.

    If, on the other hand, $m \leq a$, then the fact that $W(m)$ and $Q$ are intersecting dyadic intervals (both contain $n$) implies that $m \in 2^{k+1} \N$, that $Q \subseteq W(m)$, and that $Q - m \in \cubes_{2^{k+1}}$. It follows by the induction hypothesis that there exists $q \in Q \setminus \{a\}$ such that $q-m \in L(\ell) \cap (Q-m)$.  Thus, $|W(q-m)| \geq 2^\ell$. Since $q \in W(m)$ and $m \leq a < q$, it follows from \cref{lemma_fractal_windows} that $W(q-m) = W(q) - m$.  Therefore, $|W(q)| \geq 2^\ell$, whereby $q \in L(\ell) \cap (Q \setminus \{a\})$, as desired.
\end{proof}



Let $B \subseteq A$ be the subset of $A$ defined in \eqref{eqn_result_of_set_after_punches} in the previous section.  Since $0 \in A$, we see that $0 \in B$. It follows by combining \cref{lemma_punch_returns_vs_shift_returns} and \cref{prop_long_punch_times_are_syndetic} that the point $1_B$ is uniformly recurrent in the shift $(X,\sigma)$.  Since $0 \in B$, it follows from \cref{thm:simple_equivalent_dcS} that the set $B \setminus \{0\} \subseteq \N$ is dynamically central syndetic.  Therefore, we have that the set $A \setminus \{0\}$ is dynamically central syndetic, concluding the proof of \cref{thm_dcs_iff_cs}.

While we do not have need for it, note that it follows by combining \cref{lemma_punch_windows_and_return_times} and \cref{prop_long_punch_times_are_syndetic} that the point $(0,1_A)$ is uniformly recurrent in the shift-punch system $(\odometer \times X, \kappa)$.



\subsection{Combinatorial characterizations: proofs of Theorems \ref{mainthm_ds_equivalents} and \ref{intro_mainthm_full_chars_of_dcsyndetic}}
\label{sec_char_of_ds_and_dcs}

In this section, we bring everything together to prove Theorems \ref{mainthm_ds_equivalents} and \ref{intro_mainthm_full_chars_of_dcsyndetic} from the introduction and give a short application.  Logically, it makes more sense to prove \cref{intro_mainthm_full_chars_of_dcsyndetic} first.

\begin{proof}[Proof of \cref{intro_mainthm_full_chars_of_dcsyndetic}]
    This follows immediately by combining \cref{thm_conditions_for_central_syndetic,thm_dcs_iff_cs}.
\end{proof}

\begin{proof}[Proof of \cref{mainthm_ds_equivalents}]
    \eqref{item_intro_ds_def_condition} $\iff$ \eqref{item_intro_translate_belongs_to_sif} \ If follows by \cref{lemma_translates_of_dsyndetic_sets} that $A$ is dynamically syndetic if and only if there exists $n \in \N_0$ such that $A-n$ is dynamically central syndetic.  By \cref{intro_mainthm_full_chars_of_dcsyndetic}, a set is dynamically central syndetic if and only if it belongs to a syndetic, idempotent filter.  The desired equivalence is shown by combining these two results.

    \eqref{item_intro_translate_belongs_to_sif} $\implies$ \eqref{item_intro_ds_combo_condition} \ Suppose that there exists $n \in \N_0$ such that $A' \defeq A - n$ belongs to a syndetic, idempotent filter.  We have by \cref{intro_mainthm_full_chars_of_dcsyndetic} that there exists $B' \subseteq A'$ such that for all finite $F' \subseteq B'$, the set $B' \cap \bigcap_{f' \in F'} (B' - f')$ is syndetic.  Put $B = B' + n$, and note that $B$ is a nonempty subset of $A$.  Let $F \subseteq B$ be finite, and note that $\min F \geq n+1$.  We have that $F-n \subseteq B'$ is finite, and hence that the set
    \[\bigcap_{f' \in F-n} (B' - f')\]
    is syndetic.  But this set is contained in $\bigcap_{f \in F} (B - f)$, demonstrating syndeticity, as desired.

    \eqref{item_intro_ds_combo_condition} $\implies$ \eqref{item_intro_translate_belongs_to_sif} \ Let $b \in B$.  We will show that the set $B - b \subseteq A-b$ satisfies: for all finite $F \subseteq B-b$, the set
    \[(B-b) \cap \bigcap_{f \in F} (B - b - f)\]
    is syndetic.  It will follow by \cref{thm_conditions_for_central_syndetic} that the set $A-b$ is contained in a syndetic, idempotent filter.

    Let $F \subseteq B-b$ be finite.  By assumption, since $\{b\} \cup (F + b) \subseteq B$ is finite, we have that the set 
    \[\bigcap_{f \in \{b\} \cup (F + b)} (B - f) = (B-b) \cap \bigcap_{f \in F} (B - b - f)\]
    is syndetic, as was to be shown.
\end{proof}

As a short application of \cref{intro_mainthm_full_chars_of_dcsyndetic}, we will show that thickly syndetic sets are dynamically central syndetic.  This result has appeared in related forms several times in the literature (see \cite[Thm. 1]{Glasner-Weiss-Interpolation}, \cite[Thm. 2.4]{huang_ye_2005}, \cite[Prop. 4.4]{dong_shao_ye_2012}).

\begin{theorem}
\label{ex_ps_star_is_dcs}
    Every thickly syndetic set is dynamically central syndetic.
\end{theorem}

\begin{proof}
By \cref{lemma_condition_on_subfamily_of_ps_star}, the family $\PS^*$ is a syndetic, translation-invariant (hence, idempotent) filter.  It follows from \cref{intro_mainthm_full_chars_of_dcsyndetic} that every member of $\PS^*$ is dynamically central syndetic, as desired.
\end{proof}

Note that we upgrade \cref{ex_ps_star_is_dcs} in \cref{cor:largest_family_satisfies_cap} to: the intersection of a thickly syndetic set and a dynamically central syndetic set is dynamically central syndetic.




\section{Part II: Dynamically thick sets and sets of pointwise recurrence}
\label{sec_dthick}

In this section, we define the families of dynamically thick sets and sets of pointwise recurrence.  The main results are the combinatorial characterizations of these families (Theorems \ref{mainthm_dct_combo_characterizations} and \ref{mainthm_dt_characterizations}) in \cref{sec_dthick_comb_characterizations}; a  structure result for dynamically thick sets (\cref{mainthm_structure_for_dy_thick_sets}) in \cref{sec_general_form_of_dthick}; and partial progress in \cref{sec:sigma-compact} on the problem of partitioning dynamically thick sets.

\subsection{Definition and first characterizations}
\label{sec_dy_thick_def_and_char}

A set $A \subseteq \N$ is \emph{dynamically thick} if for all minimal systems $(X,T)$ and all points $x \in X$, the set $\{T^n x \ | \ n \in A\}$ is dense in $X$.  The set $A$ is \emph{a set of (minimal topological) pointwise recurrence} -- also, \emph{dynamically central thick} -- if for all minimal systems $(X,T)$, all points $x \in X$, and all open $U \subseteq X$ containing $x$, there exists $a \in A$ such that $T^a x \in U$. The families of dynamically thick sets and sets of pointwise recurrence are denoted $\dthick$ and $\dcthick$, respectively.

\begin{lemma}
    \label{lemma_dt_and_ds_duals}
    The families of dynamically (central) syndetic and dynamically (central) thick sets are dual, that is,
    \[
        \dthick = \dsyndetic^* \qquad \text{ and } \qquad \dcthick = \dcsyndetic^*.
    \]
\end{lemma}

\begin{proof}
    This follows immediately from the definitions. 
\end{proof}

It can be verified quickly from the definitions or from \cref{lemma_dt_and_ds_duals} that
\[\thick \subseteq \dthick \subseteq \dcthick\]
and that each of these inclusions is proper.  Here is a short list of examples of dynamically (central) thick sets that may be helpful to have in mind; more elaborate examples are given in \cref{sec_dy_thick_examples}.
\begin{enumerate}
    \item For $k \in \N$, the set of positive multiples of $k$, $k\N$, is a set of pointwise recurrence. In fact, for all thick sets $H \subseteq \N$, the set $k \N \cap H$ is a set of pointwise recurrence. Indeed, this follows from \cref{lemma_dilates_of_dsyndetic_sets} and the fact that $H / k$ is thick: if $A \in \dcsyndetic$, then $A / k \in \dcsyndetic$, whereby $A / k \cap H / k \neq \emptyset$, whereby $A \cap k\N \cap H \neq \emptyset$.
    

    \item If $A$ is not dynamically (central) syndetic, then its complement, $\N \setminus A$, is dynamically (central) thick.  This follows immediately from \cref{lemma_dt_and_ds_duals}.  The set in \eqref{eqn_intro_even_odd} from the introduction is not dynamically syndetic, hence its complement, a set of the same form, is dynamically thick.  A more general class of dynamically thick sets is described in \cref{ex:type2-special}.

    \item \label{ex_transitive_point_visit_times} If $x \in X \defeq \{0,1\}^{\N}$ has a dense orbit in the full shift $(X,\sigma)$, then for all nonempty, open $U \subseteq X$, the set $R_\sigma(x,U)$ is dynamically thick. This is a consequence of a theorem of Furstenberg \cite[Thm. II.2]{Furstenberg-disjointness}: the full shift is \emph{disjoint} (cf. \cite[Def. II.1]{Furstenberg-disjointness}) from every minimal system.  (This result is generalized in \cite[Thm. 6.2]{glasner_tsankov_weiss_zucker_2021} to arbitrary countable, infinite group actions.) Indeed, we will argue that for all minimal systems $(Y,S)$, all $y \in Y$, and all nonempty, open $V \subseteq Y$, the set
    \begin{align}
        \label{eqn_disjointness_leads_to_non_empty_set_returns_in_prod}
        R_{\sigma \times S} ((x,y), U \times V) = 
        R_\sigma(x,U) \cap R_S(y,V) \neq \emptyset.
    \end{align}
    By an equivalent definition of the disjointness of $(X,\sigma)$ and $(Y,S)$ (cf. \cite[Lemma II.1]{Furstenberg-disjointness}), the system $(X \times Y, \sigma \times S)$ is the only subsystem of the product system that projects to $X$ in the first coordinate and $Y$ in the second.  Because $x$ and $y$ have dense orbits in $X$ and $Y$, respectively, the subsystem $(\overline{(\sigma \times S)^{\N}(x,y)}, \sigma \times S)$ projects to $X$ and $Y$ in the first and second coordinates, respectively, and hence is the product system.  Since the point $(x,y)$ has a dense orbit in the product system, we see that \eqref{eqn_disjointness_leads_to_non_empty_set_returns_in_prod} holds.  We offer an alternative explanation of this example in \cref{rem:combinatorial_primes} and \cref{rmk_why_transitive_works} below.
\end{enumerate}
We call the reader's attention to the fact that dynamically central thick sets (sets of pointwise recurrence) need not be central, despite the terminology.  The example given in \cref{rem:combinatorial_primes} demonstrates this.


The robustness of the families of dynamically (central) syndetic sets was discussed in \cref{sec_dy_synd_def_and_char}.  Because of the duality shown in \cref{lemma_dt_and_ds_duals}, the families of dynamically (central) thick sets exhibit a similar robustness under changes to the combinatorial, topological, or dynamical requirements in their definitions.  Thus, the families of dynamically (central) thick subsets of $\N$ remain unchanged if \dots
\begin{enumerate}
    \item \dots \ ``for all minimal systems'' is changed to ``for all invertible minimal systems,'' ``for all zero dimensional minimal systems ,'' or ``for all (not-necessarily-metrizable) minimal systems.'' By \cref{lemma_dt_and_ds_duals}, these follow immediately from \cref{lemma_may_assume_invertible} and \cref{thm:simple_equivalent_dcS,lemma_metrizability_doesnt_matter}, respectively.

    \item \dots \ (for sets of pointwise recurrence) ``all open $U \subseteq X$ with $x \in U$'' is changed to ``all open $U \subseteq X$ with $x \in \overline{U}$.'' This is a consequence of \cref{lemma_dt_and_ds_duals} and \cref{thm:simple_equivalent_dcS}.

    \item \dots \ sets are considered up to equivalence on thickly syndetic ($\PS^*$) sets.  This is a consequence of \cref{cor:largest_family_satisfies_cap}: a set $A \subseteq \N$ is dynamically (central) thick if and only if for all thickly syndetic $B \subseteq \N$, the set $A \cap B$ is dynamically (central) thick.
\end{enumerate}



\subsubsection{First results}
\label{sec_dt_first_results}

In this section, we collect some first results concerning dynamical thickness.  As discussed in \cref{sec_intro_set_and_pt_recurrence}, that dynamically (central) thick subsets of $\N$ are piecewise syndetic has been shown several times in the literature.  We record it here for ease of reference.

\begin{theorem}
\label{thm_dy_c_thick_is_ps}
    Sets of pointwise recurrence (and, hence, dynamically thick sets) are piecewise syndetic.
\end{theorem}

\begin{proof}
    This follows immediately from the duality in \cref{lemma_dt_and_ds_duals} and \cref{ex_ps_star_is_dcs}.
\end{proof}


It is easy to see that the family of dynamically thick sets is not partition regular.  Indeed, neither the set of even nor the set of odd positive integers is dynamically thick. Host, Kra, and Maass asked in \cite[Question 6.5]{HostKraMaass2016} whether or not the family $\dcthick$ is partition regular.  The following result gives a strong negative answer to this question.


\begin{theorem}
\label{thm_partition_of_N_into_disjoint_non_sets_of_rec}
    There exist sets $A, B \subseteq \N$ that are not sets of pointwise recurrence but for which $\N = A \cup B$.  In particular, the family of sets of pointwise recurrence is not partition regular.
\end{theorem}

\begin{proof}
    Since $\N$ is a dynamically central syndetic set, according to \cref{mainthm_partition_of_dcs_sets}, there exists a partition $\N = A \cup B$ into disjoint, dynamically central syndetic sets. Because $A \cap B = \emptyset$, neither $A$ nor $B$ is a set of pointwise recurrence.
\end{proof}

We note that the partition guaranteed by \cref{mainthm_partition_of_dcs_sets} can be made quite explicit: for $\alpha \in \R \setminus \Q$, define $A \defeq \big\{n \in \N \ \big| \ \{n \alpha\} \in (0,1/2)\big\}$ and $B \defeq \big\{n \in \N \ \big| \ \{n \alpha\} \in (1/2,1)\big\}$.  Neither $A$ nor $B$ is a set of pointwise recurrence, but $\N = A \cup B$.


\begin{lemma}
    \label{lemma_intersection_of_ds_and_dt_is_infinite}
    The intersection of a dynamically (central) syndetic set and a dynamically (central) thick set is infinite.
\end{lemma}

\begin{proof}
    Let $A \in \dsyndetic$, $B \in \dthick$, and $n \in \N$.  By \cref{lemma_ds_dcs_modifyable_on_a_finite_set}, we see that $A \cap \{n+1, n+2, \ldots\} \in \dsyndetic$.  Since $B \in \dthick$, we have that $A \cap \{n+1, n+2, \ldots\} \cap B \neq \emptyset$.  Therefore, $\max (A \cap B) > n$.  Since $n \in \N$ was arbitrary, we see that $A \cap B$ is infinite.  The same argument works to show that $A \cap B$ is infinite when $A \in \dcsyndetic$ and $B \in \dcthick$.
\end{proof}

\cref{thm_dcps_and_dps_are_ps} in the next section simultaneously strengthens both \cref{thm_dy_c_thick_is_ps} and \cref{lemma_intersection_of_ds_and_dt_is_infinite} by showing that the intersection of a dynamically (central) syndetic set and a dynamically (central) thick set is, in fact, piecewise syndetic.

\subsubsection{Translates and dilates of dynamically thick sets}

We show in this section what happens to dynamically (central) thick sets under translations and dilations.  Most of the statements in \cref{lemma_translates_of_dthick_sets} can be shown directly from the definitions, but we choose to derive them from \cref{lemma_translates_of_dsyndetic_sets} using the duality in \cref{lemma_dt_and_ds_duals}.

\begin{lemma}
\label{lemma_translates_of_dthick_sets}
    Let $A \subseteq \N$.
    \begin{enumerate}
        \item
        \label{item_all_translates_are_dt}
        If $A$ is dynamically thick, then for all $n \in \N$, the set $A-n$ is dynamically thick.
        
        \item
        \label{item_some_translates_are_dct}
        If the set
        \begin{align}
            \label{eqn_translates_of_dt_set}
            \big\{ n \in \N \ \big| \ A-n \text{ is a set of pointwise recurrence} \big\}
        \end{align}
        is dynamically thick, then $A$ is dynamically thick.
        
        \item
        \label{item_dcs_translates_of_dct_are_dct}
        If the set in \eqref{eqn_translates_of_dt_set} is a set of pointwise recurrence, then $A$ is a set of pointwise recurrence.
    \end{enumerate}
    Moreover, \eqref{item_all_translates_are_dt}, \eqref{item_some_translates_are_dct}, and \eqref{item_dcs_translates_of_dct_are_dct} hold with $A-n$ replaced by $A+n$. Finally,
    \begin{align}
        \label{eqn_relationship_between_dt_and_dct_classes}
        \dthick = \bigcap_{n \in \N} \big(\dcthick - n \big) = \bigcap_{n \in \N} \big(\dcthick + n \big).
    \end{align}
\end{lemma}

\begin{proof}
    \eqref{item_all_translates_are_dt} \ Suppose $A \in \dthick$ and $n \in \N$. To see that $A - n \in \dthick$, by \cref{lemma_dt_and_ds_duals}, it suffices to show that $A-n \in \dsyndetic^*$.  Let $B \in \dsyndetic$.  By \cref{lemma_translates_of_dsyndetic_sets}, the set $B + n \in \dsyndetic$.  By \cref{lemma_intersection_of_ds_and_dt_is_infinite}, the set $A \cap (B+n)$ is infinite, so $(A-n) \cap B \supseteq (A-n) \cap (B+n - n) \neq \emptyset$, as desired.

    \eqref{item_some_translates_are_dct} \ In the language of the family algebra developed in \cref{sec_combinatorics}, this statement is equivalent to $\dthick + \dcthick \subseteq \dthick$.  Recall from \cref{lemma_translates_of_dsyndetic_sets} that $\dsyndetic \subseteq \dsyndetic + \dcsyndetic$. Taking the dual and simplifying using the algebra in \cref{lemma_fam_dual_sums}, we see that $\dthick + \dcthick \subseteq \dthick$.

    \eqref{item_dcs_translates_of_dct_are_dct} \ This statement is equivalent to $\dcthick + \dcthick \subseteq \dcthick$. Taking the dual of $\dcsyndetic \subseteq \dcsyndetic + \dcsyndetic$ from \cref{lemma_translates_of_dsyndetic_sets}, we see that $\dcthick + \dcthick \subseteq \dcthick$.
    
    That \eqref{item_all_translates_are_dt} holds with $A-n$ replaced with $A+n$ follows the same argument used in \eqref{item_all_translates_are_dt}.  That \eqref{item_some_translates_are_dct} and \eqref{item_dcs_translates_of_dct_are_dct} hold with $A-n$ replaced by $A+n$ requires different reasoning, since the set algebra explanation only treats negative translates.

    Suppose that $A$ is such that the set in \eqref{eqn_translates_of_dt_set} (with $A+n$ instead of $A-n$) is dynamically thick.  We wish to show that $A$ is dynamically thick.  It suffices by \cref{lemma_dt_and_ds_duals,lemma_may_assume_invertible} to show: for all minimal, invertible $(X,T)$, all $x \in X$, and all nonempty, open $U \subseteq X$, there exists $a \in A$ such that $T^a x \in U$.  Thus, let $(X,T)$ be a minimal, invertible system, $x \in X$, and $U \subseteq X$ open and nonempty.  Since the system $(X,T^{-1})$ is minimal (see, for example, \cite[Lemma 2.7]{glasscock_koutsogiannis_richter_2019}) and the set in \eqref{eqn_translates_of_dt_set} is dynamically thick, there exists $n \in \N$ such that $T^{-n} x \in U$ and $A + n \in \dcthick$.  By the definition of set of pointwise recurrence, there exists $m \in A+n$ (so $m > n$) such that $T^m T^{-n} x \in U$.  We see that $m - n \in A$ and $T^{m-n} x \in U$, as desired.  The argument that \eqref{item_dcs_translates_of_dct_are_dct} holds with $A-n$ replaced by $A+n$ is very similar and is left to the interested reader.

    Finally, that \eqref{eqn_relationship_between_dt_and_dct_classes} holds follows immediately from \eqref{item_all_translates_are_dt} and \eqref{item_some_translates_are_dct} for $A-n$ and $A+n$.
\end{proof}

As demonstrated in the proof of \cref{lemma_translates_of_dthick_sets}, it is interesting and useful to formulate the lemma's conclusions in terms of the family algebra developed in \cref{sec_combinatorics}. Thus,
\begin{enumerate}
    \item the family $\dthick$ is translation invariant, that is, $\dthick \subseteq \{\N\} + \dthick$;
    \item $\dthick + \dcthick \subseteq \dthick$;
    \item $\dcthick + \dcthick \subseteq \dcthick$.
\end{enumerate}

\begin{lemma}\label{lem:dilates_of_dT_dcT}
    Let $A \subseteq \N$ and $k \in \N$. If $A$ is a set of pointwise recurrence, then so are the sets $kA$ and $A/k$.  If $A$ is dynamically thick, then so is the set $A/k$.
\end{lemma}

\begin{proof}
    Suppose $A \in \dcthick$.  If $B \in \dcsyndetic$, then by \cref{lemma_dilates_of_dsyndetic_sets}, the sets $kB$ and $B/k$ are $\dcsyndetic$ sets, whereby $A \cap (kB) \neq \emptyset$ and $A \cap (B/k) \neq \emptyset$.  Since $(kB)/k = B$ and $k(B/k) \subseteq B$, it follows that $A/k \cap B \neq \emptyset$ and $kA \cap B \neq \emptyset$. Since $B \in \dcsyndetic$ was arbitrary, we have that $A/k$ and $kA$ are sets of pointwise recurrence.

    Suppose $A \in \dthick$. To show $A/k \in \dthick$ follows the same argument using \cref{lemma_dilates_of_dsyndetic_sets} as in the previous paragraph.
\end{proof}

\subsection{Combinatorial characterizations: proofs of Theorems \ref{mainthm_dct_combo_characterizations} and \ref{mainthm_dt_characterizations}}
\label{sec_dthick_comb_characterizations}

In this section, we give combinatorial characterizations of dynamical thickness and pointwise recurrence.  Recall that the empty union and empty intersection are interpreted to be the empty set and $\N$, respectively.

Theorems \ref{mainthm_dct_combo_characterizations} and \ref{mainthm_dt_characterizations} will follow immediately from Theorems \ref{lemma_our_fav_dy_thick} and \ref{thm_dct_combo_characterizations}, respectively.  In each, the second statement comes from interpreting ``has nonempty intersection with all dynamically (central) syndetic sets'' in combinatorial terms by appealing to the main results on dynamically syndetic sets, while the third, fourth, and fifth statements arise from ultrafilter dynamics, appealing ultimately to \cref{lemma_metrizability_doesnt_matter}.

\begin{theorem}
\label{lemma_our_fav_dy_thick}
    Let $A \subseteq \N$.  The following are equivalent.
    \begin{enumerate}
        \item 
        \label{item_dthick_def}
        The set $A$ is dynamically thick.

        \item 
        \label{item:dyn_thick_idempotent}
        For all $\N \supsetneq B \supseteq A$, there exists a finite set $F \subseteq \N \setminus B$ such that the set $B-F$ is thick.

        \item \label{item_condition_one}
        For all syndetic $S \subseteq \N$, there exists a finite set $F \subseteq A$ such that for all syndetic $S' \subseteq \N$, the set $F \cap (S - S')$ is nonempty.

        \item \label{item:complement_thick_dt_characterization_v2}
        For all piecewise syndetic $S \subseteq \N$, there exists a finite set $F \subseteq A$ such that the set $S-F$ is thick.

        \item
        \label{item:complement_psstar_dt_characterization_star_fixed}
        For all piecewise syndetic $S \subseteq \N$ and all minimal left ideals $L \subseteq \beta \N$ with $\overline{S} \cap L \neq \emptyset$, there exists a finite set $F \subseteq A$ such that $L \subseteq \overline{S-F}$.
    \end{enumerate}
\end{theorem}


\begin{proof}
    \eqref{item_dthick_def} $\iff$ \eqref{item:dyn_thick_idempotent}
    Because $\dsyndetic$ and $\dthick$ are dual, the set $A$ is dynamically thick if and only if the set $A' \defeq \N \setminus A$ is not dynamically syndetic.  By \cref{mainthm_ds_equivalents}, the set $A'$ is not dynamically syndetic if and only if for all nonempty subsets $B' \subseteq A'$, there exists a finite set $F \subseteq B'$ such that the set $\bigcap_{f \in F} (B' - f)$ is not syndetic.  Taking the complement of this set in $\N$, we see that this happens if and only if the set
    \[ \bigcup_{f \in F} \big( (\N \setminus B') - f \big) = \big( \N \setminus B' \big) - F\]
    is thick.  Set $B \defeq \N \setminus B'$. Note that $B' \subseteq A'$ if and only if $A \subseteq B$ and that $B' \neq \emptyset$ if and only if $B \neq \N$.
    
    Summarizing the previous paragraph, we have shown that the set $A$ is dynamically thick if and only if for all $A \subseteq B \subsetneq \N$, there exists a finite set $F \subseteq \N \setminus B$ such that the set $B - F$ is thick, as was to be shown.

    \eqref{item_dthick_def} $\implies$ \eqref{item:complement_psstar_dt_characterization_star_fixed} \ Suppose that $A \in \dthick$.  Let $S \subseteq \N$ be piecewise syndetic and $L \subseteq \beta \N$ be a minimal left ideal such that $\overline{S} \cap L \neq \emptyset$.  Note that $\overline{S} \cap L$ is nonempty and open in $L$. By the definitional robustness of $\dthick$ sets regarding non-metrizable systems (see \cref{sec_dy_thick_def_and_char}), for all $p \in L$, there exists $a \in A$ such that $a + p \in \overline{S}$.  Therefore,
    \[L \subseteq \bigcup_{a \in A} \overline{S - a}.\]
    Because $L$ is compact, there exists a finite set $F \subseteq A$ such that
    \[L \subseteq \bigcup_{f \in F} \overline{S - f}  \subseteq \overline{S-F},\]
    as was to be shown.

    \eqref{item:complement_psstar_dt_characterization_star_fixed} $\implies$ \eqref{item:complement_thick_dt_characterization_v2} \ Let $S \subseteq \N$ be a piecewise syndetic set.  It follows by \cite[Thm. 4.40]{hindman_strauss_book_2012} that there exists a minimal left ideal $L \subseteq \beta \N$ such that $\overline{S} \cap L \neq \emptyset$.  Let $F \subseteq A$ be the finite set guaranteed by \eqref{item:complement_psstar_dt_characterization_star_fixed}.  Since $\overline{S-F}$ contains a minimal left ideal, by \cite[Thm. 4.48]{hindman_strauss_book_2012}, the set $S-F$ is thick, as desired.

    \eqref{item:complement_thick_dt_characterization_v2} $\implies$ \eqref{item_condition_one} \ Let $S \subseteq \N$ be syndetic.  Let $F \subseteq A$ be the finite set guaranteed by \eqref{item:complement_thick_dt_characterization_v2}.  Since $S-F$ is thick, for all $S' \subseteq \N$ syndetic, the set $S' \cap (S-F)$ is nonempty, whereby $F \cap (S-S') \neq \emptyset$, as desired.

    \eqref{item_condition_one} $\implies$ \eqref{item_dthick_def} \ We will prove the contrapositive: if $A$ is not dynamically thick, then there exists a syndetic set $S \subseteq \N$ such that for all finite $F \subseteq A$, the set $S - F$ is not thick. Taking complements and considering $B$ as $\N \setminus A$ and $C$ as $\N \setminus S$, we must show: if a set $B \subseteq \N$ is dynamically syndetic, then there exists a set $C \subseteq \N$ that is not thick such that for all finite $F \subseteq \N \setminus B$, the set $\bigcap_{f \in F} (C - f)$ is syndetic.

    Suppose $B \subseteq \N$ is dynamically syndetic. If $B = \N$, put $C = \emptyset$ and note that the conclusion holds since the empty intersection is $\N$. Suppose $B \neq \N$. By \cref{thm:simple_equivalent_dcS}, there exists a minimal system $(X,T)$, a point $x \in X$, and a nonempty, clopen set $U \subseteq X$ such that $R(x,U) \subseteq B$.  Define $V \defeq X \setminus U$ and $C \defeq \N \setminus R(x,U) = R(x, V)$. Since $U \neq X$ and the set $R(x,U)$ is syndetic, the set $C$ is nonempty but not thick.  Let $F \subseteq \N \setminus B$ be finite.  Since $F \subseteq C$, we have that for all $f \in F$, the set $T^{-f}V$ is an open neighborhood of $x$.  We see that
    \[\bigcap_{f \in F} (C - f) = R\bigg( x, \bigcap_{f \in F}T^{-f}V \bigg),\]
    which is syndetic since the set $\bigcap_{f \in F}T^{-f}V$ is an open neighborhood of $x$, as desired.
\end{proof}



\begin{theorem}
\label{thm_dct_combo_characterizations}
    Let $A \subseteq \N$.  The following are equivalent.
    \begin{enumerate}
        \item
        \label{item_dct_def_char}
        The set $A$ is a set of pointwise recurrence.

        \item
        \label{item_dct_sif_char}
        For all $B \supseteq A$, there exists a finite set $F \subseteq \N \setminus B$ such that the set $B \cup \big( B-F \big)$ is thick.

        \item
        \label{item_condition_one_dct}
        For all syndetic $S \subseteq \N$, there exists a finite set $F \subseteq A$ such that for all syndetic $S' \subseteq S$, the set $F \cap (S - S')$ is nonempty.

        \item \label{item:complement_thick_dct_characterization_v2}
        For all syndetic $S \subseteq \N$ and all thick $H \subseteq \N$, the piecewise syndetic set $P \defeq S \cap H$ satisfies the following.  There exists a finite set $F \subseteq A$ such that for all syndetic $S' \subseteq S$, the piecewise syndetic set $P' \defeq S' \cap H$ is such that the set $P' \cap (P - F)$ is nonempty.
        
        \item
        \label{item:complement_psstar_dct_characterization_star_fixed}
        For all piecewise syndetic $S \subseteq \N$ and all minimal left ideals $L \subseteq \beta \N$ with $\overline{S} \cap L \neq \emptyset$, there exists a finite set $F \subseteq A$ such that $\overline{S} \cap L \subseteq \overline{S-F}$.
    \end{enumerate}
\end{theorem}

\begin{proof}
    \eqref{item_dct_def_char} $\iff$ \eqref{item_dct_sif_char} \ Because $\dcsyndetic$ and $\dcthick$ are dual, the set $A$ is a set of pointwise recurrence if and only if the set $A' \defeq \N \setminus A$ is not dynamically central syndetic.  By \cref{intro_mainthm_full_chars_of_dcsyndetic}, the set $A'$ is not dynamically central syndetic if and only if for all $B' \subseteq A'$, there exists a finite set $F \subseteq B'$ such that the set
    \[B' \cap \bigcap_{f \in F} (B' - f)\]
    is not syndetic.  Taking the complement of this set in $\N$, we see that this happens if and only if the set
    \[\big( \N \setminus B' \big) \cup \bigcup_{f \in F} \big( (\N \setminus B') - f \big) = (\N \setminus B') \cup \big( (\N \setminus B')  - F \big)\]
    is thick.  Set $B \defeq \N \setminus B'$. Note that $B' \subseteq A'$ if and only if $A \subseteq B$.
    
    Summarizing the previous paragraph, we have shown that the set $A$ is a set of pointwise recurrence if and only if for all $B \supseteq A$, there exists a finite set $F \subseteq \N \setminus B$ such that the set $B \cup (B - F)$ is thick, as was to be shown.

    \eqref{item_dct_def_char} $\implies$ \eqref{item:complement_psstar_dct_characterization_star_fixed} \ Suppose that $A \in \dcthick$.  Let $S \subseteq \N$ be piecewise syndetic and $L \subseteq \beta \N$ be a minimal left ideal such that $\overline{S} \cap L \neq \emptyset$.  Note that $\overline{S} \cap L$ is nonempty and clopen in $L$.  By the definitional robustness of $\dcthick$ sets regarding non-metrizable systems (see \cref{sec_dy_thick_def_and_char}), for all $p \in \overline{S} \cap L$, there exists $a \in A$ such that $a + p \in \overline{S}$. Therefore,
    \[\overline{S} \cap L \subseteq \bigcup_{a \in A} \overline{S - a}.\]
    Because $\overline{S} \cap L$ is compact, there exists a finite set $F \subseteq A$ such that
    \[\overline{S} \cap L \subseteq \bigcup_{f \in F} \overline{S - f} \subseteq \overline{S-F},\]
    as was to be shown.

    \eqref{item:complement_psstar_dct_characterization_star_fixed} $\implies$ \eqref{item:complement_thick_dct_characterization_v2} \ Let $S \subseteq \N$ be syndetic and $H \subseteq \N$ be thick, and put $P \defeq S \cap H$.  It follows by \cite[Thm. 4.48]{hindman_strauss_book_2012} that there exists a minimal left ideal $L \subseteq \overline{H}$.  Note that since $S$ is syndetic, we have that $\emptyset \neq \overline{S} \cap L = \overline{S} \cap L \cap \overline{H} = \overline{P} \cap L$.
    
    Let $F \subseteq A$ be the finite set guaranteed by \eqref{item:complement_psstar_dct_characterization_star_fixed} for the piecewise syndetic set $P$.  Let $S' \subseteq S$ be syndetic, and put $P' \defeq S' \cap H$. Since $S'$ is syndetic and $S' \subseteq S$, we have that $\emptyset \neq \overline{S'} \cap L = \overline{P'} \cap L$.  Since $\emptyset \neq \overline{P'} \cap L \subseteq \overline{P} \cap L \subseteq \overline{P-F}$, we see that $\overline{P'} \cap \overline{P-F} \neq \emptyset$, whereby $P' \cap (P-F) \neq \emptyset$, as desired.
    
    \eqref{item:complement_thick_dct_characterization_v2} $\implies$ \eqref{item_condition_one_dct} \ Let $S \subseteq \N$ be syndetic.  Get $F \subseteq A$ finite from \eqref{item:complement_psstar_dct_characterization_star_fixed} for $H = \N$.  Then for all $S' \subseteq S$ syndetic,
    \[S' \cap (S-F) \neq \emptyset,\]
    whereby $F \cap (S-S') \neq \emptyset$, as desired.

    \eqref{item_condition_one_dct} $\implies$ \eqref{item_dct_def_char} \ We will prove the contrapositive: if $A$ is not a set of pointwise recurrence, then there exists a syndetic set $S \subseteq \N$ such that for all finite $F \subseteq A$, there exists a syndetic set $S' \subseteq S$ such that $F \cap (S - S') = \emptyset$. Since $F \cap (S - S') = \emptyset$ if and only if $S' \cap (S - F) = \emptyset$, it is equivalent to show: there exists a syndetic set $S \subseteq \N$ such that for all finite $F \subseteq A$, the set $S \cap (\N \setminus (S-F))$ is syndetic.  Taking complements and considering $B$ as $\N \setminus A$ and $C$ as $\N \setminus S$, we must show: if a set $B \subseteq \N$ is dynamically central syndetic, then there exists a set $C \subseteq \N$ that is not thick such that for all finite $F \subseteq \N \setminus B$, the set $(\N \setminus C) \cap \bigcap_{f \in F} (C - f)$ is syndetic.

    Suppose $B \subseteq \N$ is dynamically central syndetic.  If $B = \N$, put $C = \emptyset$ and note that the conclusion holds since the empty intersection is $\N$. Suppose $B \neq \N$. By \cref{thm:simple_equivalent_dcS}, there exists a minimal system $(X,T)$, a point $x \in X$, and a clopen set $U \subseteq X$ with $x \in U$ such that $R(x,U) \subseteq B$.  Define $V \defeq X \setminus U$ and $C \defeq \N \setminus R(x,U) = R(x, V)$. Since $U \neq X$ and the set $R(x,U)$ is syndetic, the set $C$ is nonempty but not thick.  Let $F \subseteq \N \setminus B$ be finite.  Since $F \subseteq C$, we have that for all $f \in F$, the set $T^{-f}V$ is an open neighborhood of $x$.  We see that
    \[\big(\N \setminus C\big) \cap \bigcap_{f \in F} (C - f) = R\bigg( x, U \cap \bigcap_{f \in F}T^{-f}V \bigg),\]
    which is syndetic since the set $U \cap \bigcap_{f \in F}T^{-f}V$ is an open neighborhood of $x$, as desired.
\end{proof}



\subsubsection{Examples of dynamically thick sets via the combinatorial characterizations}
\label{sec_examples_using_combinatorics}
 

In this section, we will use the combinatorial characterizations from the previous section to give two concrete examples of dynamically thick sets.


\begin{lemma}
\label{rem:combinatorial_primes}
    Let $(p_i)_{i \in \N} \subseteq \N$ be a sequence of distinct primes, $(c_i)_{i \in \N} \subseteq \Z$, and $(H_i)_{i \in \N}$ be a sequence of thick sets. The set
    \[
        A \defeq \bigcup_{i=1}^\infty \big ((p_i \N  + c_i) \cap H_i \big)
    \]
    is dynamically thick.  Moreover, if for all $n \in \N$, we have that $c_n \not\equiv 0 \bmod {p_n}$ and, for all but finitely many distinct pairs $i, j$ of positive integers, the set $H_i \cap (H_j - n)$ is empty, then the set $A$ is not an $\IP$ set.
\end{lemma}

\begin{proof}
    We will apply \cref{lemma_our_fav_dy_thick} \eqref{item:dyn_thick_idempotent}.  Let $\N \supsetneq B \supseteq A$. We want to show there exists a finite set $F \subseteq \N \setminus B$ such that $B - F$ is thick. Consider two cases.

    Case 1: There exists $i \in \N$ for which the set $\N \setminus B$ contains a complete residue system modulo $p_i$. Let $F = \{f_0, \ldots, f_{p_i - 1}\} \subseteq \N \setminus B$ be a complete system of modulo $p_i$ residues. We see that
    \[B - F \supseteq A - F \supseteq \big((p_i \N + c_i) \cap H_i \big) - F \supseteq \bigcap_{\ell = 0}^{p_i - 1} (H_i - f_{\ell}),\]
    which is thick, as desired. 
    
    Case 2: For all $i \in \N$, the set $\N \setminus B$ avoids some modulo $p_i$ congruence class. In this case, for all $i \in \N$, there exists $a_i \in \N$ such that $B \supseteq p_i \N + a_i$. We will show that $B$ is thick. Indeed, let $k \in \N$. It suffices to show there exists $n \in \N$ such that $n + i \in p_i \N + a_i$ for all $1 \leq i \leq k$. This is equivalent to showing that there exists $n \in \N$ such that $n \equiv a_i - i \bmod p_i$ for $1 \leq i \leq k$. This is an immediate consequence of the Chinese Remainder Theorem since the $p_i$'s are distinct primes.

    To see that the set $A$ is not an $\IP$ set under the stipulated conditions, note that for all $n \in \N$, the set $A \cap (A-n)$ is contained in a finite union of sets of the form $(p \N + c ) \cap H$, where $c \not \equiv 0 \bmod p$ and $H$ is thick.  If $A$ was an $\IP$ set, there would exist $n \in \N$ for which $A \cap (A-n)$ is an $\IP$ set.  Since the family $\IP$ is partition regular, it would follow that a set of the form $(p \N + c ) \cap H$ is an $\IP$ set, which is clearly false.
\end{proof}



Recently, the second author with Koutsogiannis, Moreira, Pavlov, and Richter \cite{Koutsogiannis_Le_Moreira_Pavlov_Richter_interpolation_sets} gave an example of a set of pointwise recurrence for distal systems which is not an $\IP$ set, answering a question of Host, Kra and Maass \cite[Question 3.11]{HostKraMaass2016}. Choosing the thick sets $H_i$ appropriately, the set $A \defeq \bigcup_{i=1}^\infty \big((p_i \N + 1) \cap H_i \big)$ is, according to \cref{rem:combinatorial_primes}, dynamically thick but not an $\IP$ set.  This set provides a strong answer to the question of Host, Kra and Maass: the set $A$ is not only a set of pointwise recurrence for distal systems but a dynamically thick set.


\begin{remark}
\label{rmk_transitive_point_visit_times}
\label{rmk_why_transitive_works}

\cref{rem:combinatorial_primes} can be used to explain why the set described in \eqref{ex_transitive_point_visit_times} at the beginning of \cref{sec_dy_thick_def_and_char} is dynamically thick.  Indeed, let $x$ be a point with a dense orbit in the full shift $(\{0, 1\}^{\N},\sigma)$, and suppose $U \subseteq X$ contains the cylinder set of all of those words in $\{0,1\}^{\N}$ beginning with some word $w$ of length $\ell$. We see that
\[R(x,U) \supseteq \{n \in \N \ | \ w \text{ appears at the $n^{\text{th}}$ position in } x\}.\]
Let $p > \ell$ be prime, and let $w_1$ be any word of length $p$ starting with $w$. Let $k \in \N$, and let $w_2 = w_1 w_1 \cdots w_1$ be the word consisting of $k$ copies of $w_1$. Since $x$ is transitive, the word $w_2$ must appear in $x$: there exists $N_k \in \N$ and $c_k \in \{0, \ldots, p-1\}$ such that $p \{N_k, N_k+1, \ldots, N_k+k - 1\} + c_k \subseteq R(x,U)$.

Let $c \in \{0, \ldots, p-1\}$ be such that for infinitely many $k \in \N$, $c_k = c$.  We see that there exists a thick set $H \subseteq \N$ such that $(p \N + c) \cap H \subseteq R(x,U)$.  Since $p$ was an arbitrary prime greater than $\ell$, we see that the set $R(x,U)$ contains a set of the form in \cref{rem:combinatorial_primes} and hence is dynamically thick, as desired.
\end{remark}


The class of examples in the following lemma generalizes the dynamically thick set that appears in \eqref{eqn_intro_even_odd} in the introduction.

\begin{lemma}
\label{ex:type2-special}
    Let $k \in \N$ and $(H_i)_{i \in \N}$ be a sequence of thick subsets of $\N$. The set
    \[A \defeq \bigcup_{i=0}^{k-1} \big ((k \N  + i) \cap H_i \big)\]
    is dynamically thick.
\end{lemma}

\begin{proof}
    We will apply \cref{lemma_our_fav_dy_thick} \eqref{item:complement_thick_dt_characterization_v2}.  Let $S \subseteq \N$ be piecewise syndetic.  Since piecewise syndeticity is partition regular and $S = \bigcup_{i=0}^{k-1} \big(S \cap (k\N - i)\big)$, there exists $i \in \{0, \ldots, k-1\}$ such that $S \cap (k \N - i)$ is piecewise syndetic.  Since demonstrating \eqref{item:complement_thick_dt_characterization_v2} for a subset of $S+i$ suffices to demonstrate it for $S$, by replacing $S$ with $(S+i) \cap k\N$, we can proceed under the assumption that $S \subseteq k \N$.

      
    
    Write $S = B \cap C$ where $B$ is syndetic with gap lengths less than $\ell \in \N$ and $C$ is thick. For $i \in \{0, \ldots, k-1\}$, choose an interval $I_i$ of $H_i$ such that $|I_i| > \ell$ and $\max I_i < \min I_{i+1}$. The set 
    \[
        F \defeq \bigcup_{i=0}^{k-1} \big( (k\N + i) \cap I_i \big)
    \]
    is a finite subset of $A$. 
    We will show that for all $M \in \N$, there is a finite subset $S' \subseteq S$ for which the set $S' - F$ contains an interval of length greater than $M$. It will follow that the set $S - F$ is thick.

    Let $M \in \N$.  Choose $S' \subseteq S$ to be the set $S$ intersected with a long interval (whose length will be specified later) on which the distance between consecutive elements of $S$ is less than $\ell$.  Since $\ell < |I_i|$ and $S' \subseteq k\N$, the set $S' - \big((k\N + i) \cap I_i \big)$ contains the set $(k \N - i) \cap J_i$, where
    \[
        J_i \defeq \big\{\min S' - \max I_i, \ \ldots, \max S' - \min I_i \big\}.
    \]
    It follows that
    \[
        S' - F \supseteq \bigcup_{i=0}^{k-1} \big((k \N - i) \cap J_i\big) \supseteq \bigcap_{i=0}^{k-1} J_i = \big\{\min S' - \max I_0, \ \ldots, \max S' - \min I_{k-1} \big\},
    \]
    provided $\min S' - \max I_0 \leq \max S' - \min I_{k-1}$.  This interval has length $\max S' - \min S' - (\min I_{k-1} - \max I_0)$ which is greater than $M$ if $\max S' - \min S' > M + \min I_{k-1} - \max I_0$.
\end{proof}




\subsection{More general examples of dynamically thick sets}
\label{sec_dy_thick_examples}

In this section, we develop more sophisticated dynamical tools to generalize the examples of dynamically thick sets presented in \cref{rem:combinatorial_primes} and \cref{ex:type2-special}.  We also show that in some sense, every dynamically thick set takes the form exhibited in these examples.

\subsubsection{Dynamically thick sets from disjointness}
\label{sec_dy_thick_from_disjointness}

In this section, we generalize the example in \cref{rem:combinatorial_primes} to a wider class of dynamically thick sets.  We accomplish this by showing that every minimal system must be nearly disjoint from almost all systems in a disjoint collection of minimal systems.  This is reminiscent of the fact that in a Hilbert space, according to Bessel's inequality, every vector must be nearly orthogonal to almost all vectors in an orthonormal sequence.

We call a collection of minimal systems $(X_i, T_i)$, $i \in I$, \emph{disjoint} if the product system $(\Pi_i X_i, \Pi_i T_i)$ is minimal.  Note that by the definition of the product topology, this is equivalent to having that the product system of any finite subcollection of the systems is minimal. This definition of disjointness is the restriction of the one given in example \eqref{ex_transitive_point_visit_times} in \cref{sec_dy_thick_def_and_char} to minimal systems.

For our purposes below, a \emph{joining} of the systems $(X,T)$ and $(Y,S)$ is a subsystem $(Z \subseteq X \times Y, T \times S)$ of the product system for which $\pi_X Z = X$ and $\pi_Y Z = Y$. Thus, two minimal systems are disjoint if their only joining is the product system.  It is easy to see that if $(X,T)$ and $(Y,S)$ are both minimal, then every subsystem of $(X \times Y, T \times S)$ is a joining.  In this case, every joining of $(X,T)$ and $(Y,S)$ has a minimal subsystem that is a joining.

\begin{theorem}
\label{thm_bessel}
    Let $(Z_i,R_i)$, $i \in \N$, be a disjoint collection of minimal systems.  For all minimal $(X,T)$, all $\eps > 0$, and all sufficiently large $i$ (depending on the $(Z_i,R_i)$'s, $(X,T)$, and $\eps$), all joinings $(J,T \times R_i)$ of $(X,T)$ and $(Z_i,R_i)$ satisfy: for all $z \in Z_i$, the set $\pi_1(J \cap (X \times \{z\}))$ is $\eps$-dense in $X$.
\end{theorem}

\begin{proof}
    Let $(X,T)$ be minimal and $\eps > 0$. Suppose for a contradiction that there are infinitely many $i$'s for which there exists a joining $(J_i, T \times R_i)$ of $(X,T)$ and $(Z_i,R_i)$ and a point $y_i \in Z_i$ for which $\pi_1(J_i \cap (X \times \{y_i\}))$ is not $\eps$ dense in $X$.   To save on notation, by relabeling, we will ignore those $(Z_i,R_i)$'s which do not fall into this infinite set. By passing to minimal subsystems, we may assume without loss of generality that the joinings $(J_i, T \times R_i)$ are minimal. Since $J_i$ is closed, there exists an $\eps$-ball $U_i \subseteq X$ and an open set $V_i \subseteq Z_i$ containing $y_i$ such that $J_i \cap (U_i \times V_i) = \emptyset$.

    Fix $x_0 \in X$.  Since $(J_i,T \times R_i)$ is a joining of $(X,T)$ and $(Z_i,R_i)$ and is minimal, there exists $z_i \in Z_i$ such that $J_i = \overline{(T \times R_i)^{\N}(x_0,z_i)}$. Since $J_i \cap (U_i \times V_i) = \emptyset$, we have that
    \begin{align}
    \label{eqn_main_return-time_set_containment}
        R_{R_i}(z_i,V_i) \subseteq R_T(x_0, X \setminus U_i).
    \end{align}
    Let $c_1, c_2, \ldots, c_m \in X$ be an $\eps/2$ dense subset of $X$.  For $\ell \in \{1, \ldots, m\}$, define
    \[
        E_\ell = \{i \in \N \ \big| \ B_{\eps/2}(c_\ell) \subseteq U_i \}.
    \]
    Note that $\N = \bigcup_{\ell=1}^m E_\ell$. 
    There exists $\ell \in \{1, \ldots, m\}$ for which $E_\ell$ is infinite.  Put $U = B_{\eps/2}(c_\ell)$. It follows now from \eqref{eqn_main_return-time_set_containment} that
    \begin{align}
    \label{eqn_main_return-time_set_containment_two}
        H \coloneqq \bigcup_{i \in E_\ell} R_{R_i}(z_i,V_i) \subseteq  \bigcup_{i \in E_\ell} R_T(x_0, X \setminus U_i) \subseteq R_T(x_0, X \setminus U).
    \end{align}
    We will show that the set $H$ is thickly syndetic and, hence, thick.  This will yield a contradiction since, by the minimality of $(X,T)$, the set $R_T(x_0, X \setminus U)$ cannot be thick.

    Let $k \in \N$, and let $e_1, \ldots, e_k \in E_\ell$. We see that
    \[\bigcap_{i=1}^k (H-i) \supseteq \bigcap_{i=1}^k R_{R_{e_i}}(R_{e_i}^{i}z_{e_i},V_{e_i}) = R_{R_{e_1} \times \cdots \times R_{e_k}}\big((R_{e_1}^1z_{e_1}, \ldots, R_{e_k}^{k}z_{e_k}),V_{e_1} \times \cdots \times V_{e_k}\big),\]
    which is syndetic by the minimality of the product system $(Z_{e_1} \times \cdots \times Z_{e_k},R_{e_1} \times \cdots \times R_{e_k})$. Therefore, intervals of length $k$ appear in $H$ syndetically. Since $k \in \N$ was arbitrary, this shows that the set $H$ is thickly syndetic, as desired.
\end{proof}

\begin{theorem}
\label{thm_quantitative_strengthening_of_type_1_example}
    Let $(X_i,T_i)$, $i \in \N$, be a disjoint collection of minimal systems.  For all minimal systems $(Y,S)$ and all nonempty, open $V \subseteq Y$, there exists $i \in \N$ such that the following holds.  For all nonempty, open $U \subseteq X_i$, there exists $s \in \N$ such that for all $(x,y) \in X_i \times Y$, the set $R_{T_i \times S}((x,y), U \times V) = R_{T_i}(x,U) \cap R_S(y,V)$ is syndetic with gap size bounded by $s$.
\end{theorem}

\begin{proof}
    Let $(Y,S)$ be minimal and $V \subseteq Y$ be nonempty and open.  Let $\eps > 0$ be such that $V$ contains an $\eps$-ball.  According to \cref{thm_bessel}, there exists $i \in \N$ such that all joinings $(J,T_i \times S)$ of $(X_i,T_i)$ and $(Y,S)$ satisfy: for all $x \in X_i$, the set $\pi_2(J \cap (\{x\} \times Y))$ is $\eps$-dense in $Y$.

    Let $U \subseteq X$ be nonempty and open.  We claim that for all $(x,y) \in X_i \times Y$, there exists $n \in \N$ such that $(T_i \times S)^n (x,y) \in U \times V$.  Indeed, let $(x,y) \in X_i \times Y$. By a theorem of Auslander \cite{auslander_1960} and Ellis \cite{ellis_1957}, the point $(x,y)$ is proximal to a point $(x_0,y_0) \in X_i \times Y$ that is uniformly recurrent under $T_i \times S$.  Since $J \defeq \overline{(T_i \times S)^{\N_0}(x_0,y_0)}$ is a (minimal) joining of $(X_i,T_i)$ and $(Y,S)$, it is $\eps$-dense in every fiber over $X_i$.  Since $V$ contains an $\eps$-ball, we see that $J \cap (U \times V) \neq \emptyset$.  Because $(x,y)$ is proximal to $(x_0,y_0)$ and $R_{T_i \times S}((x_0,y_0), U \times V)$ is syndetic, we see that $R_{T_i \times S}((x,y),U \times V)$ is piecewise syndetic, hence nonempty.
    
    We have shown that for all $(x,y) \in X_i \times Y$, there exists $n \in \N$ such that $(T_i \times S)^n (x,y) \in U \times V$. It follows that $\bigcup_{n=1}^\infty (T_i \times S)^{-n} (U \times V) = X_i \times Y$.  By compactness, there exists $s \in \N$ such that $\bigcup_{n=1}^s (T_i \times S)^{-n} (U \times V) = X_i \times Y$. It follows that for all $(x,y) \in X_i \times Y$, there exists $n \in \{1, \ldots, s\}$ such that $(T_i \times S)^n(x,y) \in U \times V$, whereby for all $(x,y) \in X_i \times Y$, the set $R_{T_i \times S}((x,y), U\times V)$ is syndetic with gap size bounded by $s$, as was to be shown.
\end{proof}

The following is a generalization of \cref{rem:combinatorial_primes}.

\begin{theorem}%[Type 1 Example]
\label{lemma_more_general_example_of_dy_thick_not_IP}
    Let $(X_i,T_i)$, $i \in \N$, be a disjoint collection of minimal systems.  For each $i \in \N$, let $x_i \in X_i$, $U_i \subseteq X_i$ be nonempty and open, and $H_i \subseteq \N$ be thick.   The set
    \[
        A \defeq \bigcup_{i=1}^\infty \big( R_{T_i}(x_i, U_i) \cap H_i \big)
    \]
    is dynamically thick.
\end{theorem}

\begin{proof}
    To see that $A \in \dthick$, we must show that for all minimal $(Y,S)$, all nonempty, open $V \subseteq Y$, and all $y \in Y$, the set $A \cap R_S(y,V)$ is nonempty.  This follows immediately from \cref{thm_quantitative_strengthening_of_type_1_example}.
\end{proof}

\subsubsection{Dynamically thick sets from distal points}



In this section, we offer a simultaneous generalization of the examples in \cref{rem:combinatorial_primes} and \cref{ex:type2-special}. 


\begin{lemma}
\label{lemma_orbit_along_ps_is_somewhere_dense}
    Let $(X,T)$ be a minimal system and $x \in X$.  For all piecewise syndetic sets $P \subseteq \N$, the set $T^P x$ is somewhere dense, ie., $(\overline{T^P x})^\circ \neq \emptyset$.
\end{lemma}

\begin{proof}
    Since the set $P$ is piecewise syndetic, there exists $k \in \N$ such that $H \defeq \bigcup_{i=1}^k (P-i)$ is thick.  Since $(X,T)$ is minimal, the set $T^H x$ is dense.  We see that
    \[X = \overline{\bigcup_{i=1}^k T^{-i} T^{P}x} = \bigcup_{i=1}^k T^{-i} \overline{T^{P}x}.\]
    By the Baire Category Theorem, there exits $i \in \{1, \ldots, k\}$ such that the set $T^{-i} \overline{T^{P}x}$ has nonempty interior.  By semiopenness of minimal maps \cite[Thm 2.4]{kolyada_snoha_trofimchuk_2001}, applying $i$ many times the map $T$, the set $\overline{T^{P}x}$ has nonempty interior, as desired.
\end{proof}

The following theorem generalizes the examples in both \cref{rem:combinatorial_primes} and \cref{ex:type2-special}.  This can be seen by considering a periodic system in the latter case and an infinite product of periodic systems of prime cardinality in the former case.  Recall the definition of a distal point from \cref{sec_top_dynamics}.


\begin{theorem}%[Type 2 Example]
\label{prop_description_of_dy_thick}
    Let $(X, T)$ be a minimal system and $U_1, U_2, \ldots$ be nonempty, open subsets of $X$. Let $x \in X$ be a distal point and $H_1, H_2, \ldots \subseteq \N$ be thick sets. The set
    \[
        A \defeq \bigcup_i \big(R(x,U_i) \cap H_i \big)
    \]
    is dynamically thick if and only if the set $\bigcup_{i=1}^\infty U_i$ is dense in $X$.
\end{theorem}

\begin{proof}
    If the set $\bigcup_{i=1}^\infty U_i$ is not dense in $X$, then it is disjoint from a nonempty, open set $V \subseteq X$. The set $A$ is clearly disjoint from the set $R(x,V)$, whereby $A$ is not dynamically thick.

    Suppose that $\bigcup_{i=1}^\infty U_i$ is dense. Let $(Y,S)$ be a minimal system, $y \in Y$, and $V \subseteq Y$ be nonempty and open.  Since the set $R_S(y,V)$ is syndetic, by \cref{lemma_orbit_along_ps_is_somewhere_dense}, the set $T^{R_S(y,V)} x$ is somewhere dense. Therefore, there exists $i \in \N$ such that $T^{R_S(y,V)} x \cap U_i \neq \emptyset$.  It follows that the set $R_{T \times S} ((x,y), U_i \times V)$ is nonempty. Since $x$ is a distal point, by \cite[Thm. 9.11]{furstenberg_book_1981}, the point $(x,y)$ is uniformly recurrent under $T \times S$.  Since $U_i \times V$ is open, we have that $R_{T \times S} ((x,y), U_i \times V)$ is (dynamically) syndetic. Therefore, the set $R_{T \times S} ((x,y), U_i \times V) \cap H_i$ is nonempty, which implies that $A \cap R_S(y,V)$ is nonempty, as desired.
\end{proof}



\subsubsection{A general form for dynamically thick sets: a proof of \cref{mainthm_structure_for_dy_thick_sets}}
\label{sec_general_form_of_dthick}

In this section, we draw on a future result, \cref{thm_dcps_and_dps_are_ps}, to describe the form of any dynamically thick set.  (For the astute reader concerned about circular logic: while \cref{cor_weak_structure_result} relies on \cref{thm_dcps_and_dps_are_ps}, we do not invoke \cref{cor_weak_structure_result} anywhere in the paper.)


\begin{definition}
\label{def_robust_syndetic}
A collection $\collection$ of subsets of $\N$ is \emph{robustly syndetic} if for all dynamically syndetic sets $A \subseteq \N$, there exists $B \in \collection$ such that $A \cap B$ is syndetic.
\end{definition}

The dynamically thick sets described so far all take the form $\bigcup_{B \in \collection} (B \cap H_B)$, where $\collection$ is a robustly syndetic collection of sets and $(H_B)_{B \in \collection}$ is a collection of thick sets.  Indeed, in \cref{lemma_more_general_example_of_dy_thick_not_IP}, the collection $\{R_{T_i}(x_i,U_i) \ | \ i \in \N\}$ is robustly syndetic by virtue of the disjointness of the collection of systems $(X_i, T_i)$, $i \in \N$.  In \cref{prop_description_of_dy_thick}, the collection $\{R_{T}(x,U_i) \ | \ i \in \N\}$ is robustly syndetic by virtue of the fact that $x$ is a distal point and $\bigcup_{i=1}^\infty U_i$ is dense in $X$.  \cref{mainthm_structure_for_dy_thick_sets} shows that every dynamically thick set has an underlying robustly syndetic collection.



\begin{proof}[Proof of \cref{mainthm_structure_for_dy_thick_sets}]
    Suppose that $\collection$ is robustly syndetic and that, for each $B \in \collection$, the set $H_B \subseteq \N$ is thick such that \eqref{eqn_A_contains_family_cap_thick} holds. To see that $A \in \dthick$, let $C \in \dsyndetic$.  There exists $B \in \collection$ such that $B \cap C$ is syndetic. Thus, the set $B \cap C \cap H_B$ is nonempty, whereby $A \cap C \neq \emptyset$, as desired.

    Conversely, suppose that $A$ is dynamically thick.  For all $S \in \dsyndetic$, we will define a thick set $G_S \subseteq \N$ and a set $B_S \subseteq \N$ in such a way that the collection $\collection \defeq \{ B_S \ | \ S \in \dsyndetic\}$ is robustly syndetic and
    \[A \supseteq \bigcup_{S \in \dsyndetic} \big( B_S \cap G_S \big).\]
    This suffices to reach the conclusion: the collection $\collection \cup \{A\}$ is robustly syndetic, and setting $H_A \defeq \N$ and, for $B = B_S \in \collection$, setting $H_B \defeq G_{S}$, we see that equality in \eqref{eqn_A_contains_family_cap_thick} holds.

    Let $S \in \dsyndetic$. The set $A \cap S$ is an intersection of a dynamically thick set and a dynamically syndetic set and, hence, by \cref{thm_dcps_and_dps_are_ps}, is piecewise syndetic.  There exist $\ell \in \N$ and finite, disjoint intervals $I_1$, $I_2$, \dots of consecutive positive integers satisfying: a) $\lim_{i \to \infty} |I_i| \to \infty$ and, for all $i \in \N$, b) $\max I_i  + i< \min I_{i+1}$; and c) the set $A \cap S$ has nonempty intersection with all subintervals of $I_i$ of length at least $\ell$.  
    
    Define $G_S \defeq \bigcup_{i=1}^\infty I_i$, and note that it is thick.  Define
    \begin{align}
    \label{eq:B_S}
        B_S \defeq (A \cap G_S) \cup (\N \setminus G_S).
    \end{align}
    Since the set $\N \setminus G_S$ is a disjoint union of finite intervals whose lengths tend to infinity, the set
    \begin{align}
    \label{eq:B_ScapS}
        B_S \cap S = (A \cap G_S \cap S) \cup \big((\N \setminus G_S) \cap S \big) \text{ is syndetic.}
    \end{align}
    Indeed, the set $S$ has nonempty intersection with any interval of length $\ell' \in \N$.  Therefore, on $G_S$, the set $B_S \cap S$ has nonempty intersection with any interval of length at least $\ell$, while on $\N \setminus G_S$, the set $B_S \cap S$ has nonempty intersection with any interval of length $\ell'$.
    
    Now \eqref{eq:B_ScapS} shows that the collection $\collection \defeq \{B_S \ | \ S \in \dsyndetic\}$ is robustly syndetic, and \eqref{eq:B_S} shows that $A \supseteq \bigcup_{S \in \dsyndetic} (B_S \cap G_S)$, as desired.
\end{proof}

It would be interesting and useful to improve \cref{cor_weak_structure_result} by saying more about the robustly syndetic collection $\collection$.  We discuss this further in \cref{quest_robust_synd_family}.



\subsection{The partition problem and \texorpdfstring{$\sigma$}{sigma}-compactness}
\label{sec:sigma-compact}

It is natural to ask whether any dynamically thick set can be partitioned into two disjoint dynamically thick sets.  While we were not able to answer this question -- see \cref{ques:partition_into_two_pointwise_recurrence} for some discussion on how the structure result in \cref{cor_weak_structure_result} may help -- we were able to show that the established technique for answering this type of question does not work.

The established technique goes as follows.  Let $\family$ be a family of subsets of $\N$, and suppose we wish to show that for all $A \in \family$, there exists a disjoint union $A = A_1 \cup A_2$ such that $A_1, A_2 \in \family$.  Let $N_1 < N_2 < \cdots$ be a rapidly increasing sequence of positive integers.  Assign the set $A \cap [N_1]$ to $A_1$, the set $A \cap ([N_2] \setminus [N_1])$ to $A_2$, the set $A \cap ([N_3] \setminus [N_2])$ to $A_1$, and so forth.  If the family $\family$ is ``suitable'' and the $N_i$'s are chosen appropriately, the sets $A_1$ and $A_2$ will both belong to $\family$. This outline is realized in detail in the proof of \cref{prop:sigma-compact-implies-partition-to-two}.

There are several instances of this in the literature \cite{Forrest_recurrence-in-dynamical-systems,Le-interpolation-first,Le_interpolation_for_nilsequences}, but to our knowledge, sufficient and necessary conditions for ``suitability'' have never been explained.  In all instances that we know of, ``suitability'' can be succinctly captured in terms of compactness.  We show in \cref{prop:sigma-compact-implies-partition-to-two} below that if the dual family $\family^*$ is $\sigma$-compact (\cref{def:compact}), then all members of $\family$ can be partitioned into two disjoint members of $\family$.  We give several examples of $\sigma$-compact families that arise naturally in the subject.

The main result in this section is \cref{theorem_dy_syndetic_is_not_sigma_compact}: the families $\dthick^* = \dsyndetic$ and $\dcthick^* = \dcsyndetic$ are not $\sigma$-compact.  This demonstrates why the usual procedure for splitting dynamically (central) thick sets into two such sets fails.  It also indicates that these families are more complex than others usually considered in the subject. 

\subsubsection{Compact and \texorpdfstring{$\sigma$}{sigma}-compact families}
\label{sec:compact-sigma-compact-first}

The terminology in the following definition is motivated in \cref{prop:sigma-compact-equivalent} below.

\begin{definition}
\label{def:compact}
\label{def:uniformity}
Let $\family$ be a family of subsets of $\N$.
\begin{enumerate}
    \item The family $\family$ is \emph{compact} if for all $B \in \family^*$, there exists $M \in \N$ such that $B \cap [M] \in \family^*$.

    \item The family $\family$ is \emph{$\sigma$-compact} if it is the union of countably many compact families.
\end{enumerate}
\end{definition}

Before giving some examples of compact and $\sigma$-compact families, we move to explain the terminology.  For a family $\class$ of subsets of $\N$, define
\[X_{\class} \defeq \big\{1_A \in \{0, 1\}^{\N} \ \big| \ A \in \class\big\}.\]
Thus, the set $X_{\class}$ is a subset of the compact metric space $\{0, 1\}^{\N}$.



\begin{theorem}
\label{prop:sigma-compact-equivalent}
    A family $\family$ is compact (resp. $\sigma$-compact) if and only if the set $X_{\class}$ is compact (resp. $\sigma$-compact).
\end{theorem}

\begin{proof}
Since $X_{\bigcup_{N=1}^\infty \family_N} = \bigcup_{N = 1}^\infty X_{\family_N}$, it suffices to show that $\family$ is compact if and only if $X_{\class}$ is compact.


Suppose $X_{\class}$ is a compact subset of $\{0, 1\}^{\N}$. Assume for the sake of a contradiction that $\class$ is not a compact family. Thus, there exists $B \in \family^*$ such that for every $M \in \N$, there exists $A_M \in \family$ such that
\begin{align}
\label{eqn_f_not_compact}
    B \cap [M] \cap A_M = \emptyset.
\end{align}

Since $X_{\family}$ is compact, the sequence $(1_{A_M})_{M \in \N} \subseteq X_{\family}$ has a subsequence that converges to some point $1_A \in X_{\family}$. By the product topology on $\{0,1\}^{\N}$ and \eqref{eqn_f_not_compact}, we see that $B \cap A = \emptyset$, in contradiction with the fact that $A \in \family$ and $B \in \family^*$. 

Conversely, suppose that $\family$ is a compact family. We will show that $X_{\family}$ is compact by showing that it is closed. Let $(x_k)_{k \in \N} \subseteq X_{\family}$ be a convergent sequence with limit point $x \in \{0,1\}^{\N}$. We will show that $x \in X_{\family}$ by showing that for all $B \in \family^*$, the set $\supp(x) \cap B$ is nonempty.

Let $B \in \family^*$.  Since $\family$ is compact, there exists $M \in \N$ such that for all $k \in \N$,
\[
    B \cap [M] \cap \supp(x_k) \neq \emptyset.
\]
Since $\lim_{k \to \infty}x_k = x$ and $B \cap [M]$ is a finite set, we see that $B \cap [M] \cap \supp(x) \neq \emptyset$. In particular, the set $B \cap \supp(x)$ is nonempty, as was to be shown.
\end{proof}


Next we provide some examples of compact and $\sigma$-compact families.

\begin{example}
\label{ex_syndetic_is_sigma_compact}
    For $N \in \N$, denote by $\syndetic_N$ the family of sets that have nonempty intersection with every interval in $\N$ of length $N$.  The family $\syndetic_N$ is compact.  Indeed, given $B \in \syndetic_N^*$, the set $B$ is thick.  If $M \in \N$ is large so that $B \cap [M]$ contains an interval of length at least $N$, then $B \cap [M] \in \syndetic_N^*$.  We see that the family of syndetic sets, $\syndetic$, is equal to $\bigcup_{N=1}^\infty \syndetic_N$, whereby it is $\sigma$-compact.
\end{example}


\begin{example}
\label{ex_ip_k_sets_are_compact}
    Let $N \in \N$. Denote by $\IP_N$ the upward closure of the collection of sets of the form
    \[\FS(x_i)_{i=1}^N \defeq \left\{ \sum_{f \in F} x_f \ \middle | \ \emptyset \neq F \subseteq \{1, \ldots, N\} \right\}, \quad x_1, \ldots, x_N \in \N.\]
    As in the previous example, the family $\IP_N^*$ is easily seen from the definition to be compact.  The dual of the family $\IP_0 \defeq \bigcap_{N = 1}^\infty \IP_N$ is $\sigma$-compact, since $\IP_0^* = \bigcup_{N=1}^\infty \IP_N^*$.  The family $\IP_0^*$ is important in quantitative strengthenings of recurrence theorems and appears again in this paper in \cref{sec_open_quests_dps_sets}.
\end{example}

\begin{example}
\label{example:set-topological-partition}
    The family of sets of return times in minimal systems is $\sigma$-compact. More precisely, let $\class$ be the upward closure of collection of sets of the form
    \begin{equation}\label{eq:sigma_compact_topological}
        \big\{n \in \N \ \big| \ U \cap T^{-n} U \neq \emptyset \big\},
    \end{equation}
    where $(X, T)$ is a minimal system and $U \subseteq X$ is a nonempty, open set. The family $\class$ is $\sigma$-compact, as shown in \cite[Prop. 5.8]{boshernitzan_glasner_2009}. The argument is short, so we give it here.
    
    It is a fact \cite[Thm. 2.4]{bergelson_mccutcheon_1998} that $A \in \class$ if and only if there exists a syndetic set $S \subseteq \N$ such that $A \supseteq S-S$. Thus, using the notation from \eqref{ex_syndetic_is_sigma_compact}, if for $N \in \N$ we define
    \[\class_N \defeq \upclose \big\{ S - S \ \big| \ S \in \syndetic_N \big\},\]
    we have that $\class = \bigcup_{N=1}^{\infty} \class_N$.  We have only to show that for all $N \in \N$, the family $\class_N$ is compact. Let $N \in \N$ and $B \in \class_N^*$.  We must show that there exists $M \in \N$ such that for all $A \in \class_N$, the set $A \cap B \cap [M]$ is nonempty.

    Assume for the sake of a contradiction that no such $M$ exists.  Thus, for all $M \in \N$, there exists $S_M \in \syndetic_N$ such that
    \begin{align}
    \label{eqn_sm_limit}
        (S_M - S_M) \cap B \cap [M] = \emptyset.
    \end{align}
    Since $\syndetic_N$ is compact, passing to a subsequence of $(S_M)_{M \in \N}$, there exists $S \in \syndetic_N$ such that $\lim_{M \to \infty} 1_{S_M} = 1_S$.  Since $S \in \syndetic_N$, the set $S-S$ is a member of $\class_N$.  It follows from \eqref{eqn_sm_limit} that $(S-S) \cap B = \emptyset$, contradicting the fact that $B \in \class_N^*$.
\end{example}

\begin{example}
\label{example:set-measurable-partition}
    In analogy to \cref{example:set-topological-partition}, the family of sets of return times in measure preserving systems is $\sigma$-compact. More precisely, the upward closure $\family$ of the collection of sets of the form
    \begin{equation}\label{eq:sigma_compact_measurable}
        R(E,E) \defeq \big\{n \in \N \ \big| \ \mu(E \cap T^{-n} E)  > 0 \big\}, 
    \end{equation}
    where $(X, \mu, T)$ is a measure preserving system (see \cref{sec_dcps_and_set_recurrence}) and $E \subseteq X$ is a set of positive measure, is $\sigma$-compact.  This is most readily seen as a consequence of a result of Forrest \cite{Forrest_recurrence-in-dynamical-systems}: for all sets $B \subseteq \N$ of single measurable recurrence (i.e. for all $B \in \class^*$) and all $\delta > 0$, there exists $M \in \N$ such that for all sets $E$ with measure at least $\delta$, the set $B \cap R(E,E) \cap [M]$ is nonempty. For $N \in \N$, defining
    \[
        \class_N \defeq \upclose \big\{R(E,E) \subseteq \N \ \big| \ \text{$(X,\mu,T)$ is a measure preserving system, $E \subseteq X$, $\mu(E) > 1/N$} \big\},
    \]
    we see that $\class_N$ is compact and $\family = \bigcup_{N =1}^\infty \family_N$ is $\sigma$-compact.
\end{example}


The previous two examples concerned set recurrence, while the following ones concern pointwise recurrence.
\begin{example}
    Bohr$_0$ sets and  Nil$_k$-Bohr$_0$ sets are special dynamically central syndetic sets for which the systems in consideration are minimal rotations on compact abelian groups and minimal $k$-step nilsystems, respectively. (See \cite{Host-Kra_nilbohr} for a detailed definition of Nil$_k$-Bohr$_0$ sets.)
    It can be shown that these families are $\sigma$-compact using the results in \cite[Prop. 1.4]{Le-interpolation-first} and \cite[Lemma 3.3]{Le_interpolation_for_nilsequences}, respectively.
\end{example}




\subsubsection{The partition problem in dual compact families}

The following result offers, under a mild assumption on $\family \classcap \family^*$, a solution to the partition problem in families whose duals are $\sigma$-compact.

\begin{theorem}
\label{prop:sigma-compact-implies-partition-to-two}
    Let $\family$ be a family of subsets of $\N$ with the property that all members of $\family \classcap \family^*$ are infinite.  If $\family^*$ is $\sigma$-compact, then for all $A \in \family$, there exists a disjoint partition $A = A_1 \cup A_2$ with $A_1, A_2 \in \family$.
\end{theorem}

\begin{proof}
    Suppose $\family^*$ is $\sigma$-compact. There exists $\family^* = \bigcup_{N=1}^{\infty} \familytwo_N$ where each family $\familytwo_N$ is compact. By taking finite unions, we may assume without loss of generality that $\familytwo_1 \subseteq \familytwo_2 \subseteq \cdots$.
    
    Let $A \in \family$. Since $\familytwo_1 \subseteq \family^*$, the set $A$ belongs to $\familytwo_1^*$. Since $\familytwo_1$ is compact, there exists a finite set $A_1' \subseteq A$ such that $A_1' \in \familytwo_1^*$. Since all members of $\family \classcap \family^*$ are infinite, for all $B \in \family^*$, the set $A \cap B$ is infinite, whereby
    \[
        \big(A \setminus A_1' \big) \cap B \neq \emptyset. 
    \]
    Since $B \in \family^*$ was arbitrary, we see that $A \setminus A_1' \in \family$. Repeating the argument for $A \setminus A_1'$ in place of $A$, we get a finite set $A_2' \subseteq A \setminus A_1'$ such that $A_2' \in \familytwo_2^*$. By the same argument as before, the set $A \setminus (A_1' \cup A_2')$ is a member of $\family$, allowing us to continue on as before.
    
    Repeating ad infinitum, we get a sequence $A_1'$, $A_2'$, \dots of disjoint, finite subsets of $A$ such that for all $N \in \N$, the set $A_N' \in \familytwo_N^*$.  Define $A_1 \defeq A_1' \cup A_3' \cup \cdots$ and $A_2 \defeq A_2' \cup A_4' \cup \cdots$.  By appending any elements in $A \setminus (A_1 \cup A_2)$ to $A_1$, we may assume that $A = A_1 \cup A_2$.
    
    We need only to show that $A_1$ and $A_2$ are members of $\family$.  Indeed, since $\family^* = \bigcup_{N=1}^{\infty} \familytwo_N$, we see that $\family = (\family^*)^* = \bigcap_{N=1}^{\infty} \familytwo_N^*$.  Since $\familytwo_1^* \supseteq \familytwo_2^* \supseteq \cdots$, by the construction of the $A_i$'s, both $A_1$ and $A_2$ belong to $\bigcap_{N=1}^{\infty} \familytwo_N^*$.
\end{proof}



\begin{remark}
    In view of \cref{prop:sigma-compact-implies-partition-to-two} and the examples of $\sigma$-compact families presented in \cref{sec:compact-sigma-compact-first}, every set in the following list of families can be partitioned into two sets in the same family:
    \begin{enumerate}
        \item sets of topological recurrence (dual of sets of the form in \eqref{eq:sigma_compact_topological}),
        \item sets of measurable recurrence (dual of sets of the form in \eqref{eq:sigma_compact_measurable}), 
        \item sets of Bohr recurrence (dual of Bohr$_0$ sets),
        \item sets of pointwise recurrence for $k$-step nilsystems (dual of Nil$_k$-Bohr$_0$ sets).
    \end{enumerate}
    These facts are well known to experts, but the framing in terms of compactness is, as far as we know, new.
\end{remark}



\subsubsection{The family of dynamically syndetic sets is not \texorpdfstring{$\sigma$}{sigma}-compact}

We show in this section that the families $\dsyndetic$ and $\dcsyndetic$ are not $\sigma$-compact.

\begin{theorem}
\label{thm_suff_condition_for_dual_not_sigma_compact}
    Let $\classtwo$ be a family of subsets of $\N$.  If
    \begin{align}
        \label{eqn_key_property_to_refute_sigma_compactness_one}
        \begin{gathered}
            \text{there exist $B_1, B_2, \ldots \in \classtwo$, disjoint, such that for all $B \in \classtwo$} \\ \text{with $B \subseteq \bigcup_{i=1}^\infty B_i$, there exists $i \in \N$ such that $|B \cap B_i| = \infty$,}
        \end{gathered}
    \end{align}
    then the family $\classtwo^*$ is not $\sigma$-compact.
\end{theorem}


\begin{proof}
    Suppose property \eqref{eqn_key_property_to_refute_sigma_compactness_one} holds, and suppose for a contradiction that $\classtwo^* = \bigcup_{i=1}^\infty \class_i$, where $\class_1$, $\class_2$, \dots are compact families.

    Let $i \in \N$.  We claim that there exists a function $\varphi_i: \classtwo \to \N$ such that for all $A \in \class_i$ and for all $B \in \classtwo$,
    \[\min (A \cap B) \leq \varphi_i(B).\]
    Indeed, fix $B \in \classtwo$, and consider the function $X_{\class_i} \to \N$ defined by $1_A \mapsto \min (A \cap B)$.  It is locally constant since, more generally, the function $X_{\classtwo} \times X_{\classtwo^*} \to \N$ defined by $(1_A,1_B) \mapsto \min (A \cap B)$ is locally constant. Since $\class_i$ is compact, the function $X_{\class_i} \to \N$ is bounded from above.  We define $\varphi_i(B)$ to be an upper bound.

    Define the set
    \[A \defeq \left( \N \setminus \bigcup_{i=1}^\infty B_i \right) \cup \bigcup_{i=1}^\infty \big( B_i \setminus \{1, \ldots, \varphi_i(B_i) \} \big).\]
    We claim that $A \in \classtwo^*$.  Indeed, if $B \in \classtwo$, then either $B \cap ( \N \setminus \bigcup_{i=1}^\infty B_i ) \neq \emptyset$, in which case $A \cap B \neq \emptyset$, or $B \subseteq \bigcup_{i=1}^\infty B_i$.  In the latter case, by \eqref{eqn_key_property_to_refute_sigma_compactness_one}, there exists $i \in \N$ such that $|B \cap B_i| = \infty$.  It follows that $B \cap (B_i \setminus \{1, \ldots, \varphi_i(B_i) \}) \neq \emptyset$, whereby $A \cap B \neq \emptyset$.

    Since $A \in \classtwo^*$, there exists $j \in \N$ such that $A \in \class_j$.  It follows that $\min (A \cap B_j) \leq \varphi_j(B_j)$. But since the $B_j$'s are disjoint, we see from the definition of $A$ that
    \[A \cap B_j = B_j \setminus \{1, \ldots, \varphi_j(B_j) \}\]
    whereby $\min (A \cap B_j) > \varphi_j(B_j)$, a contradiction.
\end{proof}

The following lemma gives a sufficient condition to satisfy the property in \eqref{eqn_key_property_to_refute_sigma_compactness_one}.  Recall that all difference sets are computed in $\N$.

\begin{lemma}
\label{lemma_suff_condition_for_key_property_for_dual_not_sigma_compact}
    Let $\classtwo$ be a family of subsets of $\N$.  If
    \begin{align}
        \label{eqn_suff_condition_for_key_prop_one}
        \text{for all $B \in \classtwo$, there exists $n \in \N$ such that $|B \cap (B-n)| = \infty$,}
    \end{align}
    and if
        \begin{align}
        \label{eqn_suff_condition_for_key_prop_two}
        \begin{gathered}
            \text{there exist $B_1, B_2, \ldots \in \classtwo$, disjoint, such that} \\ \lim_{d \to \infty} \min \bigcup_{\substack{i, j = 1 \\ \max(i,j) > d}}^\infty (B_i - B_j) = \infty,
        \end{gathered}
    \end{align}
    then the property in \eqref{eqn_key_property_to_refute_sigma_compactness_one} holds.
\end{lemma}

\begin{proof}
    Suppose the properties in \eqref{eqn_suff_condition_for_key_prop_one} and \eqref{eqn_suff_condition_for_key_prop_two} hold.  We will show that the property in \eqref{eqn_key_property_to_refute_sigma_compactness_one} holds for the same sets $B_1$, $B_2$, \dots.

    Let $B \in \classtwo$ be such that $B \subseteq \bigcup_ {i=1}^\infty B_i$.  By \eqref{eqn_suff_condition_for_key_prop_one}, there exists $n \in \N$ such that $|B \cap (B-n)| = \infty$.  Writing $B = B \cap \bigcup_{i=1}^\infty B_i$, we see that
    \[B \cap (B-n) = \bigcup_{i,j = 1}^\infty \big( B \cap B_i \cap (B-n) \cap (B_j - n)\big).\]
    By \eqref{eqn_suff_condition_for_key_prop_two}, there exists $d \in \N$ such that if $\max(i,j) > d$, then $B_i \cap (B_j - n) = \emptyset$.  Therefore,
        \[B \cap (B-n) = \bigcup_{i,j = 1}^d \big( B \cap B_i \cap (B-n) \cap (B_j - n)\big) \subseteq \bigcup_{i=1}^d \big( B \cap B_i \big).\]
    By assumption, the set on the left hand side is infinite, so there must exist an $i \in \{1, \ldots, d\}$ such that $B \cap B_i$ is infinite.  This verifies \eqref{eqn_key_property_to_refute_sigma_compactness_one}.
\end{proof}


The result in the following lemma is well known.  We include a proof for completeness.

\begin{lemma}
\label{lem:piecewise_syndetic_shift}
    If $B \subseteq \N$ is piecewise syndetic, then there exists $n \in \N$ such that the set $B \cap (B - n)$ is piecewise syndetic.
\end{lemma}

\begin{proof}
    Suppose $B \subseteq \N$ is piecewise syndetic. There exists a syndetic set $S \subseteq \N$ and a thick set $H \subseteq \N$ such that $B = S \cap H$.  There exists $N \in \N$ such that $\bigcup_{i=1}^N (S-i) = \N$, and so $\bigcup_{i=1}^N (B-i) \supseteq \bigcap_{i=1}^N (H-i)$. It follows that the set 
    \[
        B \cap \bigcup_{i=1}^N (B-i) \supseteq S \cap H \cap \bigcap_{i=1}^N (H-i)
    \]
    is piecewise syndetic.  Since the family $\PS$ is partition regular, there exists $n \in \{1, \ldots, N\}$ such that the set $B \cap (B-n)$ is piecewise syndetic.
\end{proof}

\begin{theorem}
\label{theorem_dy_syndetic_is_not_sigma_compact}
    The families $\dthick^* = \dsyndetic$ and $\dcthick^* = \dcsyndetic$ are not $\sigma$-compact.
\end{theorem}

\begin{proof}
    We will show that $\dthick$ and $\dcthick$ satisfy the conditions in \cref{lemma_suff_condition_for_key_property_for_dual_not_sigma_compact} and hence in \eqref{eqn_key_property_to_refute_sigma_compactness_one}. It will follow by \cref{thm_suff_condition_for_dual_not_sigma_compact} that the families $\dthick^* = \dsyndetic$ and $\dcthick^* = \dcsyndetic$ are not $\sigma$-compact.
    
    To see that \eqref{eqn_suff_condition_for_key_prop_one} holds for $\dthick$ and $\dcthick$, let $B \in \dthick$ or $B \in \dcthick$. By \cref{thm_dy_c_thick_is_ps}, the set $B$ is piecewise syndetic. Thus, by \cref{lem:piecewise_syndetic_shift}, there exists $n \in \N$ for which the set $B \cap (B-n)$ is piecewise syndetic and hence infinite. 
    
    
    We will show that property \eqref{eqn_suff_condition_for_key_prop_two} in \cref{lemma_suff_condition_for_key_property_for_dual_not_sigma_compact} holds for the family $\dthick$. Since $\dthick \subseteq \dcthick$, it will follow that $\dcthick$ also satisfies \eqref{eqn_suff_condition_for_key_prop_two}. Let $\{T_{i,j} \ | \ i, j \in \N\}$ be a doubly-indexed family of thick subsets of $\N$ that are well separated, in the sense that
    \begin{align}
    \label{eqn_very_well_separated_thick_sets}
        \lim_{d \to \infty} \min \bigcup_{\substack{ i, j, k, \ell \in \N \\ (i,j) \neq (k, \ell) \\ \max(i,j,k,\ell) > d}} \min \big(T_{i,j} - T_{k,\ell} \big) = \infty,
    \end{align}
    where the notation $(i,j)$ denotes an element of $\N^2$.
    Let $p_1, p_2, \ldots$ be an increasing enumeration of the primes.  For $i \in \N$, define
    \[B_i = \bigcup_{j = i}^\infty \big(T_{i,j} \cap (p_j \N + 1) \big).\]
    The sets $B_1$, $B_2$, \dots are disjoint, and by \cref{rem:combinatorial_primes}, each is dynamically thick.  To verify the limit in \eqref{eqn_suff_condition_for_key_prop_two}, note that for all $i, k \in \N$,
    \[B_i - B_k = \bigcup_{\substack{j=i\\ \ell = k}}^\infty \Big ( \big( T_{i,j} \cap (p_j \N + 1) \big) - \big( T_{k,\ell} \cap (p_\ell \N + 1) \big) \Big) \subseteq \bigcup_{\substack{j=i\\ \ell = k}}^\infty \big( T_{i,j} - T_{k,\ell} \big).\]
    For all $d \in \N$, considering the cases $i=k$ and $i \neq k$ separately, we see that
    \[\min \bigcup_{\substack{i,k = 1 \\ \max(i,k) > d}}^\infty \big( B_i - B_k\big) \geq \min \bigg( \bigcup_{\substack{ i, j, k, \ell \in \N \\ (i,j) \neq (k, \ell) \\ \max(i,j,k,\ell) > d}} \min \big(T_{i,j} - T_{k,\ell} \big), p_{d} \bigg).\]
    The limit in \eqref{eqn_suff_condition_for_key_prop_two} follows now from the limit in \eqref{eqn_very_well_separated_thick_sets} and the fact that $\lim_{d \to \infty} p_d = \infty$, as desired.
\end{proof}

\begin{remark}
Using \cref{lemma_suff_condition_for_key_property_for_dual_not_sigma_compact}, one can show that the families 
\[
    \IP, \IP_0, \central, \PS, \dthick, \dcthick, \thick, \IP^*, \central^*, \PS^*
\]
are not $\sigma$-compact. However, members of the families $\dsyndetic = \dthick^*$, $\dcsyndetic = \dcthick^*$, $\syndetic = \thick^*$, $\IP = (\IP^*)^*$, $\central = (\central^*)^*$, $\PS = (\PS^*)^*$ can be partitioned into two sets belonging to the same family. This demonstrates that $\sigma$-compactness of $\family^*$ is sufficient, but not necessary, for every member of $\family$ to be partitioned into two members of $\family$.
\end{remark}



\section{Part III: Dynamically piecewise syndetic sets}
\label{sec_dps_sets}

In this section, in analogy to the family of piecewise syndetic sets, we introduce and study the families of dynamically (central) piecewise syndetic sets. The main results describe the relationships these families have with other established families (\cref{thm:relations_C_dPS_PS-intro}) and show that dynamically central piecewise syndetic sets are sets of polynomial multiple measurable recurrence for commuting transformations (\cref{thm:dcps_multiple_recurrence}).


\subsection{Definitions and first results}
\label{sec_dcps_def_and_first_results}

Recall that $\dsyndetic$ and $\dcsyndetic$ denote the families of dynamically syndetic and dynamically central syndetic subsets of $\N$.  Their duals, the families of dynamically thick sets and sets of pointwise recurrence, are denoted $\dthick$ and $\dcthick$.

\begin{definition}
    Let $A \subseteq \N$.
    \begin{enumerate}
        \item The set $A$ is \emph{dynamically piecewise syndetic} if there exists a dynamically syndetic set $B \subseteq \N$ and a dynamically thick set $C \subseteq \N$ such that $A = B \cap C$.  The family of dynamically piecewise syndetic sets is
        \[\dPS \defeq \dsyndetic \classcap \dthick.\]

        \item The set $A$ is \emph{dynamically central piecewise syndetic} if there exists a dynamically central syndetic set $B \subseteq \N$ and a set of pointwise recurrence $C \subseteq \N$ such that $A = B \cap C$.  The family of dynamically central piecewise syndetic sets is
        \[\dcPS \defeq \dcsyndetic \classcap \dcthick.\]
    \end{enumerate} 
\end{definition}

Examples of dynamically (central) piecewise syndetic sets can be gotten by intersecting the examples of dynamically (central) syndetic and dynamically (central) thick sets given in \cref{sec_dy_synd_def_and_char,sec_dy_thick_def_and_char}, respectively. We will show next that return times of distal points to neighborhoods of themselves provide nice dynamical examples of members of the family $\dcPS^*$.

\begin{example}
\label{example_dcps_star_point_returns}
    Let $(X, T) $ be a system, $x \in X$ be a distal point, and $U \subseteq X$ be an open neighborhood of $x$. We claim that the set $R_T(x,U)$ is a $\dcPS^*$ set.  Indeed, by \cref{lemma_class_cap_is_pr}, it suffices to show that for all $B \in \dcsyndetic$, the set $R_T(x,U) \cap B \in \dcsyndetic$.

    Let $B \in \dcsyndetic$.  There exists a minimal system $(Y,S)$, a point $y \in Y$, and an open set $V \subseteq Y$ containing $y$ such that $R_S(y,V) \subseteq B$.  By \cite[Thm. 9.11]{furstenberg_book_1981}, the point $(x,y) \in X \times Y$ is $T \times S$ uniformly recurrent, whereby the set $R_{T \times S}((x,y), U \times V) = R_T(x,U) \cap R_S(y,V)$ is dynamically central syndetic.  It follows that $R_T(x,U) \cap B \in \dcsyndetic$, as desired.
\end{example}



\subsubsection{Containment and partition regularity}

As $\dcsyndetic \subseteq \dsyndetic$ and $\dcthick \supseteq \dthick$, the relationship between the families $\dcPS$ an $\dPS$ is, apriori, not clear.  We show in the following lemma that $\dcPS \subseteq \dPS$, but the proof uses \cref{thm_ps_is_dps,thm_dcps_and_dps_are_ps}, which ultimately rely on the main result in this paper, \cref{thm_dcs_iff_cs}. It would be of interest to find a simpler proof of \cref{lemma_basic_results}. (A note to the astute reader concerned about circular logic.  Prior to \cref{thm_ps_is_dps,thm_dcps_and_dps_are_ps}, on which \cref{lemma_basic_results} depends, we use \cref{lemma_basic_results} only in the proof of \eqref{eqn_relationship_between_dps_and_dcps_classes} in \cref{lemma_translates_of_dps_sets}. \cref{thm_ps_is_dps,thm_dcps_and_dps_are_ps} depend only on item \eqref{item_all_translates_are_dps} in \cref{lemma_translates_of_dps_sets}, and hence does not rely on \cref{lemma_basic_results}.)

\begin{theorem}
\label{lemma_basic_results}
    Dynamically central piecewise syndetic subsets are dynamically piecewise syndetic, that is, $\dcPS \subseteq \dPS$.
\end{theorem}

\begin{proof}
    It follows by combining \cref{thm_ps_is_dps,thm_dcps_and_dps_are_ps} below that $\dcPS \subseteq \PS \subseteq \dPS$, as was to be shown.
\end{proof}

\begin{lemma}
\label{lemma_dcps_dps_are_pr}
    The families $\dcPS$ and $\dPS$ are partition regular, and their duals, $\dcPS^*$ and $\dPS^*$, are filters.
\end{lemma}

\begin{proof}
    That the families $\dcPS$ and $\dPS$ are partition regular follows immediately from the definition and \cref{lemma_class_cap_is_pr}.  It follows then from the discussion in \cref{sec_filter_pr_and_ufs} that the families $\dcPS^*$ and $\dPS^*$ are filters.
\end{proof}



\subsubsection{Translates and dilates of dynamically piecewise syndetic sets}

The following lemmas clarify how members of the families $\dcPS$ and $\dPS$ behave under translations and dilations.  Point \eqref{item_some_translates_are_dcps} is upgraded in \cref{thm:relations_translates_of_dps_to_central} below.

\begin{lemma}
\label{lemma_translates_of_dps_sets}
    Let $A \subseteq \N$ be dynamically piecewise syndetic.
    \begin{enumerate}        
        \item
        \label{item_all_translates_are_dps}
        For all $n \in \N$, the set $A-n$ is dynamically piecewise syndetic.
        
        \item
        \label{item_some_translates_are_dcps}
        The set
        \begin{align}
            \label{eqn_translates_of_dps_set}
            \big\{ n \in \N \ \big| \ A-n \text{ is dynamically central piecewise syndetic} \big\}
        \end{align}
        is dynamically syndetic.
    \end{enumerate}
    Moreover, points \eqref{item_all_translates_are_dps} and \eqref{item_some_translates_are_dcps} hold with $A-n$ replaced by $A+n$. Finally,
    \begin{align}
        \label{eqn_relationship_between_dps_and_dcps_classes}
        \dPS = \bigcup_{n \in \N} \big(\dcPS - n \big) = \bigcup_{n \in \N} \big(\dcPS + n \big).
    \end{align}
\end{lemma}


\begin{proof}
    Because the families $\dsyndetic$ and $\dthick$ are translation invariant (by \cref{lemma_translates_of_dsyndetic_sets,lemma_translates_of_dthick_sets}), it follows by \cref{lemma_idempotent_classcap_is_idempotent} that the family $\dPS = \dsyndetic \classcap \dthick$ is translation invariant, showing \eqref{item_all_translates_are_dps}.

    To see \eqref{item_some_translates_are_dcps}, because $A$ is dynamically piecewise syndetic, there exists a dynamically syndetic set $B \subseteq \N$ and a dynamically thick set $C \subseteq \N$ such that $A = B \cap C$.  By \cref{lemma_translates_of_dsyndetic_sets}, there are dynamically syndetically many $n$'s for which $B-n$ is dynamically central syndetic.  By \cref{lemma_translates_of_dthick_sets}, for all $n \in \N$, the set $C-n$ is dynamically thick.  Therefore, the set in \eqref{eqn_translates_of_dps_set} is dynamically syndetic, as desired. The same argument can be shown abstractly using prior results about dynamically syndetic sets (from \cref{lemma_translates_of_dsyndetic_sets}) and family algebra (from \cref{sec_family_algebra}):
    \[\dPS = \dsyndetic \classcap \dthick \subseteq (\dsyndetic + \dcsyndetic) \classcap (\dsyndetic + \dcthick) \subseteq (\dsyndetic \classcap \dsyndetic) + (\dcsyndetic \classcap \dcthick) = \dsyndetic + \dcPS.\]

    The arguments in the previous two paragraphs combine with the positive-translate results in \cref{lemma_translates_of_dsyndetic_sets,lemma_translates_of_dthick_sets} to show that points \eqref{item_all_translates_are_dps} and \eqref{item_some_translates_are_dcps} hold with $A-n$ replaced by $A+n$. Finally, that \eqref{eqn_relationship_between_dps_and_dcps_classes} holds follows immediately and \eqref{item_all_translates_are_dps} and \eqref{item_some_translates_are_dcps} for $A-n$ and $A+n$ and from the fact that $\dcPS \subseteq \dPS$ (from \cref{lemma_basic_results}).
\end{proof}

The conclusions of \cref{lemma_translates_of_dps_sets} \eqref{item_all_translates_are_dps} and \eqref{item_some_translates_are_dcps} formulated in terms of the family algebra developed in \cref{sec_combinatorics} read: $\dPS \subseteq \{\N\} + \dPS$ and $\dPS \subseteq \dsyndetic + \dcPS$, respectively. In analogy to \cref{lemma_translates_of_dsyndetic_sets} \eqref{item_dcs_translates_of_dcs_are_dcs}, we would like to show that $\dcPS \subseteq \dcsyndetic + \dcPS$, that is, that if $A$ is dynamically central piecewise syndetic, then the set in \eqref{eqn_translates_of_dt_set} is dynamically central syndetic.  We were not able to show this.  It should be easier to determine whether or not the family $\dcPS$ is idempotent: $\dcPS \subseteq \dcPS + \dcPS$. See \cref{quest_is_dcps_idempotent} below for more discussion.


\begin{lemma}
\label{lemma_dilates_of_dps_sets}
    Let $A \subseteq \N$ and $k \in \N$. If $A$ is dynamically central piecewise syndetic, then so are the sets $kA$ and $A/k$. If $A$ is dynamically piecewise syndetic, then so is the set $kA$.
\end{lemma}

\begin{proof}
    Suppose $A \in \dcPS$, and write $A = B \cap C$ where $B \in \dcsyndetic$ and $C \in \dcthick$. We see that $kA = kB \cap kC$ and $A/k = B/k \cap C/k$, both of which are $\dcPS$ sets in view of  \cref{lemma_dilates_of_dsyndetic_sets,lem:dilates_of_dT_dcT}.

    Suppose $A \in \dPS$, and write $A = B \cap C$ where $B \in \dsyndetic$ and $C \in \dthick$.  To see that $kA \in \dPS$, we will write
    \[kA = kB \cap kC = kB \cap \big( kC \cup ( \N \setminus k \N ) \big)\]
    and show that $kB \in \dsyndetic$ and $C' \defeq kC \cup ( \N \setminus k \N ) \in \dthick$.
    
    That $kB \in \dsyndetic$ follows from \cref{lemma_dilates_of_dsyndetic_sets}.  To see that $C' \in \dthick$, let $D \in \dsyndetic$.  If $D \not\subseteq k\N$, then $D \cap C' \neq \emptyset$.  Otherwise, we have that $D \subseteq k \N$.   By \cref{lemma_translates_of_dsyndetic_sets}, there exists $n \in \N$ such that $D - n \in \dcsyndetic$. Since every $\dcsyndetic$ set contains a multiple of $k$ (by \cref{lemma_dilates_of_dsyndetic_sets}), it must be that $n \in k \N$. By \cref{lemma_dilates_of_dsyndetic_sets}, the set $(D - n)/k = D/k - n/k \in \dcsyndetic$. According to \cref{lemma_translates_of_dsyndetic_sets}, the set $D/k \in \dsyndetic$. It follows that $(D/k) \cap C \neq \emptyset$, whereby $D \cap C' \neq \emptyset$.  Since $D \in \dsyndetic$ was arbitrary, it follows that $C' \in \dthick$, as desired.
\end{proof}

It is a consequence of \cref{thm:relations_C_dPS_PS-intro} (proved below independently of \cref{lemma_dilates_of_dps_sets}) that the families of piecewise syndetic sets and dynamically piecewise syndetic sets are the same.  It will follow easily from that that dilates of $\dPS$ sets are $\dPS$ sets, offering another proof of the latter statement in \cref{lemma_dilates_of_dps_sets}.

\subsection{Filter algebra for central and piecewise syndetic family duals}
\label{sec_filter_algebra_for_central_duals}

In this section, we draw on \cref{cor_central_syndetic_is_strongly_central,thm_dcs_iff_cs} to prove some abstract results concerning special classes of filters. These results ultimately will allow us to relate the families $\dcPS$ and $\dPS$ to themselves and to the families $\central$ and $\PS$ in the next section.




\begin{theorem}
\label{prop:general_abstract_intersection_filter_nonsense}
    If $\filter$ is a syndetic, idempotent filter and $\filtertwo$ is a $\central^*$, idempotent filter, then $\filter \familycap \filtertwo$ is a syndetic, idempotent filter.
\end{theorem}
\begin{proof}
    Since $\filter$ and $\filtertwo$ are filters, the family $\filter \familycap \filtertwo$ is a filter.  Thus, to see that it is a syndetic filter, it suffices to show that for all $A \in \filter$ and $B \in \filtertwo$, the set $A \cap B$ is syndetic.  Let $A \in \filter$, $B \in \filtertwo$, and $H \subseteq \N$ be thick.  Since the set $A$ is central syndetic, by \cref{cor_central_syndetic_is_strongly_central}, the set $A \cap H$ is central. Since $B \in \central^*$, the set $A \cap H \cap B$ is nonempty.  Because $H$ was an arbitrary thick set, we see that the set $A \cap B$ is syndetic, as desired. Finally, that $\filter \classcap \filtertwo$ is idempotent follows immediately from \cref{lemma_idempotent_classcap_is_idempotent} and the idempotency of $\filter$ and $\filtertwo$.
\end{proof}
   

\begin{theorem}
\label{cor_central_class_equalities}
\label{cor_shifted_class_equalities}
Let $\familytwo$ be a filter on $\N$.
\begin{enumerate}
    \item
    \label{item_central_star_class_cap}
    If $\filtertwo$ is a $\central^*$, idempotent filter, then $\filtertwo \subseteq \dcPS^*$.

    \item
    \label{item_syndetic_trans_inv_star_class_cap}
    If $\filtertwo$ is a syndetic, translation-invariant filter, then $\filtertwo \subseteq \dPS^*$.
\end{enumerate}
\end{theorem}

\begin{proof}
    \eqref{item_central_star_class_cap} \ Suppose that $\filtertwo$ is a $\central^*$, idempotent filter. To see that $\filtertwo \subseteq \dcPS^*$, it suffices by \cref{lemma_class_cap_is_pr} to show that for all $A \in \dcsyndetic$ and $B \in \filtertwo$, the set $A \cap B \in \dcsyndetic$.  Let $A \in \dcsyndetic$ and $B \in \filtertwo$. By \cref{thm_dcs_iff_cs}, the set $A$ belongs to a syndetic, idempotent filter $\filter$.  By \cref{prop:general_abstract_intersection_filter_nonsense}, the family $\filter \classcap \filtertwo$ is a syndetic, idempotent filter.  Since $A \cap B \in \filter \classcap \filtertwo$, we have by \cref{thm_dcs_iff_cs} that $A \cap B \in \dcsyndetic$, as desired.

    

    \eqref{item_syndetic_trans_inv_star_class_cap} \ Suppose that $\filtertwo$ is a syndetic, translation-invariant filter. By \cref{lemma_condition_on_subfamily_of_ps_star}, the family $\PS^*$ contains all syndetic, translation-invariant filters, so we have that $\filtertwo \subseteq \PS^* \subseteq \central^*$.  Note that $\filtertwo$ is idempotent because it is translation invariant.
   

    To see that $\filtertwo \subseteq \dPS^*$, it suffices by \cref{lemma_class_cap_is_pr} to show that for all $A \in \dsyndetic$ and $B \in \filtertwo$, the set $A \cap B \in \dsyndetic$.  Let $A \in \dsyndetic$ and $B \in \filtertwo$. By \cref{lemma_translates_of_dsyndetic_sets}, there exists $n \in \N$ such that $A - n \in \dcsyndetic$. Since $B - n \in \filtertwo - n \subseteq \filtertwo \subseteq \central^*$ and $\filtertwo$ is idempotent, it follows from \eqref{item_central_star_class_cap} and \cref{lemma_class_cap_is_pr} that
    \[
        (A \cap B) - n = (A - n) \cap (B - n) \in \dcsyndetic \classcap \filtertwo = \dcsyndetic.
    \]
    Adding $n$, by \cref{lemma_translates_of_dsyndetic_sets}, the set $A \cap B \in \dsyndetic$, as was to be shown.
\end{proof}


It is tempting to read \cref{cor_central_class_equalities} \eqref{item_central_star_class_cap} as saying that the intersection of a $\central^*$ set with a $\dcsyndetic$ set is a $\dcsyndetic$ set.  This does not follow from \cref{cor_central_class_equalities} \eqref{item_central_star_class_cap} (and is, in fact, false) since not every $\central^*$ set belongs to a $\central^*$, idempotent filter.  Indeed, the set constructed in \cref{rem:combinatorial_primes} is a dynamically thick set that is not $\IP$.  The complement is an example of an $\IP^*$ set, hence $\central^*$ set, that is not $\dsyndetic$.  By \cref{thm_dcs_iff_cs}, this set cannot belong to a $\central^*$, idempotent filter.

It is true, however, that every $\PS^*$ set belongs to a syndetic, translation-invariant filter; indeed, by \cref{lemma_condition_on_subfamily_of_ps_star}, the family $\PS^*$ is itself such a filter.  Therefore, it follows from \cref{cor_central_class_equalities} \eqref{item_syndetic_trans_inv_star_class_cap} and \cref{lemma_class_cap_is_pr} that the intersection of a $\PS^*$ and a $\dcsyndetic$ set is a $\dcsyndetic$ set. We expand on this in \cref{cor:largest_family_satisfies_cap} below.  



\subsubsection{Combinatorial and dynamical piecewise syndetic families and a proof of \cref{thm:relations_C_dPS_PS-intro}}
\label{sec_dcps_relations}

In this section, we explain the relationships between the dynamical families $\dcPS$ and $\dPS$ and the combinatorial families $\central$ and $\PS$. In particular, we will show that
\begin{align}
    \label{eqn_main_containments}
    \central \subsetneq \dcPS \subsetneq \dPS = \PS,
\end{align}
thus completing the proof of \cref{thm:relations_C_dPS_PS-intro} and the containments stated in \cref{fig:containmentdiagram}.

\begin{theorem}
\label{thm_c_is_dcps}
\label{thm_ps_is_dps}
    The family of central sets is properly contained in the family of dynamically central piecewise syndetic sets, ie., $\central \subsetneq \dcPS$, and piecewise syndetic sets are dynamically piecewise syndetic, ie., $\PS \subseteq \dPS$.
\end{theorem}

\begin{proof}
    By \cref{lem:central_characterization} and the fact that $\thick \subseteq \dthick \subseteq \dcthick$, we have that
    \begin{align*}
        \central = \dcsyndetic \classcap \thick &\subseteq \dcsyndetic \classcap \dcthick = \dcPS,\\
        \PS = \dsyndetic \classcap \thick &\subseteq \dsyndetic \classcap \dthick = \dPS.
    \end{align*}
    By \cref{rem:combinatorial_primes}, there is a set $A \in \dthick$ which is not an $\IP$ set. Since every central set is an $\IP$ set, we see that the set $A \in \dcPS$ but $A \not \in \central$. Thus, the containment $\central \subseteq \dcPS$ is proper.
\end{proof}

\begin{remark}
    Let us offer an alternative, combinatorial proof of the containment $\PS \subseteq \dPS$ in \cref{thm_c_is_dcps}.  We will show that $\dPS^*$ is a syndetic, translation-invariant filter and appeal to \cref{lemma_condition_on_subfamily_of_ps_star} to see that $\dPS^* \subseteq \PS^*$, whereby $\PS \subseteq \dPS$.  Translation invariance of $\dPS^*$ follows from the positive translation invariance of $\dPS$ shown in \cref{lemma_translates_of_dps_sets} \eqref{item_all_translates_are_dps}. That $\dPS^*$ is a filter is shown in \cref{lemma_dcps_dps_are_pr}.  That every member of $\dPS^*$ is syndetic follows from the fact that $\thick \subseteq \dthick \subseteq \dPS$, whereby the family $\dPS^*$ is syndetic.
\end{remark}



The following theorem, when combined with \cref{thm_c_is_dcps}, completes the proof of the containments in \eqref{eqn_main_containments} and hence the proof of \cref{thm:relations_C_dPS_PS-intro}.  It is interesting to note that the containment $\dPS \subseteq \PS$ follows immediately from \cref{cor_weak_structure_result}.  Note, however, that the containment $\dPS \subseteq \PS$ was used to prove \cref{cor_weak_structure_result}, so it is best to see \cref{cor_weak_structure_result} as an enhancement to the containment $\dPS \subseteq \PS$.

\begin{theorem}
\label{thm_dcps_and_dps_are_ps}
    The family of dynamically (central) piecewise syndetic sets is (properly) contained the family of piecewise syndetic sets, ie., $\dcPS \subsetneq \PS$ and $\dPS \subseteq \PS$.
\end{theorem}

\begin{proof}
    By \cref{lemma_condition_on_subfamily_of_ps_star}, the family $\PS^*$ is a syndetic, translation-invariant filter.  Since $\PS^* \subseteq \central^*$ and translation-invariance implies idempotency, we have that $\PS^*$ is also a $\central^*$, idempotent filter.  By \cref{cor_shifted_class_equalities}, we see that $\PS^* \subseteq \dcPS^*$ and $\PS^* \subseteq \dPS^*$.  It follows that $\dcPS \subseteq \PS$ and $\dPS \subseteq \PS$.

    To see that $\dcPS \subsetneq \PS$, note that it follows from \cref{example_dcps_star_point_returns} that the set $2\N$ is $\dcPS^*$.  Therefore, the set $2\N - 1$, which is $\PS$, is not $\dcPS$.
\end{proof}



\begin{corollary}
\label{thm:relations_translates_of_dps_to_central}
    For all dynamically piecewise syndetic sets $A \subseteq \N$, the set
    \begin{align}
        \label{eqn_translates_of_dps_set_to_central}
        \big\{ n \in \N \ \big| \ A-n \text{ is central} \big\}
    \end{align}
    is dynamically syndetic, that is, $\dPS \subseteq \dsyndetic + \central$. The same statement holds with $A-n$ replaced by $A+n$. Finally,
    \begin{align}
        \label{eqn_relationship_between_dps_and_c_classes}
        \dPS = \bigcup_{n \in \N} (\central - n) = \bigcup_{n \in \N} (\central + n).
    \end{align}
        
\end{corollary}

\begin{proof}
    Let $A \in \dPS$.  By \cref{thm_dcps_and_dps_are_ps}, the set $A$ is piecewise syndetic.  By \cref{lem:central_characterization}, there exists a minimal system $(X,T)$, a point $x \in X$, a nonempty, open set $U \subseteq X$, and a thick set $H \subseteq \N$ such that $A \supseteq R(x,U) \cap H$.  For all $n \in R(x,U)$, we see that the set $A - n \supseteq R(T^n x, U) \cap (H-n)$, which, by \cref{lem:central_characterization}, is central.  Thus, the set in \eqref{eqn_translates_of_dps_set_to_central} contains $R(x,U) \in \dsyndetic$, as was to be shown.

    To arrive at the same result when $A-n$ is replaced by $A+n$, begin by writing $A \supseteq R(x,U) \cap H$ in the same way, noting that by \cref{lemma_may_assume_invertible}, we may assume that the system $(X,T)$ is invertible. We claim that the set $R_{T^{-1}}(x,U)$ is contained in the set in \eqref{eqn_translates_of_dps_set_to_central} with $A-n$ replaced by $A+n$.  Since the system $(X,T^{-1})$ is minimal (see, for example, \cite[Lemma 2.7]{glasscock_koutsogiannis_richter_2019}), this will demonstrate that the set in \eqref{eqn_translates_of_dps_set_to_central} is dynamically syndetic.
    
    Suppose $n \in R_{T^{-1}}(x,U)$ so that $T^{-n} x \in U$.  It is quick to check that $R(T^{-n}x,U) \cap \{n+1, n+2, \ldots\} \subseteq R(x,U) + n$.  Since $R(T^{-n}x,U) \in \dcsyndetic$, it follows by \cref{lemma_ds_dcs_modifyable_on_a_finite_set} that the set $R(T^{-n}x,U) \cap \{n+1, n+2, \ldots\} \in \dcsyndetic$. It follows by \cref{lem:central_characterization} that the set
    \[A + n \supseteq \big(R(x,U) + n\big) \cap \big(H + n\big) \supseteq \big(R(T^{-n}x,U) \cap \{n+1, n+2, \ldots\}\big) \cap \big(H + n\big),\]
    is central, as desired.

    Finally, that \eqref{eqn_relationship_between_dps_and_c_classes} holds follows immediately from \eqref{eqn_translates_of_dps_set_to_central} for $A-n$ and $A+n$.
\end{proof}


\begin{theorem}
\label{cor:largest_family_satisfies_cap}
    Let $A \subseteq \N$.  The following conditions are equivalent to the set $A$ being thickly syndetic ($A \in \PS^*$):
    \begin{enumerate}
        \item
        for all $B \in \dsyndetic$, the set $A \cap B \in \dsyndetic$;

        \item
        for all $B \in \dthick$, the set $A \cap B \in \dthick$;

        \item
        for all $B \in \PS$, the set $A \cap B \in \PS$;

        \item for all $B \in \PS^*$, the set $A \cap B \in \PS^*$.
    \end{enumerate}
    Thus, the family $\PS^*$ is the largest family of sets that satisfies any one condition (equivalently, all conditions).  Moreover, if the set $A$ is thickly syndetic, then
    \begin{enumerate}
    \setcounter{enumi}{4}
        \item
        for all $B \in \dcsyndetic$, the set $A \cap B \in \dcsyndetic$;

        \item
        for all $B \in \dcthick$, the set $A \cap B \in \dcthick$;

        \item
        for all $B \in \dcPS$, the set $A \cap B \in \dcPS$;

        \item for all $B \in \dcPS^*$, the set $A \cap B \in \dcPS^*$.
    \end{enumerate}
\end{theorem}
\begin{proof}
    Combining \cref{thm_ps_is_dps,thm_dcps_and_dps_are_ps}, we have that $\PS^* = \dPS^* \subseteq \dcPS^*$.  That the first four statements are equivalent to $A$ being a member of $\PS^*$ then follows immediately from \cref{lemma_class_cap_is_pr}.  If $A \in \PS^*$, then $A \in \dcPS^*$.  The final four statements follow since, by \cref{lemma_class_cap_is_pr}, they are all equivalent to $A$ being a member of $\dcPS^*$.
\end{proof}



\subsection{Polynomial set recurrence and Brauer configurations: a proof of \texorpdfstring{\cref{thm:dcps_multiple_recurrence}}{Theorem H}}
\label{sec_dcps_and_set_recurrence}


\subsubsection{Polynomial set recurrence: a proof of \texorpdfstring{\cref{thm:dcps_multiple_recurrence}}{Theorem H}}

In this section, we will show that -- in a very general sense made precise by \cref{thm:dcps_multiple_recurrence} -- dynamically central piecewise syndetic sets are sets of recurrence in ergodic theory.  In what follows, a \emph{commuting probability measure preserving system} $(X, \mu, T_1, \ldots, T_k)$ is a tuple consisting of a probability measure space $(X, \mu)$ together with commuting measure preserving transformations $T_i: X \to X$, $i = 1, \ldots, k$.  The system is \emph{invertible} if the maps $T_i$, $i = 1, \ldots, k$, are invertible.

Denote by $\filtertwo$ the upward closure of the family of all subsets of $\N$ of the form
\begin{align}
    \label{eqn_poly_return_set}
    \polyret \big(\mu, T_i, p_{i,j}, E \big)_{\substack{i = 1, \ldots, k \\ j = 1, \ldots, \ell}} \defeq \left\{ n \in \N \ \middle | \ \mu\left( \bigcap_{j=1}^{\ell} \left( \prod_{i=1}^{k} T_i^{p_{i,j}(n)} \right)^{-1} E \right) > 0 \right\},
\end{align}
where $k, \ell \in \N$, the tuple $(X, \mu, T_1, \ldots, T_k)$ is an invertible, commuting probability measure preserving system, the set $E \subseteq X$ satisfies $\mu(E) > 0$, and, for $1 \leq i \leq k$, $1 \leq j \leq \ell$, the polynomial $p_{i, j} \in \Q[n]$ satisfies $p_{i,j}(\Z) \subseteq \Z$ and $p_{i,j}(0) = 0$.  Thus, the family $\filtertwo$ consists of times of returns of any set of positive measure in any probability measure space under polynomial iterates of a finite collection of invertible, commuting measure preserving transformations.  Members of the dual family, $\familytwo^*$, are called \emph{sets of polynomial multiple measurable recurrence for invertible commuting transformations}.



\begin{lemma}
\label{lemma_localization_is_syndetic_idempotent_filter}
    The family $\filtertwo$ is an $\IP^*$, idempotent filter.
\end{lemma}

\begin{proof}
    By a corollary of the IP Polynomial \Szemeredi{} Theorem of Bergelson and McCutcheon \cite[Thm. 7.12]{Bergelson-McCutcheon-IPpolynomialSzemeredi}, every member of $\filtertwo$ is an $\IP^*$ set.  Therefore, the family $\filtertwo$ is $\IP^*$.

    
    To show that $\filtertwo$ is a filter, we must show that the set
    \begin{align}
        \label{eqn_set_to_show_is_filter}
        \polyret \big(\mu, T_i, p_{i,j}, E \big)_{\substack{i = 1, \ldots, k \\ j = 1, \ldots, \ell}} \cap \polyret \big(\nu, S_i, q_{i,j}, F \big)_{\substack{i = 1, \ldots, r \\ j = 1, \ldots, t}}
    \end{align}
    belongs to $\filtertwo$, where $(X, \mu, T_1, \ldots, T_k)$, $p_{i, j} \in \Q[n]$, $E \subseteq X$, $(Y, \nu, S_1, \ldots, S_r)$, $q_{i, j} \in \Q[n]$, and $F \subseteq Y$ are as described above.  Denote by $I$ the identity transformation on both $X$ and $Y$, and consider the commuting probability measure preserving system
    \[\big(X \times Y, \mu \otimes \nu, T_1 \otimes I, \ldots, T_k \otimes I, I \otimes S_1, \ldots, I \otimes S_r \big).\]
    The set $E \times F \subseteq X \times Y$ satisfies $(\mu \otimes \nu)(E \times F) = \mu(E) \nu(F) > 0$. Moreover,
    \begin{gather*}
        \big(\mu \otimes \nu \big) \left( \bigcap_{j_1=1}^{\ell} \bigcap_{j_2=1}^{t} \left( \prod_{i=1}^{k} (T_i \otimes I)^{p_{i,j_1}(n)} \ \cdot \ \prod_{i=1}^{r} (I \otimes S_i)^{q_{i,j_2}(n)}\right)^{-1} (E \times F) \right) = \\
        \big( \mu \otimes \nu \big) \left( \bigcap_{j=1}^{\ell} \left( \prod_{i=1}^{k} (T_i \otimes I)^{p_{i,j}(n)} \right)^{-1} (E \times F) \ \cap \ \bigcap_{j=1}^{t} \left( \prod_{i=1}^{r} (I \otimes S_i)^{q_{i,j}(n)} \right)^{-1} (E \times F) \right) = \\
        \mu\left( \bigcap_{j=1}^{\ell} \left( \prod_{i=1}^{k} T_i^{p_{i,j}(n)} \right)^{-1} E \right) \ \cdot \ \nu\left( \bigcap_{j=1}^{t} \left( \prod_{i=1}^{r} S_i^{q_{i,j}(n)} \right)^{-1} F \right).
    \end{gather*}
    It follows that the set in \eqref{eqn_set_to_show_is_filter} is equal to
    \[\polyret \big( \mu \otimes \nu, U_i, m_{i,j}, E \times F \big)_{\substack{i = 1, \ldots, k+r \\ j = 1, \ldots, \ell t}},\]
    where
    \[U_i \defeq \begin{cases} T_i \otimes I & \text{ if $i \leq k$} \\ I \otimes S_{i-k} & \text{ if $i > k$}\end{cases}, \qquad \text{ and } \qquad m_{i,j} \defeq \begin{cases} p_{i, (j \text{ mod } \ell) + 1} & \text{ if $i \leq k$} \\ q_{i-k, \lceil j / \ell \rceil} & \text{ if $i > k$}\end{cases}.\]
    This shows that the set in \eqref{eqn_set_to_show_is_filter} belongs to $\filtertwo$, as desired.
    
    
    
    Next, we will show that $\filtertwo$ is idempotent. It suffices by \cref{rmk_showing_idempotency} to show that if $D$ is a set of the form in \eqref{eqn_poly_return_set}, then $D \subseteq D - \filtertwo$, that is, for all $m \in D$, the set $D-m \in \filtertwo$.
    
    Let $D$ be a set of the form in \eqref{eqn_poly_return_set}, and let $m \in D$.
    Define
    \[
        E' \defeq \bigcap_{j=1}^{\ell} \left( \prod_{i=1}^{k} T_i^{p_{i,j}(m)} \right)^{-1} E.
    \]
    Since $m \in D$, we have that $\mu(E') > 0$. For $i = 1, \ldots, k$ and $j = 1, \ldots, \ell$, define $q_{i,j} \in \Q[n]$ by
    \[
        q_{i,j}(n) \defeq p_{i,j}(n + m) - p_{i,j}(m).
    \]
    Note that $q_{i,j}(\Z) \subseteq \Z$ and $q_{i,j}(0) = 0$.

    By expanding the definition of $E'$, we get
    \begin{align*}
        \bigcap_{j=1}^{\ell} \left( \prod_{i=1}^{k} T_i^{q_{i,j}(n)} \right)^{-1} E' &= \bigcap_{j=1}^{\ell} \left( \prod_{i=1}^{k} T_i^{p_{i,j}(n+m) - p_{i,j}(m)} \right)^{-1} E' \\
        &= \bigcap_{j=1}^{\ell} \left( \prod_{i=1}^{k} T_i^{p_{i,j}(n+m) - p_{i,j}(m)} \right)^{-1} \left( \bigcap_{j'=1}^{\ell} \left( \prod_{i'=1}^{k} T_{i'}^{p_{i',j'}(m)} \right)^{-1} E \right) \\
        &\subseteq \bigcap_{j=1}^{\ell} \left( \prod_{i=1}^{k} T_i^{p_{i,j}(n+m)} \right)^{-1} E.
    \end{align*}
    It follows that
    \[
        \polyret \big(\mu, T_i, q_{i,j}, E' \big)_{\substack{i = 1, \ldots, k \\ j = 1, \ldots, \ell}} \subseteq \left\{n \in \N \ \middle | \ \mu \left( \bigcap_{j=1}^{\ell} \left( \prod_{i=1}^{k} T_i^{p_{i,j}(n+m)} \right)^{-1} E \right) > 0 \right\} = D - m,
    \]
    whereby $D-m \in \filtertwo$, as desired.
\end{proof}

\begin{theorem}
\label{thm_dcps_is_poly_mult_comm_rec}
    Every $\dcPS$ set is a $\filtertwo^*$ set, that is, dynamically central piecewise syndetic sets are sets of polynomial multiple measurable recurrence for invertible commuting transformations.
\end{theorem}

\begin{proof}
    By \cref{lemma_localization_is_syndetic_idempotent_filter}, the family $\familytwo$ is a $\central^*$, idempotent filter.  Therefore, by \cref{cor_central_class_equalities}, we have that $\familytwo \subseteq \dcPS^*$.  Taking the dual, we see $\dcPS \subseteq \familytwo^*$, as desired.
\end{proof}

To prove \cref{thm:dcps_multiple_recurrence} from the introduction, we need only to remove the assumption on invertibility in \cref{thm_dcps_is_poly_mult_comm_rec}.

\begin{proof}[Proof of \cref{thm:dcps_multiple_recurrence}]
    Denote by $\filter$ the upward closure of the family of all subsets of $\N$ of the form in \eqref{eqn_poly_return_set}, where $k, \ell \in \N$, the tuple $(X, \mu, T_1, \ldots, T_k)$ is a commuting probability measure preserving system, the set $E \subseteq X$ satisfies $\mu(E) > 0$, and, for $1 \leq i \leq k$, $1 \leq j \leq \ell$, the polynomial $p_{i, j} \in \Q[n]$ satisfies $p_{i,j}(\N_0) \subseteq \N_0$ and $p_{i,j}(0) = 0$. To prove \cref{thm:dcps_multiple_recurrence}, we must show that $\dcPS \subseteq \family^*$.
    
    Let $A \in \dcPS$ and let $B \in \family$.  The set $B$ contains a set of the form $\polyret \big(\mu, T_i, p_{i,j}, E \big)$ as described in the previous paragraph.  By the measurable natural extension (cf. \cite[Lemma 7.11]{Bergelson-McCutcheon-IPpolynomialSzemeredi}), there exists an invertible, commuting probability measure preserving system $(Y, \nu, S_1, \ldots, S_k)$ and a measurable set $F \subseteq Y$ such that for all $n_{i, j} \in \N_0$, $1 \leq i \leq k$, $1 \leq j \leq \ell$,
    \begin{align}
        \label{eqn_measurable_natural_extension}
        \nu \left( \bigcap_{j=1}^\ell \left( \prod_{i=1}^{k} S_i^{n_{i,j}} \right)^{-1} F \right) = \mu \left( \bigcap_{j=1}^\ell \left( \prod_{i=1}^{k} T_i^{n_{i,j}} \right)^{-1} E \right).
    \end{align}
    By \cref{thm_dcps_is_poly_mult_comm_rec}, there exists $n \in A$ such that
    \[\nu\left( \bigcap_{j=1}^{\ell} \left( \prod_{i=1}^{k} S_i^{p_{i,j}(n)} \right)^{-1} F \right) > 0.\]
    It follows from \eqref{eqn_measurable_natural_extension} that $n \in \polyret \big(\mu, T_i, p_{i,j}, E \big)$, showing that $A \cap B \neq \emptyset$, as desired.
\end{proof}


\begin{remark}
The $\IP$ polynomial \Szemeredi{} theorem \cite[Thm 7.12]{Bergelson-McCutcheon-IPpolynomialSzemeredi} shows that $\filtertwo$ is an $\IP^*$ family, while \cref{thm_dcps_is_poly_mult_comm_rec} shows that $\filtertwo$ is an $\dcPS^*$ family. Because there is no containment in either direction between the families $\IP^*$ and $\dcPS^*$, \cref{thm_dcps_is_poly_mult_comm_rec} is neither an upgrade nor an immediate corollary of the $\IP$ polynomial \Szemeredi{} theorem.

An upgrade to the conclusion of the $\IP$ polynomial \Szemeredi{} theorem is available in the case that the polynomials in \eqref{eqn_poly_return_set} are required to be linear.  In that case, the family $\familytwo$ can be shown to be $\IP_0^*$ (see \cref{sec_open_quests_dps_sets} for a definition of the family $\IP_0^*$) by a corollary of a result of Furstenberg and Katznelson \cite[Thm. 10.3]{furstenberg_katznelson_1985}.  The relationship between the families $\dcPS^*$ and $\IP_0^*$ is not known; it may be the case that $\IP_0^* \subseteq \dcPS^*$, for example.  See \cref{ques:pointwise_recurrence_IP_0} below for more discussion in this direction.
\end{remark}

\subsubsection{Dynamically central piecewise syndetic sets contain polynomial Brauer configurations}

Central sets are known to contain an abundance of combinatorial configurations (cf. \cite[Ch. 14]{hindman_strauss_book_2012}). In this section, we demonstrate that dynamically central piecewise syndetic sets also contain some combinatorial configurations.  (Recall from \cref{thm_c_is_dcps} that central sets are dynamically central piecewise syndetic.)  We speculate in \cref{sec_open_quests_dps_sets} about the extent to which $\dcPS$ sets exhibit further combinatorial richness.

\begin{theorem}
\label{thm:Brauer_and_dcPS}
    Let $A \subseteq \N$ be dynamically central piecewise syndetic. For all $k \in \N$ and all $p_1, p_2, \ldots, p_k \in \Q[x]$ with $p_i(\Z) \subseteq \Z$ and $p_i(0) = 0$, there exist $x, y \in \N$ such that
    \begin{align}
        \label{eqn_poly_brauer_in_dcps}
        x, y, x + p_1(y), \ldots, x + p_k(y) \in A.
    \end{align}
\end{theorem}

\begin{proof}
    By \cref{thm_dcps_and_dps_are_ps}, the set $A$ is piecewise syndetic and hence has positive upper Banach density, $d^*(A) > 0$. By a standard application of the Furstenberg Correspondence Principle (\cite[Thm. 1.1]{Furstenberg-ErgodicBehavior}, see also \cite[Thm. 3.2.5]{mccutcheon_1999}), there exists an invertible probability measure preserving system $(X, \mu, T)$ and a set $E \subseteq X$ with $\mu(E) \geq d^*(A) > 0$ such that for all $n_1, \ldots, n_k \in \Z$,
    \begin{align*}
        % \label{eqn_correspondence_result}
        \mu \big(T^{-n_1} E \cap \cdots \cap T^{-n_k} E \big) \leq d^* \big( (A-n_1) \cap \cdots \cap (A-n_k) \big).
    \end{align*}
    Setting $p_0 \equiv 0$, by \cref{thm_dcps_is_poly_mult_comm_rec}, there exists $y \in A$ such that
    \[
        0 < \mu \big(E \cap T^{-p_1(y)} E \cap \cdots \cap T^{-p_k(y)} E \big) \leq d^* \big( A \cap (A-p_1(y)) \cap \cdots \cap (A-p_k(y)) \big).
    \]
    Therefore, there exist (many) $x \in A$ for which \eqref{eqn_poly_brauer_in_dcps} holds, as was to be shown.
\end{proof}

\section{Open questions}
\label{sec:open_questions}

We collect below a number of open questions and directions suggested by the results in this paper. They are roughly categorized as concerning the families of dynamically syndetic, thick, and piecewise syndetic sets, with the exception of the following.

\begin{question}
\label{quest_generalizations_to_other_grps}
    To which other group and semigroup actions do the main results in this paper generalize?  For example, a generalization of \cref{thm_dcs_iff_cs} to $\Z$ might read: a subset of $\Z$ is dynamically central syndetic if and only if it is a member of a syndetic, idempotent filter on $\Z$.
\end{question}

Before attempting to answer \cref{quest_generalizations_to_other_grps}, it would be necessary to decide how the family of dynamically central syndetic subsets of $\Z$ is defined.  We discuss this at more length in point \eqref{item_generalization_to_z} at the beginning of \cref{sec_dsyndetic}.

\subsection{Dynamically syndetic sets}

By definition, a subset of $\N$ is dynamically central syndetic if it contains a set of the form $R(x,U)$, where $(X,T)$ is a minimal system, $x \in X$, and $U \subseteq X$ is an open neighborhood of $x$.  It can happen, however, that a set of the form $R(x,U)$ where $x \not \in \overline{U}$ is dynamically central syndetic.  Recall example \eqref{item_dcs_two_adic_example} in \cref{sec_dy_synd_def_and_char}: the set $A$ of positive integers with even 2-adic valuation is dynamically central syndetic because the point $1_{A \cup \{0\}}$ is uniformly recurrent in the full shift $(\{0,1\}^{\N_0},\sigma)$ and $A = R_\sigma(1_{A \cup \{0\}}, [1]_0)$ where $1_{A \cup \{0\}} \in [1]_0$. Since $\{0,1\}^{\N_0}$ is the disjoint union of the cylinders $[0]_0$ and $[1]_0$, by properties of the 2-adic valuation, we see that
\[A / 2 = \N \setminus A = R_\sigma(1_{A \cup \{0\}}, [0]_0).\]
By \cref{lemma_dilates_of_dsyndetic_sets}, this set is dynamically central syndetic despite the fact that $1_{A \cup \{0\}} \not\in [0]_0$.  This leads us naturally to the following question.


\begin{question}
    \label{quest_nec_and_suff_for_rxu_dcs}
    Let $(X,T)$ be a minimal system, $x \in X$, and $U \subseteq X$ be nonempty and open.  Give necessary and sufficient conditions for $x$ and $U$ so that the set $R(x,U)$ is dynamically central syndetic.
\end{question}

An answer to \cref{quest_nec_and_suff_for_rxu_dcs} would be useful in addressing several of the questions posed below. There is precedent for the family of central sets.  Indeed, if $U \subseteq X$ is clopen, then the set $R(x,U)$ is central if and only if $U$ contains a point proximal to $x$ (recall the definition of proximal from \cref{sec_ip_and_central_sets}).  It is a fact \cite[Lem. 19.22, Thms. 19.23, 19.25]{hindman_strauss_book_2012} that a pair of points $x, y \in X$ is proximal if and only if there exists an idempotent ultrafilter $p \in \beta \N$ such that $T^p x = y$.  If $R(x,U)$ is central, then it belongs to a minimal, idempotent ultrafilter $p$, and so $T^p x \in \overline{U} = U$.  On the other hand, if $x$ is proximal to a point $y \in U$, then there is a minimal, idempotent ultrafilter $p$ such that $T^p x = y$, and so $R(x,U) \in p$.  It seems plausible that an analogous answer to \cref{quest_nec_and_suff_for_rxu_dcs} could be given in terms of some kind of proximality.\\
  
Our main result, \cref{thm_dcs_iff_cs}, gives that if $A \subseteq \N$ belongs to a syndetic, idempotent filter, then there exists a minimal system $(X,T)$, a point $x \in X$, and an open set $U \subseteq X$ containing $x$ such that $R(x,U) \subseteq A$.  The proof of \cref{thm_dcs_iff_cs}, however, does not seem to give any indication as to the dynamical nature of the system $(X,T)$.  When, for example, can it be guaranteed that $(X,T)$ is totally minimal, weakly mixing, distal, point distal, a nilsystem, or equicontinuous?  Here, for example, is a concrete question along these lines.  Recall the definition of a distal point from \cref{sec_dcps_def_and_first_results}.

\begin{question}
\label{question_cstar_idempo_filter_implies_pt_distal}
    Does every member of an $\IP^*$, idempotent filter contain a set of the form $R(x,U)$, where $(X,T)$ is a minimal system, $x \in X$ is a distal point, and $U \subseteq X$ is an open neighborhood of $x$?
\end{question}

\cref{question_cstar_idempo_filter_implies_pt_distal} is motivated by a result of Furstenberg \cite[Thm 9.11]{furstenberg_book_1981}: a point $x \in X$ is a distal point if and only if for all open neighborhoods $U \subseteq X$ of $x$, the set $R(x,U)$ belongs to $\IP^*$ if and only if for all open neighborhoods $U \subseteq X$ of $x$, the set $R(x,U)$ belongs to $\central^*$.  Note that the family of thickly syndetic sets, $\PS^*$, is a $\central^*$, idempotent filter.  Thus, an ostensibly easier variant of \cref{question_cstar_idempo_filter_implies_pt_distal} is attained by replacing $\IP^*$ with $\PS^*$.

While $\IP$ sets are sets of pointwise recurrence for distal points, recall from \cref{rem:combinatorial_primes} presented in \cref{sec_dthick} that even dynamically thick sets need not be $\IP$.  A positive answer to \cref{question_cstar_idempo_filter_implies_pt_distal} would show, via the following lemma, that a set of pointwise recurrence for distal points must be a translate of an $\IP$ set.

\begin{lemma}
    Suppose the answer to \cref{question_cstar_idempo_filter_implies_pt_distal} is positive. If $A \subseteq \N$ is a set of pointwise recurrence for all distal points (ie., for all systems $(X,T)$, all distal $x \in X$, and all open $U \subseteq X$ containing $x$, the set $A \cap R(x,U) \neq \emptyset$), then there exists $n \in \N$ such that $A - n \in \IP$.
\end{lemma}

\begin{proof}
    Consider the family
    \begin{align*}
        \class &\defeq \big\{ A \subseteq \N \ \big | \ \exists \ \ell \in \N, \ \ A-\ell \in \IP \big\}.
    \end{align*}
    Using the fact that the family $\IP$ is partition regular, it is not hard to show that
    \begin{align*}
        \class^* &= \big\{ B \subseteq \N \ \big | \ \forall \ \ell \in \N, \ B \cap (B-1) \cap \cdots \cap (B-\ell) \in \IP^* \big\}.
    \end{align*}
    The family $\class^*$ is an $\IP^*$, translation-invariant (hence, idempotent) filter.

    Suppose $A \subseteq \N$ is a set of pointwise recurrence for all distal points.  We will use the positive answer to \cref{question_cstar_idempo_filter_implies_pt_distal} to show that $A \in \class$ and conclude the proof of the lemma.

    To show that $A \in \class$, we will show that $A \in (\class^*)^*$.  Let $B \in \class^*$.  We must show that $A \cap B \neq \emptyset$.  By the positive answer to \cref{question_cstar_idempo_filter_implies_pt_distal}, the set $B$ contains the times of returns of a distal point to a neighborhood of itself.  By assumption, the set $A \cap B \neq \emptyset$, as was to be shown.
\end{proof}

It was shown in \cref{lemma_translates_of_dsyndetic_sets} \eqref{item_dcs_translates_of_dcs_are_dcs} that the family $\dcsyndetic$ is idempotent: $\dcsyndetic \subseteq \dcsyndetic + \dcsyndetic$.  Is this an equality, and, if not, how close is it to one?

\begin{question}
\label{quest_is_dcs_plus_dcs_equal_to_dcs}
    Is it the case that
    \[\dcsyndetic = \dcsyndetic + \dcsyndetic?\]
    Because it is known that $\dcsyndetic \subseteq \dcsyndetic + \dcsyndetic$, it is equivalent to ask: If $A \subseteq \N$ is such that the set
    \[A - \dcsyndetic \defeq \big\{ n \in \N \ \big | \ A - n \in \dcsyndetic \big\}\]
    is dynamically central syndetic, is $A$ dynamically central syndetic?    
\end{question}

We were not able to show that a set $A \subseteq \N$ is dynamically central syndetic even under the assumption that every shift $A - n$, $n \in \N$, is dynamically central syndetic.  It may be the case that a set with this property is actually thickly syndetic.

The dual form of \cref{quest_is_dcs_plus_dcs_equal_to_dcs} is equally appealing.  By \cref{lemma_filter_algebra} \eqref{item_sum_dual}, the equivalent, dual form asks: \emph{Is $\dcthick \subseteq \dcthick + \dcthick$, that is, is the family of sets of pointwise recurrence idempotent?}  A positive answer to this question would show that the family $\dcPS$ is idempotent, giving a positive answer to \cref{quest_is_dcps_idempotent} below.


\subsection{Dynamically thick sets}
\label{sec_dy_thick_sets}

It is simple to see from the definitions that if $A \subseteq \N$ is a set of pointwise recurrence, then for all dynamically central syndetic sets $B \subseteq \N$, there exists $n \in A$ such that the set $B-n$ is dynamically central syndetic.  Does the converse hold, that is, does this property characterize being a set of pointwise recurrence?

\begin{question}
\label{quest_alt_char_of_dcthick}
    Is it true that a set $A \subseteq \N$ is a set of pointwise recurrence if and only if for all $B \in \dcsyndetic$, there exists $n \in A$ such that $B-n \in \dcsyndetic$?
\end{question}

If we denote by $\dcsyndetic - \dcsyndetic$ the family of sets of the form $B - \dcsyndetic$ where $B \in \dcsyndetic$, then \cref{quest_alt_char_of_dcthick} asks whether or not $\dcthick = (\dcsyndetic - \dcsyndetic)^*$.  The equivalent, dual form of this equality is $\dcsyndetic = \dcsyndetic - \dcsyndetic.$ Thus, an equivalent form of \cref{quest_alt_char_of_dcthick} asks whether or not a set $A \subseteq \N$ is dynamically central syndetic if and only if there exists $B \in \dcsyndetic$ such that $B - \dcsyndetic \subseteq A$.  An answer to \cref{quest_nec_and_suff_for_rxu_dcs} would likely be very helpful in addressing this question.

There are several other relatives of \cref{quest_alt_char_of_dcthick} that are of interest.
\begin{enumerate}
    \item Is it true that a set $A \in \dthick$ if and only if for all $B \in \dsyndetic$, there exists $n \in A$ such that $B-n \in \dcsyndetic$?  Equivalently: is it true that $A \in \dsyndetic$ if and only if there exists $B \in \dsyndetic$ such that $B - \dcsyndetic \subseteq A$?

    \item Is it true that a set $A \in \dcthick$ if and only if for all $B \in \central$, there exists $n \in A$ such that $B-n \in \central$?  Equivalently: is it true that $A \in \dcsyndetic$ if and only if there exists $B \in \central$ such that $B - \central \subseteq A$?

    \item Is it true that a set $A \in \dthick$ if and only if for all $B \in \PS$, there exists $n \in A$ such that $B-n \in \central$?  Equivalently: is it true that $A \in \dsyndetic$ if and only if there exists $B \in \PS$ such that $B - \central \subseteq A$?\\
\end{enumerate}



It is a simple exercise to partition a thick set into a disjoint union of two thick sets.  The same task for dynamically thick sets or sets of pointwise recurrence does not appear to be quite so simple.

\begin{question}
\label{ques:partition_into_two_pointwise_recurrence}
    Can every dynamically thick set (resp. set of pointwise recurrence) be partitioned into two disjoint dynamically thick sets (resp. sets of pointwise recurrence)?
\end{question}

There is a positive answer to the analogue of \cref{ques:partition_into_two_pointwise_recurrence} for sets of topological recurrence (\cref{example:set-topological-partition}), sets of measurable recurrence (\cite{Forrest_recurrence-in-dynamical-systems}, \cref{example:set-measurable-partition}), sets of Bohr recurrence \cite[Prop. 1.4]{Le-interpolation-first}, and sets of pointwise recurrence for nilsystems \cite[Lemma 3.3]{Le_interpolation_for_nilsequences}. Though not explicitly written in these terms, the proofs of all of these results follow by demonstrating the $\sigma$-compactness of the dual families, as defined and described in \cref{sec:sigma-compact}.  Since the family of dynamically central syndetic sets is not $\sigma$-compact (\cref{theorem_dy_syndetic_is_not_sigma_compact}), a new strategy must be devised to answer \cref{ques:partition_into_two_pointwise_recurrence}.\\

We showed in \cref{cor_weak_structure_result} that all dynamically thick sets take the form $\bigcup_{B \in \family} (B \cap H_B)$, where $\family$ is a robustly syndetic collection of sets (recall \cref{def_robust_syndetic}) and $(H_B)_{B \in \family}$ is a collection of thick sets.  This result seems of limited use, however, since we were not able to say much about the family $\family$.



\begin{question}
\label{quest_robust_synd_family}
    Is it true that for all dynamically thick sets $A \subseteq \N$, there exists a countable, robustly syndetic collection $\family$ of subsets of $\N$ and, for all $B \in \family$, a thick set $H_B \subseteq \N$, such that
    \begin{align}
        \label{eqn_enhanced_dt_form}
        A = \bigcup_{B \in \family} \big( B \cap H_B \big)?
    \end{align}
\end{question}

A positive answer to \cref{quest_robust_synd_family} would strengthen \cref{cor_weak_structure_result} and would open the door to applications.  As an example, we show in the following lemma how a positive answer could be used to give a positive answer to \cref{ques:partition_into_two_pointwise_recurrence}.


\begin{lemma}
\label{lemma_solution_to_partition_problem}
    If the answer to \cref{quest_robust_synd_family} is positive, then every dynamically thick set can be partitioned into two disjoint dynamically thick sets.
\end{lemma}


\begin{proof}
    Suppose the answer to \cref{quest_robust_synd_family} is positive.  Let $A \in \dthick$, and write the set $A$ in the form given in \eqref{eqn_enhanced_dt_form}. Enumerate $\family$ as $\{B_1, B_2, \ldots\}$. For ease of notation, for each $i \in \N$, write $H_i$ instead of $H_{B_i}$. By passing to subsets of the $H_i$'s, we may assume without loss of generality that $H_i = \bigcup_{j=1}^{\infty} I_{i, j}$ where $I_{i,j}$ are finite intervals in $\N$ satisfying $\lim_{j \to \infty} |I_{i,j}| = \infty$ and, for all $i, j \in \N$, $\max I_{i,j} < \min I_{i, j+1}$.

    It is a simple exercise to choose positive integers
    \[n^{(1)}_1 < n^{(1)}_2 < n^{(2)}_1 < n^{(1)}_3 < n^{(2)}_2 < n^{(3)}_1 < n^{(1)}_4 < n^{(2)}_3 < n^{(3)}_2 < n^{(4)}_1 < \cdots\]
    so that the intervals $I_{i, n^{(i)}_j}$, $i, j \in \N$, are disjoint.  For $i \in \N$, define two thick sets
    \[
        H^{(1)}_i = \bigcup_{j=1}^{\infty} I_{i, n^{(i)}_{2j-1}} \quad \text{ and } \quad  H^{(2)}_i = \bigcup_{j=1}^{\infty} I_{i, n^{(i)}_{2j}}.
    \]
    Let
    \[  
        A_1 = \bigcup_{i=1}^{\infty}  \big(B_i \cap H^{(1)}_i \big) \quad \text{ and } \quad A_2 = \bigcup_{i=1}^{\infty}  \big(B_i \cap H^{(2)}_i \big).
    \]
    By construction, the sets $A_1$ and $A_2$ are disjoint subsets of $A$.  By appending $A \setminus (A_1 \cup A_2)$ to $A_1$, we get that $A = A_1 \cup A_2$.  Moreover, both $A_1$ and $A_2$ have the form in \cref{cor_weak_structure_result}, and so they are dynamically thick.
\end{proof}


In the event of a negative answer to \cref{quest_robust_synd_family}, one could weaken the requirement on the collection $\family$ while keeping enough to salvage the reasoning in \cref{lemma_solution_to_partition_problem}.  Such a modification to \cref{quest_robust_synd_family} could be: \emph{Is there is a collection $\family$ with the property that $\{ 1_B \ | \ B \in \family\}$ is a $\sigma$-compact subset of $\{0,1\}^{\N}$ and, for all $B \in \family$, a thick set $H_B \subseteq \N$ such that the map $\family \to \{0,1\}^{\N}$ given by $B \mapsto H_B$ is continuous such that $A$ has the form in \eqref{eqn_enhanced_dt_form}?}



\subsection{Dynamically piecewise syndetic sets}
\label{sec_open_quests_dps_sets}

It was shown in \cref{sec_dcps_relations} that family of dynamically central piecewise syndetic sets sits between the families of central and piecewise syndetic sets:
\[\central \subseteq \dcPS \subseteq \PS.\]
Central sets are known to have an abundance of combinatorial configurations \cite[Ch. 14]{hindman_strauss_book_2012}. In \cref{thm:Brauer_and_dcPS}, we show that every $\dcPS$ set contains ``Brauer''-type polynomial configurations $x, y, x + p_1(y), \ldots, x + p_k(y)$. On the other hand, \cref{rem:combinatorial_primes} shows that $\dcPS$ sets (in fact, dynamically thick sets) need not be $\IP$ sets. This naturally begs the question: to what extent do $\dcPS$ sets contain finite and infinite combinatorial configurations?  The following is a specific example of a question along these lines.  Recall the definition of an $\IP_0$ set from \cref{ex_ip_k_sets_are_compact}.

\begin{question}
\label{ques:pointwise_recurrence_IP_0}
    Must every dynamically central piecewise syndetic set be an $\IP_0$ set?
\end{question}

It follows from \cref{thm:Brauer_and_dcPS} that $\dcPS$ sets are $\IP_2$ sets -- they contain configurations of the form $\{x, y, x+y\}$ -- but we were not able to iterate in order to find higher-order finite sums.  A simpler variant of \cref{ques:pointwise_recurrence_IP_0} is: \emph{Must dynamically thick sets contain configurations of the form $\{x,y,z, x+y,x+z,y+z, x+y+z\}$?}

A positive answer to \cref{ques:pointwise_recurrence_IP_0} would combine with deep results of Furstenberg and Katznelson to give more context to the result in \cref{thm_dcps_is_poly_mult_comm_rec}, which says that a general family of times of set returns in ergodic theory is $\dcPS^*$.
A corollary of a result of Furstenberg and Katznelson \cite[Thm. 10.3]{furstenberg_katznelson_1985} (see the discussion and derivation in \cite[Sec. 7.1]{bergelson_glasscock_2020}) gives that a narrower family of times of set returns is $\IP_0^*$ and hence, by a positive answer to \cref{ques:pointwise_recurrence_IP_0}, $\dcPS^*$. Thus, Furstenberg and Katznelson's result would go part of the way toward explaining the result in \cref{thm_dcps_is_poly_mult_comm_rec}.  To recover the full result in \cref{thm_dcps_is_poly_mult_comm_rec} via this line of reasoning would require an answer to a long-standing open question: \emph{Can the results of Bergelson and McCutcheon (eg. \cite[Thm. 7.12]{Bergelson-McCutcheon-IPpolynomialSzemeredi}) be upgraded from $\IP^*$ to $\IP_0^*$?}\\

Finite sums structure is related to idempotency, which is addressed in following question.

\begin{question}
    \label{quest_is_dcps_idempotent}
    Is the family of dynamically central piecewise syndetic sets idempotent, that is,
    \[\dcPS \subseteq \dcPS + \dcPS?\]
\end{question}

Note that the families $\central$ and $\dcsyndetic$ are idempotent. Were the family $\dcthick$ idempotent (ie., if we had a positive answer to \cref{quest_is_dcs_plus_dcs_equal_to_dcs}), then by \cref{lemma_idempotent_classcap_is_idempotent,lemma_translates_of_dsyndetic_sets}, we would have a positive answer to \cref{quest_is_dcps_idempotent}.

Combining \cref{lemma_basic_results,lemma_translates_of_dps_sets}, we have that
\[\dcPS \subseteq \dPS \subseteq \dsyndetic + \dcPS \subseteq \dPS + \dcPS.\]
This is superficially close to demonstrating a positive answer to \cref{quest_is_dcps_idempotent}.  In fact, it is natural to guess that perhaps $\dcPS \subseteq \dcsyndetic + \dcPS$.  Note, however, that \cref{lemma_basic_results} relies on the shift-punch machinery in \cref{sec_char_of_dy_central_synd_sets}.  It would be interesting and potentially useful to find a proof of \cref{lemma_basic_results} that avoids this heavy machinery.


\bibliographystyle{abbrv}
\bibliography{dynthick}

%AFFILIATION
\bigskip
\bigskip
\footnotesize
\noindent
Daniel Glasscock\\
\textsc{University of Massachusetts Lowell}\par\nopagebreak
\noindent
\href{mailto:daniel_glasscock@uml.edu}
{\texttt{daniel{\_}glasscock@uml.edu}}

\bigskip
\noindent
Anh N. Le\\
\textsc{University of Denver}\par\nopagebreak
\noindent
\href{mailto:daniel_glasscock@uml.edu}
{\texttt{anh.n.le@du.edu}}



\end{document}